\input{preamble}

\begin{document}

\title{Algebraic Geometry}
\maketitle

\phantomsection
\label{section-phantom}

\tableofcontents

\section{Algebraic sets}
\label{section-algebraic-sets}

\noindent
The {\it affine space} $\mathbb{A}^n$ is the
Cartesian
product of the field $n$ times.
An {\it algebraic set} is the zero set of an
ideal
in $k[x_1,…,x_n]$.

$V(I)$ and lots of properties…

\begin{lemma}
\label{lemma-irreducible-variety-prime-ideal}
An algebraic set is irreducible if and only
if
its defining ideal is prime.
\end{lemma}

\begin{proof}
First notice that in
the case that $Y$ is an irreducible
subvariety,
this ideal is prime: if $fg \in I_Y(U)$,
then $V(f) \cup  V(g) \supset Y\cap U$
so that $(V(f) \cup  V(g))\cap Y = Y\cap U$,
and then it must be true that $V(f)\cap
Y\subseteq V(g)\cap Y$
or that $V(f)\cap Y \supseteq V(g)\cap Y \cap
U$
since $Y$ is irreducible, and thus, say,
$V(f)\supseteq Y \cap U$,
that is, $I(V(f))\subseteq Y \cap U$
so that $f \in I(Y \cap U)$.
\end{proof}

\begin{theorem}[Weak Nullstellensatz]
\label{theorem-weak-nullstellensatz}
(Misha's formulation.) Let $A \subset
\mathbb{C}^n$
be an affine variety and $\mathcal{O}_A$ the
ring of polynomial
functions on $A$.
Then every maximal ideal in $A$
is an ideal of a point $a \in A$.
\end{theorem}

\noindent
The idea behind the radical idea is that
we can take roots of elements:

\begin{definition}
\label{definition-radical-ideal}
$I \subset R$ is {\it radical}
if for every $x \in R$ we have $x^n \in I
\implies x \in I$.
\end{definition}

\begin{definition}
\label{definition-nilpotent}
An element $a \in R$ is called {\it
nilpotent}
if $a^n = 0$.
\end{definition}

\begin{exercise}
\label{exercise-no-nilpotents-iff-radical}
$A/I$ has no nilpotents if and only if $I$ is
a radical ideal.
\end{exercise}

\noindent
We may understand the Nullstellensatz
as saying that algebraic subsets of an affine
variety
are the radical ideals of its coordinate
algebra.

\begin{theorem}[Nullstellensatz]
\label{theorem-nullstellensatz}
$$
I(V(I))=\sqrt{I}
$$
where $I$ is an ideal of $k[x_1,…,x_n]$.
\end{theorem}

\begin{proof}
Prove the Weak Nullstellensatz first
and then use Rabinowitsch trick
(which uses localization).
\end{proof}

\noindent
(Prove this:)
$k[x_1,…,x_n]/I$ has no nilpotents if and
only if
$I = \sqrt{I}$.

\medskip\noindent
Now we discuss how to decompose a variety in
irreducible components.

\begin{definition}
\label{definition-irreducible-variety}
An algebraic set is {\it irreducible}
if it cannot be decomposed as a union of two
properly
contained algebraic sets.
\end{definition}

\begin{lemma}
\label{lemma-irred-iff-coord-ring-reduc}
An affine variety is irreducibly if and only
if
its coordinate ring has no zerodivisors.
\end{lemma}

\begin{proof}
A zero divisor gives an expression $fg=0$.
\end{proof}

\begin{exercise}
\label{exercise-surje-map-irred-varie}
If $\varphi:A \to B$ is a surjective morphism
of varieties
and $B$ is irreducible, so is $A$.
\end{exercise}


To prove the decomposition in
irreducible ideals we must discuss
Noetherianity.

\begin{definition}
\label{definition-noetherian-ring}
A ring $R$ is {\it noetherian}
if any ascending chain of ideal stabilizes.
\end{definition}

\begin{theorem}[Hilbert's basis theorem]
\label{theorem-hilberts-basis}
If $R$ is a Noetherian ring, so is $R[x]$.
Equivalently, every ideal in $R[x]$ is
finitely generated.
\end{theorem}

\begin{lemma}
\label{lemma-noeth-iff-ideal-finit-gener}
A ring $R$ is Noetherian
if and only if
all its ideals are finitely generated.
\end{lemma}

\begin{proof}
Try!
\end{proof}

\begin{theorem}[Hilbert's basis theorem,
Misha's version]
\label{theorem-hilberts-basis-misha}
Any finitely generated ring is Noetherian.
\end{theorem}

\noindent
The idea is to use this theorem
to show that any algebraic set may be
decomposed
as a finite union of irreducible varieties.
The idea of the proof is to reduce the
original algebraic set
to a union of two proper subvarieties, and
repeat the process until we arrive at
irreducible varieties.

\begin{definition}
\label{definition-noeth-topol-space}
A topological space is {\it Noetherian}
if any descendending sequence of closed
subsets stabilizes.
\end{definition}

\begin{definition}
\label{definition-zariski-topology}
{\it Zariski topology} is the topology
whose closed sets are algebraic sets.
\end{definition}

\begin{exercise}[Misha]
\label{exercise-varie-zaris-compa}
Any variety is Zariski compact.
\end{exercise}

\noindent
To prove that any variety has a decomposition
as a finite union of irreducible components,
we first show pointwise that every point
is contained in a finite union of irreducible
components.
With existence at hand we just prove
in a similar way the existence of the
decomposition
for all the variety.

\medskip\noindent
I don't know where to put this yet:

\begin{definition}
\label{definition-degree}
\begin{reference}
\href{https://en.wikipedia.org/wiki/%
Degree_of_an_algebraic_variety}{Wikipedia}.
\end{reference}
The {\it degree} of an affine or projective
variety $X$ of dimension $n$ is the number of
points
of intersection of $X$ and $n$ hyperplanes in
general position.
\end{definition}

\begin{exercise}
\label{exercise-hypersurfaces-of-pn}
\begin{reference}
\cite[Exercise I.2.8]{hart}
\end{reference}
A projective variety $Y \subseteq
\mathbb{P}^n$
has dimension $n-1$ if and only if
it is the zero set of a single irreducible
homogeneous polynomial
$f$ of positive degree.
\end{exercise}

\begin{proof}
$Y$ has dimension $n-1$ if and only if there
is a chain
of irreducible subsets $Y \supset Y_2\supset…
\supset Y_{n-1}$
if and only if there is a chain of
prime ideals $I(Y) \subset I_2
\subset…\subset I_{n-1}$.
But $S:=k[x_0,…,x_n]$ has Krull dimension
$n+1$,
so that $I(Y)$ has height 1
(not considering the irrelevant ideal) and
thus principal
since $S$ is a UFD.
\end{proof}

\begin{proof}
($\implies$) Suppose $Y$ ir irreducible
and $\dim Y=n-1$. Pick any nonzero
homogeneous $g \in I(Y)$.
Factor $g$ into homogeneous irreducible
facts and choose one such facto $f$.
Then $f$ is homogeneous, irreducible and
vanishes on $Y$,
so $Y \subseteq V(f)$.
Both $Y$ and $V(f)$ are irreducible
projective
varieties and have the same dimension $n-1$,
hence $Y=V(I(f))$.

($\impliedby$) Suppose $Y=V(f)$ for some
nonconstant homogeneous polynomial $f$.
Then it has dimension $n-1$.
\end{proof}

\section{Regular functions, rational
functions and morphisms}
\label{section-regular-functions}

\noindent
We start with the intuitive definition, yet
different from Harshorne's,
which works well for affine varieties.
A {\it regular} or {\it polynomial map}
is $\varphi:V \subset \mathbb{A}^n \to W
\subset \mathbb{A}^m$
if there exist $P_1,…,P_m \in k[x_1,…,x_n]$
such that
$\varphi(a_1,…,a_n)=(P_1(a_1,…,a_n),…,
P_m(a_1,…,a_n))$
and $\varphi(V) \subset W$.
If $\varphi$ is polynomial then there is an
induced
map in coordinate rings
$\varphi^*:k[W] \to k[V]$ given by pullback.
This map is well defined since
$\varphi^*(I(W))\subset I(V)$.

A polynomial map $\varphi:V \to W$ is an {\it
isomorphism}
if it has a left and right inverse that is
also a polynomial function.
$V \cong W \iff k[V] \cong k[W]$ for affine
varieties.

\medskip\noindent
Here is Hartshorne's definition:

\begin{definition}
\label{definition-regular-function}
Let $Y$ be a quasi-affine variety
(i.e. an open set of an affine variety).
A function $Y \to k$ is called {\it regular}
at $P \in Y$
if there is an open set $U \ni P$ and
$g,h \in k[x_1,…,x_n]$ such that
$f=g/h$ on $U$. $f$ is {\it regular} if it is
regular
at all points of $Y$.
\end{definition}

\noindent
And the definition for projective varieties:

\begin{definition}
\label{definition-regul-funct-proje}
Let $Y \subset \mathbb{P}^N$ be a
(quasi)projective variety.
A function $f:Y \to k$ is regular at $P \in
Y$
if there is an open set $U \ni P$ such that
$f=g/h$ for two homogeneous polynomials of
the same degree,
with $h$ nonzero in $U$.
\end{definition}

\noindent
Note that although homogeneous polynomials
are not functions on $\mathbb{P}^n$
(because we can rescale points, but the
functions will give different outputs),
the quotient of two homogeneous polynomials
is (because the rescaling factor
cancels out).

\begin{definition}
\label{definition-morphism}
A {\it morphism} of varieties $\varphi:V \to
W$
is a function that pulls back polynomial
functions
$f:W \to k$ to polynomial functions
$\varphi^*f:V \to k$.
\end{definition}

\begin{exercise}
\label{exercise-morphisms-def}
$\varphi$ is a morphism of affine varieties
if and only if
it is a polynomial function.
\end{exercise}

\begin{proof}
If $\varphi$ is polynomial then it is a
morphism:
if $f:W \to k$ is polynomial,
then
$\varphi^*f(a_1,…,a_n)=f(P_1(a_1,…,a_n),…,
P_m(a_1,…,a_n))$
is also polynomial.
Conversely, if $\varphi$ is a morphism
we see that its coordinate functions are
polynomials
by computing $y_i\varphi^*$ where $y_i:W \to
k$
is $i$-th coordinate function (which is
clearly a polynomial map).
\end{proof}

\medskip\noindent
Now we introduce the notion of local ring
(stalk).
Here's Olivier's (intuitive) definition:
Let $V$ be an affine variety (:= irreducible
algebraic set)
and $p \in V$.
$$
\mathcal{O}_p(V)
:= \{f \in k(V):\exists g,h \in k[V]
\text{ s.t. }f=g/h \text{ and } h(p) \neq
0\}\subset k(V).
$$

Compare with Hartshorne's:

\begin{definition}
\label{definition-local-ring}
Let $Y$ be a variety (affine or projective)
and $P \in Y$.
The {\it local ring} of $P$ on $Y$ is
the ring of germs of regular functions on $Y$
near $P$.
\end{definition}

\begin{exercise}
\label{exercise-local-ring-is-local}
The local ring is local with maximal ideal
given by functions vanishing at $P$.
\end{exercise}


\begin{exercise}
\label{exercise-stalk-affine-varieties}
Show that if $V$ is an affine variety,
$\mathcal{O}_p(V) =
k[V]_{\mathfrak{m}_p}$
where $\mathfrak{m}_p$ is the set of regular
functions
that vanish in $p$.
\end{exercise}

\begin{proof}
Consider the obvious map: take a quotient of
classes
$[g]/[h]$ in the localization of the
coordinate ring of $V$
and map it to the rational function $g/h$.
This is well defined because if $g-g' \in
I(V)$,
then they obviously coincide in any open set
of $V$
because they coincide in all of $V$, hence
forming the same germ.
The key fact is that $\mathcal{O}_p(V)$ is
the germs of rational
functions of $V$.
It is injective because if a germ $g/h$ is
zero, then there's
an open of $V$ where $g$ vanishes, implying
that $g=0$
since it is a polynomial (or holomorphic)
map.
It is obviously surjective.
\end{proof}

\begin{definition}
\label{definition-degree-zero-localization}
Let $S=\bigoplus_{d\geq 0} S_d$ be a
graded ring and let $f \in S$ be
homogeneous.
The localization $S_f$ is again graded by
\[
\deg(a/f^m)=\deg(a)-m\deg(f)
\]
for homogeneous $a \in S$.
We write $(S_f)_0$ for its degree-zero
part.
Thus $(S_f)_0$ consists of the fractions
$a/f^m$ such that
\[
\deg(a)=m\deg(f).
\]
More generally, if $\mathfrak{p}$ is a
homogeneous prime of $S$, we write
$S_{(\mathfrak{p})}$ for the degree-zero
part of the localization of $S$ at the
homogeneous elements outside
$\mathfrak{p}$.
\end{definition}

\begin{proposition}
\label{proposition-standard-open-proj}
Let $Y=\operatorname{Proj} S$ and let
$f \in S$ be homogeneous.
Then the open set
\[
D_+(f)=\{\mathfrak{p}\in Y:f\notin
\mathfrak{p}\}
\]
is affine, with coordinate algebra
\[
\mathcal{O}_Y(D_+(f))=(S_f)_0.
\]
In other words, restricting to the chart
where $f$ does not vanish corresponds to
localizing at $f$ and taking degree zero.
\end{proposition}

\begin{theorem}
\label{theorem-functions-projective-case}
\begin{reference}
\cite[Theorem 3.4]{har}
\end{reference}
Let $Y \subseteq \mathbb{P}^n$ be a
projective variety
with homogeneous coordinate ring $S(Y)$.
Then:
\begin{enumerate}
\item $\mathcal{O}(Y)=k$,
\item
$\mathcal{O}_P=S(Y)_{(\mathfrak{m}_P)}$,
where $\mathfrak{m}_P$ is the ideal generated
by the homogeneous
$f \in S(Y)$ not vanishing in $P$,
\item $K(Y)\cong S(Y)_{((0))}$.
\end{enumerate}
\end{theorem}

\noindent
Now let's discuss the equivalence of
categories between
affine varieties and finitely generated
reduced algebras.

\begin{proposition}
\label{proposition-coordinate-algebra}
The coordinate algebra of an affine variety
$V$,
i.e. $k[x_1,…,x_n]/I(V)$ is isomorphic to
the ring of regular functions on $V$, i.e.
$\mathcal{O}(V)$.
\end{proposition}
The following proposition is also at
\href{https://stacks.math.columbia.edu/tag/%
01I1}{01I1}

\begin{proposition}
\label{proposition-irreducible-prime}
Let $V,W \subset \mathbb{A}^n$ are
irreducible algebraic set.
There is a bijection
$$
\Hom_{k-alg}(k[W],k[V]))
\xrightarrow{bij}\underbrace{\Mor(V,W)}_{
\text{polynomial maps}}.
$$
\end{proposition}

\begin{proof}
We map
$$
\begin{aligned}
\tilde{\varphi}:
\underbrace{k[W]}_{=k[y_1,…,y_m]/I(W)}
&\longrightarrow
\underbrace{k[V]}_{=k[x_1,…,x_n]/I(V)} \\
 y_1&\longmapsto f_1\\
&\vdots \\
y_m & \longmapsto f_m
\end{aligned}
$$

\noindent
to
$$
\begin{aligned}
\varphi: V &\longrightarrow W \\
(x_1,…,x_n) &\longmapsto (f_1(x_1,…,x_n), …,
f_m(x_1,…,x_n)).
\end{aligned}
$$

\noindent
We check that this is well-defined, i.e.
that it doesn't depend on the choice of
representatives,
namely the $f_i$. But choosing other $f_i$
which differ
by an element of the ideal would not make a
difference
when evaluating at $V$ because the ideal is
precisely $I(V)$.

The inverse map we already described above.
\end{proof}

\noindent
This shows that we have an equivalence of
categories:

\begin{lemma}
\label{lemma-algeb-varie-categ-equiv}
There is an equivalence of categories between
finitely generated reduced
$k$-algebras
and irreducible algebraic sets.
\end{lemma}


\medskip\noindent
Now let's discuss rational functions.
Again let's start with a more reasonable
explanation.
Since $k[V]$ for an irreducible affine
algebraic set is a domain,
there exists a fraction field, which is the
set of {\it rational functions}.
Functions here look like $f/g$ for $f,g \in
k[V]$,
which are not defined in $V(g)$.
Since this can lead to different ways
of defining the same function
--- the correct definition shall be that
we take the greatest set where the function
can be defined.

\begin{definition}
\label{definition-rational-ring}
The {\it ring of rational
functions} is the set of regular functions
defined
on open sets with the identifications that
$(U,f)\sim(V,g)$
if they coincide in an intersection. (Like
the stalk,
but there's no point.)
\end{definition}

Compare this to Voisin's way of describing
a meromorphic function:
it's a function that can be written locally
as a quotient of two holomorphic functions.



\section{Dimension: Krull dimension,
transcendence degree}
\label{section-dimension}

\begin{definition}
\label{definition-irreducible-topological}
A nonempty subset $Y$ of a topological space
$X$
is {\it irreducible} if it cannot be
expressed
as a union of two proper subsets,
each of which is closed in $Y$.
\end{definition}

\begin{definition}
\label{definition-topol-space-dimen}
If $X$ is a topological space,
the {\it dimension} of $X$ is the supremum of
all integers $n$
such that there exists a chain
$Z_0 \subset Z_1\subset … \subset Z_n$
of distinct irreducible closed subsets of
$X$.
The {\it dimension of an algebraic variety}
is its dimension as a topological space.
\end{definition}

\begin{definition}
\label{definition-krull-dimension}
In a ring $A$, the {\it height} of a prime
ideal $\mathfrak{p}$
is the supremum of all integers $n$
such that there exists a chain
$\mathfrak{p}_0
\subset\mathfrak{p}
_1\subset…\subset\mathfrak{p}_n=\mathfrak{p}$
of distinct prime ideals.
The {\it Krull dimension} of $A$ is the
supremum of the heights
of all prime ideals.
\end{definition}

\begin{proposition}
\label{proposition-dimension}
If $Y$ is an affine algebraic set,
the dimension of $Y$ is equal to the
dimension
of its coordinate ring.
\end{proposition}

\begin{theorem}
\label{theorem-dimen-trans-degre}
Let $k$ be a field and let $B$ be an integral
domain
which is a finitely generated $k$-algebra.
Then:
\begin{enumerate}
\item the dimension of $B$ is equal to the
transcendence degree
of the quotient field $K(B)$ of $B$ over $k$;
\item For any prime ideal $\mathfrak{p}$ in
$B$, we have
$$
\text{height }\mathfrak{p}+ \dim B /
\mathfrak{p}=\dim B.
$$

\end{enumerate}
\end{theorem}

\begin{exercise}
\label{exercise-affine-twisted-cubic}
Let $Y\subseteq \mathbb{A}^3$ be the set
$Y= \{(t,t^2,t^3:t \in k\}$.
Show that $Y$ is an affine variety of
dimension 1.
Find generators for the ideal $I(Y)$.
Show that $A(Y)$ is isomorphic to a
polynomial ring in one variable over $k$.
We say that $Y$ is given by the {\it
parametric representation}
$x=t$, $y=t^2$, $z=t^3$.
\end{exercise}

\begin{proof}
It's immediate that $x^2-y$ and $x^3-z$ must
be in
the ideal defining $Y$. Let
$J=(x^2-y,x^3-z)$.
Consider the ring map
$$
\begin{aligned}
\varphi:k[x,y,z] &\longrightarrow k[t] \\
x &\longmapsto t\\
y &\longmapsto t^2\\
z &\longmapsto t^3
\end{aligned}
$$

\noindent
Since $\varphi$ is given literally by the
parametrization of $Y$,
it follows from definition that the ideal of
$Y$
is the kernel of $\varphi$.

It's clear the in the kernel of $\varphi$
lie the polynomials $x^2-y$ and $x^3-z$.
But how to make sure that $J$ is the complete
kernel?

The procedure by AI is the following.
First notice that there is an isomorphism
$k[x,y,z]/J \cong k[x]$
given by $f(x,y,z)+J \mapsto f(x,x^2,x^3)$.
This shows that $k[x,y,z]/J$ is a domain, and
thus $J$ is prime.
Further, the projection $k[x,y,z] \to
k[x,y,z]/J$,
which obviously has kernel $J$,
is exactly the map I introduced before from
$k[x,y,z]$ to $k[t]$ when we see $k[x,y,z]/J$
as $k[t]$.
Since the kernel of such map is the ideal of
$Y$ by definition,
we are done.
\end{proof}

\section{Smooth points}
\label{section-smooth-points}

\noindent
The notion of nonsingularity for affine
varieties is what we usually call
``the Jacobian criterion'',
i.e. it's given by the vanishing
of partial derivatives.
We shall see how to translate
that to commutative algebra,
which allows for a definition applicable
to non-affine varieties.

\begin{definition}
\label{definition-nonsingular-affine}
Let $Y \subseteq \mathbb{A}^n$ be ann affine
variety,
and let $f_1,…,f_t \in R=k[x_1,…,x_n]$ be a
set of generators
for the ideal of $Y$.
$Y$ is {\it nonsingular at a point $P \in Y$}
if the rank of the matrix
$\left(\frac{\partial f_i}{\partial
x_j}(P)\right)$
is $n-r$, where $r$ is the dimension of $Y$.
$Y$ is {\it nonsingular} if it is nonsingular
at every point.
\end{definition}

\noindent
It can be checked that this definition is
independent of the choice of generators.

\begin{definition}
\label{definition-residue-field}
Let $R$ be a ring and $\mathfrak{m}$ a
maximal ideal of $R$.
The {\it residue field} is  $R/\mathfrak{m}$.
\end{definition}

\begin{definition}
\label{definition-regular-local-ring}
Let $R$ be a noetherian local ring with
maximal ideal $\mathfrak{m}$
and residue field $k:=R/\mathfrak{m}$.
$R$ is a {\it regular local ring} if
$\dim_k \mathfrak{m}/\mathfrak{m}^2=\dim R$.
\end{definition}

\noindent
Upshot. $\mathfrak{m}$ is the ideal of
functions vanishing at $P$.
This means that the constant term of $h \in
\mathfrak{m}$ is zero.
Differentiate $h$.
Now the constant term is what originally was
the linear part.
If $h'$ has no constant term, that is, if $h$
vanishes at $P$
up to order 1,
then $h$ can be seen as a product $h=fg$ for
$f,g \in \mathfrak{m}$.
Indeed, the smallest nonzero term of $h$ is
of degree 2,
so we can make it look like a product of two
things of degree 1.
Conversely, products $h=fg$ vanish up to
order 1 since
the first-order coefficient is given by the
derivative
(Taylor expansion)
and $dh=df\cdot g+fdg =0$ at $P$.

Then we can interpret
$\mathfrak{m}/\mathfrak{m}^2$ as follows.
Among all the functions which vanish at $P$,
i.e. inside $\mathfrak{m}$,
we identify any two functions if their
difference
is in $\mathfrak{m}^2$,
that is, if their difference vanishes up to
order 1 at $P$,
that is, if their linear parts coincide at
$P$.

Thus, $\mathfrak{m}/\mathfrak{m}^2$ is the
space
of linear parts of polynomials vanishing at
$P$.

\begin{theorem}[Jacobian criterion]
\label{theorem-nonsi-jacob-iff-regul-local}
Let $Y \subseteq \mathbb{A}^n$ be an affine
variety.
Let $P \in Y$ be a point.
Then $Y$ is nonsingular at $P$
if and only if
the local ring $\mathcal{O}_{P,Y}$ is a
regular local ring.
\end{theorem}

\begin{proof}
Still hard to grasp completely,
but an important insight is that the map
$$
\begin{aligned}
\nabla: k[x_1,…,x_n]  &\longrightarrow k^n \\
 f &\longmapsto \nabla f=\left(\frac{\partial
 f}{\partial x_1},
…,\frac{\partial f}{\partial x_n}\right)
\end{aligned}
$$

\noindent
has kernel $\mathfrak{m}^2$ and is surjective
when restricted to $\mathfrak{m}$
(because $\nabla(x_i-p_i)=(0,…,0,1,0,…0)$, so
we construct a base of $k^n$)
so we have the isomorphism
$$
\mathfrak{m}/\mathfrak{m}^2=k^n,
$$
which we understand as the space of linear
directions at $P$.
\end{proof}

\begin{definition}
\label{definition-nonsingular}
Let $Y$ be any variety,
$Y$ is {\it nonsingular} at a point $P \in Y$
if the local ring $\mathcal{O}_{P,Y}$ is a
regular local ring,
that is,
if $\dim_k\mathfrak{m}/\mathfrak{m}^2=\dim
\mathcal{O}_{P,Y}$.
\end{definition}

\noindent
Actually the Jacobian criterion
also works for projective hypersurfaces
(by simply looking at the what the
differential of a homogeneous polynomial is):

\begin{lemma}[Jacobian test for projective hypersurfaces]
\label{lemma-projective-jacobian}
Let
$$
  X=V(F)\subset \mathbb P^n
$$
be a projective hypersurface, where
$F\in k[X_0,\ldots,X_n]$ is homogeneous.
Let
$$
  p=[a_0:\cdots:a_n]\in X.
$$
Then $p$ is singular on $X$ iff
$$
  \frac{\partial F}{\partial X_0}(p)
  =
  \cdots
  =
  \frac{\partial F}{\partial X_n}(p)
  =
  0.
$$
Equivalently, $p$ is smooth iff at least
one homogeneous partial derivative of
$F$ does not vanish at $p$.
\end{lemma}

\begin{proof}
Choose an affine chart containing $p$,
say $X_0\neq 0$. Put
$$
  x_i=\frac{X_i}{X_0}
  \qquad
  i=1,\ldots,n.
$$
On this chart the hypersurface is defined
by the dehomogenized equation
$$
  f(x_1,\ldots,x_n)
  =
  F(1,x_1,\ldots,x_n).
$$
By the affine Jacobian criterion, $p$ is
singular iff
$$
  \frac{\partial f}{\partial x_i}(p)=0
  \qquad
  i=1,\ldots,n.
$$
But these partials are exactly the
corresponding homogeneous partials
$$
  \frac{\partial F}{\partial X_i}(p),
  \qquad
  i=1,\ldots,n.
$$

It remains only to include the missing
partial with respect to $X_0$. Since $F$
is homogeneous, Euler's formula gives
$$
  \sum_{i=0}^n
  X_i\frac{\partial F}{\partial X_i}
  =
  dF,
$$
where $d=\deg F$. At points of $X$ we
have $F=0$. Hence
$$
  \sum_{i=0}^n
  X_i\frac{\partial F}{\partial X_i}
  =
  0.
$$
If $X_0\neq 0$ and all the other partials
vanish, then the $X_0$-partial vanishes
too. Thus the affine chart criterion is
equivalent to the vanishing of all
homogeneous partial derivatives.
\end{proof}


\section{Smooth points of curves}
\label{section-smooth-points-curves}

\noindent
Singularity theory for plane curves
is particularly important for exam questions.
It can also be developed slightly differently
from the general construction so we put it
in this separate section.

\medskip\noindent
Upshot: the multiplicity of a point
is how many tangent lines there are at the
point. A point is smooth if and only if
it has multiplicity 1.

\medskip\noindent
The multiplicity of a plane curve $F$ at the
origin
is the degree of the lowest nonzero
homogeneous term $F_m$
of $F=F_m+F_{m+1}+…+F_n$.
The tangent lines at the origin are the
linear factors of $F_m$, which decomposes in
linear components
by the following homogenization trick:

\begin{exercise}
\label{exercise-factor-in-linear-terms}
Let $f_m$ be a homogeneous polynomial of
degree $m$.
Show that $f_m$ is a product of $m$ linear
factors
$f_m=\prod (b_ix+c_iy)$.
\end{exercise}

\begin{proof}
Upshot: fix $y=1$, factor as product of
linear coefficients
using fundamental theorem of algebra, and
then rehomogenize.
The last is because a 1-variable polynomial
$F$ is homogeneous of degree $m$ if and only
if
$F(\lambda x)=\lambda^mF(x)$. So if you have
$F(x)=f(x,1)$
you rehomogenize by putting
$y^mF(x/y)=y^mf(x/y,1)$.
Which is strange because $y^mF(x/y)$ is in
$k[x,y,1/y]$,
but that doesn't matter because when you plug
it in you get
 $$
y^mf(x/y,1)=y^mc\prod(\frac{x}{y}-\alpha_i)=
c\prod(x-\alpha_iy),
$$
(where $c$ is a constant coming from making
the product having no
coefficients on $x$ after using fundamental
theorem of algebra)
which is already a product of linear
polynomials.
We conclude by showing that in fact
$y^mf(x/y,1)=f(x,y)$,
but that's just literally
$$
y^mf(x/y,1)=y^m\sum_{i}a_i \frac{x^i}{y^i}
=\sum_ia_i x^iy^{m-i}.
$$
\end{proof}

\noindent
Thus tangent lines can have multiplicity.
A {\it node} is an ordinary double point,
which means that there are two tangent lines
and they are not the same.
A {\it cusp} is a double point with a double
tangent line.

\begin{exercise}[Intersection multiplicity]
\label{exercise-intersection-multiplicity}
If $Y,Z \subseteq \mathbb{A}^2$ are two
distinct curves,
given by equations $f$ and $g$, and if $P \in
Y \cap Z$,
we define the {\it intersection multiplicity}
$(Y \cdot Z)_P$
of $Y$ and $Z$ at $P$ to be the length of the
$\mathcal{O}_P$-module $\mathcal{O}_P/(f,g)$.
\begin{itemize}
\item Show that $(Y\cdot Z)_P$ is finite,
and $(Y\cdot Z)_P \geq  \mu_P(Y) \mu_P(Z)$.
\end{itemize}
\end{exercise}


\medskip\noindent
The upshot about Bézout theorem is that the
``size'' of the intersection
of two curves is the product of their
degrees.
In nice cases this is just the number of
intersection points.
In even nicer cases one of the curves is a
line,
in which case the product of degrees is
just the degree of the curve,
so that we compute the degree of a curve by
{\it the number of intersection points with a
hyperplane}.

But the notion of ``size'' of the
intersection
has to be made precise, which is done via the
notion of multiplicity.
The idea is that intersecting with a line
$\ell=ax +by+c$ where, say,  $b \neq 0$,
can be done by substituting $y$
for $-\alpha x-b$.
This makes the ``resultant''
$f_\ell(x):=f(x,\alpha x+\beta)$
be a polynomial in one variable,
which can be factored as a product of linear
terms, each of the roots of which
are the intersection points,
and are now equipped with a multiplicity.

\begin{definition}
\label{definition-multiplicity-line}
The {\it multiplicity} or {\it intersection
index}
of $f\in k[x,y]\setminus k$ and
the line $\ell=ax+by+c$ (where $a \neq 0$ or
$b\neq 0$) is
$$
(\ell,f)_p =
\begin{cases}
0\qquad &\text{if }p \not\in V(f,\ell) \\
\infty \qquad &\text{if }p \in V(\ell) \text{
and }V(\ell) \subset V(f)\\
m_i \qquad & \text{if }p=(x_i,\alpha
x_i+\beta)
\end{cases}
$$
where $f_\ell(x):=f(x,\alpha x+\beta)$
factors as $c\prod(x-x_i)^{m_i}$.

If $V(\ell)\not \subset V(f)$, the integer
$$
\text{deg}(f)-\sum_{i=1}^rm_i
$$
is the {\it multiplicity of $\ell,f$ at
infinity}.
\end{definition}

\begin{exercise}
\label{exercise-asymptotic-direction}
If $f=\sum_{i=0}^df_i \in k[x,y]$ for $f_i$
homogeneous of degree $i$,
an {\it asymptotic direction} of $f$ is
$ax+by$
which divides $f_d$.

Show that  $m_\infty \neq 0$ if and only if
$y-\alpha x$
is an asymptotic direction of $f$.
\end{exercise}

\begin{proof}
This uses a theorem using the resultant…
\end{proof}

\begin{theorem}[Bézout]
\label{theorem-bezout}
Let $k$ be an algebraically closed field.
Let  $C,P \subset \mathbb{P}_k^2$ be two
curves without
irreducible components in common
and suppose that $C \cap D= \{p_1,…,p_r\}$
with intersection multiplicities  $m_1,…,m_r$
in $p_1,…,p_r$.
Then
$$
\text{deg}f\text{deg}g=\sum_{i=1}^r m_i.
$$
\end{theorem}


\begin{proposition}
\label{proposition-multiplicity-point}
Let $f \in k[x,y]\setminus k$ and $p \in
V(f)$.
There exists an integer $m_p(f)\geq 1$ called
the {\it multiplicity of $p$ in $f$}
such that for every line $\ell$ passing
through  $p$
 $$
(\ell,f)_p \geq m
$$
and $(\ell,f)_p>m$ for at most $m$ lines and
at least one line.
\end{proposition}

\begin{definition}
\label{definition-tangent-line}
Let $p=(x_0,y_0) \in V(f)$.
Write $f(x+x_0,y+y_0)=f_m(x,y)+\text{lower
degree terms} $,
i.e. translate $f$ so that the origin is in
its zero set
and write as homogeneous polynomials.
Then put
$f_m(x,y)=\prod_{i=1}^r(a_ix+b_iy)^{m_i}$.
The lines $\ell_i=a_ix+b_iy$ are called
{\it tangent lines} to $f$ and each has a
multiplicity $m_i$.
\end{definition}

\begin{definition}
\label{definition-smoot-ordin-point-curve}
A point $p \in V(f)$ is {\it smooth} if
$m_p(f)=1$
and  {\it ordinary} if all its tangent lines
have multiplicity 1.
\end{definition}

\noindent
Thus, ordinary points are singular,
but these are mild singularities.
Essentially they are points where $m$
straight lines intersect transversally.

Definition of intersection index of two plane
curves
without common components.
Definition of intersection index of two plane
curves
with common components.
Properties of intersection index (7 of them).
These 7 properties characterize uniquely the
intersection index.

The multiplicity of a root $\alpha$ of a
power series
$f\in k[[x]]$,
i.e. the greatest integer such that
$(x-\alpha)^r|f$,
equals $\dim_k k[[x]]/(x^r)=r$.
(\href{https://www.youtube.com/live/%
EspNBiSU8F4?si=UEcdxrysAS47ZA9H&t=5230}{
Here}'s why.)
This leads to the alternative definition
of intersection index as
$\dim_k
k[[x,y]]/(f(x-x_0,y-y_0),g(x-x_0,y-y_0))$.
This satisfies the 7 properties of the
intersection index.


\noindent
\begin{proposition}[Jacobian criterion]
\label{proposition-jacobian-criterion}
The singular locus of $V(f)$ is
$V(f,\partial_x,\partial_y)$.
\end{proposition}



\begin{theorem}[Plücker formulas]
\label{theorem-plucker-formulas} Let $C
\subset \mathbb{P}^2$ be an irreducible curve
of degree $d$ whose only singularities are
$\delta$ nodes and $\chi$ ordinary cusps
(cusps such that the tangent line intersects
with multiplicity 3 at the singular point).
Let \begin{itemize} \item $i$ be the number
of inflection lines (lines $\ell$ such that
$(\ell,C)_p \geq 3$, and, in this case, do
not intersect another point with multiplicity
$\geq 2$), and

\item $d^*$ the number or lines tangent to
$C$ passing through a point $p \not\in C$
such that $p$ is not in any tangent lines at
the singular points, tangent inflection lines
or bitangent lines (lines that are tangent to
two different points). \end{itemize}

\noindent Then $$
\begin{aligned}
d(d-1)&=d^*+2\delta+3\chi\\
3d(d-2)&=i+6\delta+8\chi. \end{aligned}
$$
\end{theorem}


\section{Zariski tangent space}
\label{section-zariski-tangent-space}

\noindent
Upshot: the Zariski tangent space is the set
of vectors killed simultaneously by the
differentials of all the generators of the
ideal.

\medskip\noindent
Let $X$ be an affine variety over $k$,
and let $p\in X$ be a $k$-point.
Let $\mathcal O_{X,p}$ be the local
ring at $p$, and let $\mathfrak m_p$
be its maximal ideal.

\begin{definition}
\label{definition-zariski-cotangent}
The Zariski cotangent space of $X$
at $p$ is
$$
  T_p^*X:=\mathfrak m_p/\mathfrak m_p^2.
$$
\end{definition}

\begin{definition}
\label{definition-zariski-tangent}
The Zariski tangent space of $X$
at $p$ is
$$
  T_pX
  :=
  \operatorname{Hom}_k
  (\mathfrak m_p/\mathfrak m_p^2,k).
$$
\end{definition}

The slogan is:
$T_pX$ is the space of first-order
directions at $p$ which preserve the
equations of $X$.

\begin{proposition}
\label{proposition-jacobian-tangent}
Let
$$
  X=V(f_1,\ldots,f_r)
  \subset \mathbb A_k^n,
$$
and let $p=(p_1,\ldots,p_n)\in X$.
Then, under the identification
$T_p\mathbb A^n\simeq k^n$,
$$
  T_pX=\ker J(p),
$$
where
$$
  J(p)
  =
  \left(
    \frac{\partial f_j}{\partial x_i}(p)
  \right)_{j,i}.
$$
Equivalently, $v=(v_1,\ldots,v_n)$
belongs to $T_pX$ iff
$$
  \sum_{i=1}^n
  \frac{\partial f_j}{\partial x_i}(p)
  v_i
  =
  0
$$
for every $j=1,\ldots,r$.
\end{proposition}

\begin{proof}
A vector $v=(v_1,\ldots,v_n)\in k^n$
defines a derivation at $p$
$$
  D_v:k[x_1,\ldots,x_n]\to k
$$
by
$$
  D_v(g)
  =
  \sum_{i=1}^n
  \frac{\partial g}{\partial x_i}(p)
  v_i.
$$
This is the algebraic version of
evaluating $dg_p$ on $v$.

For $v$ to be tangent to $X$, the
equations defining $X$ must remain true
to first order. Hence $D_v$ must vanish
on the ideal $(f_1,\ldots,f_r)$.

It is enough to impose this on the
generators. Thus
$$
  D_v(f_j)=0
$$
for every $j$. But
$$
  D_v(f_j)
  =
  \sum_{i=1}^n
  \frac{\partial f_j}{\partial x_i}(p)
  v_i.
$$
Therefore $v\in T_pX$ exactly when
it lies in the kernel of the Jacobian
matrix evaluated at $p$.
\end{proof}

\begin{remark}
\label{remark-residue-field}
The definition
$$
  T_pX
  =
  \operatorname{Hom}_k
  (\mathfrak m_p/\mathfrak m_p^2,k)
$$
assumes that $p$ is a $k$-point, i.e.
that its residue field is
$$
  \kappa(p)=k.
$$
In general, the correct definition is
$$
  T_pX
  =
  \operatorname{Hom}_{\kappa(p)}
  (\mathfrak m_p/\mathfrak m_p^2,\kappa(p)).
$$

Thus the only change is:
replace $k$ by the residue field
$\kappa(p)$.

Over an algebraically closed field $k$,
closed points are exactly $k$-points.
So, for classical varieties 
(zero-sets of ideals, affine or projective)
over $\mathbb C$, one may safely write
$$
  T_pX
  =
  \operatorname{Hom}_{\mathbb C}
  (\mathfrak m_p/\mathfrak m_p^2,\mathbb C).
$$

Scheme-theoretically there are also
non-closed points, such as generic
points. At such points, the residue
field need not be $k$.
\end{remark}



\section{Smoothness exercises}
\label{section-smoothness-exercises}

\noindent
This is the result of a study session.
For now I'll leave it separate from the notes
in the previous section --- task: merge
the two sections.

\medskip\noindent
Let
$$
X=V(f_1,\dots,f_r)\subset \mathbb{A}^n.
$$
The Zariski tangent space at $p\in X$ is
$$
T_pX=\ker J(p),
$$
where
$$
J(p)=
\left(\frac{\partial f_i}{\partial x_j}(p)\right).
$$
Thus
$$
\dim T_pX=n-\operatorname{rank}J(p).
$$
The point $p$ is smooth if and only if
$$
\dim T_pX=\dim_pX.
$$
Equivalently,
$$
\operatorname{rank}J(p)=n-\dim_pX,
$$
the local codimension of $X$ in the ambient affine
space.

\medskip\noindent
For a hypersurface
$$
X=\{F=0\}\subset \mathbb{A}^n,
$$
a point $p\in X$ is smooth if and only if
$$
dF_p\neq 0.
$$
Equivalently, at least one partial derivative of $F$
is nonzero at $p$. A point is singular if
$$
F(p)=0,
\qquad
F_{x_1}(p)=\cdots=F_{x_n}(p)=0.
$$

\medskip\noindent
For a plane curve singularity at the origin, write
$$
f=f_m+f_{m+1}+\cdots+f_d,
$$
where each $f_i$ is homogeneous of degree $i$ and
$f_m\neq 0$. Then
$$
\operatorname{mult}_0(f)=m,
$$
and the tangent cone is
$$
TC_0V(f)=V(f_m).
$$
The multiplicity is the degree of the tangent cone,
counted with multiplicity. The number of distinct
tangent directions may be smaller.

\begin{exercise}
\label{exercise-smoothness-node}
Classify the singularity of
$$
C=\{y^2=x^2+x^3\}\subset \mathbb{A}^2
$$
at the origin using multiplicity and tangent cone.
\end{exercise}

\begin{proof}
Write
$$
f=y^2-x^2-x^3.
$$
The lowest homogeneous part is
$$
f_2=y^2-x^2=(y-x)(y+x).
$$
Thus the multiplicity is $2$ and the tangent cone is
$$
TC_0C=\{y=x\}\cup \{y=-x\}.
$$
There are two distinct tangent directions, so the
origin is a node.
\end{proof}

\begin{exercise}
\label{exercise-zariski-tangent-node}
For
$$
C=\{y^2=x^2+x^3\},
$$
compute the Zariski tangent space at the origin and
compare it with the tangent cone.
\end{exercise}

\begin{proof}
The local ring is
$$
\mathcal{O}_{C,0}
=
\left(k[x,y]/(y^2-x^2-x^3)\right)_{(x,y)}.
$$
Let $\mathfrak m=(x,y)$. Since the defining equation
has no linear part, it imposes no linear relation in
$\mathfrak m/\mathfrak m^2$. Therefore
$$
\mathfrak m/\mathfrak m^2
=\langle \bar x,\bar y\rangle_k,
$$
so
$$
\dim_k\mathfrak m/\mathfrak m^2=2.
$$
Equivalently, the Jacobian vanishes at the origin, so
$$
T_0C=\mathbb{A}^2.
$$
But the tangent cone is
$$
TC_0C=\{y^2-x^2=0\},
$$
a union of two lines inside $\mathbb{A}^2$.
The tangent space sees only linear equations, while the
tangent cone sees the first nonzero homogeneous part.
\end{proof}

\begin{exercise}
\label{exercise-cusp-local-ring-not-regular}
Let
$$
C=\{y^2=x^3\}\subset \mathbb{A}^2.
$$
Show that $\mathcal{O}_{C,0}$ is not regular.
Compare this with the normalization
$$
\mathbb{A}^1\to C,
\qquad
 t\mapsto (t^2,t^3).
$$
\end{exercise}

\begin{proof}
We have
$$
\mathcal{O}_{C,0}
=
\left(k[x,y]/(y^2-x^3)\right)_{(x,y)}.
$$
The maximal ideal is
$$
\mathfrak m=(x,y).
$$
The equation $y^2-x^3$ has no linear part, so it gives
no relation in $\mathfrak m/\mathfrak m^2$. Therefore
$$
\mathfrak m/\mathfrak m^2
=\langle \bar x,\bar y\rangle_k,
\qquad
\dim_k\mathfrak m/\mathfrak m^2=2.
$$
But $C$ is a curve, so
$$
\dim \mathcal{O}_{C,0}=1.
$$
Hence
$$
\dim_k\mathfrak m/\mathfrak m^2
\neq
\dim \mathcal{O}_{C,0},
$$
and the local ring is not regular.

The parametrization induces
$$
k[x,y]/(y^2-x^3)\hookrightarrow k[t],
\qquad
x\mapsto t^2,
\quad
 y\mapsto t^3.
$$
Thus the coordinate ring is $k[t^2,t^3]\subset k[t]$.
The element $t$ is integral over $k[t^2,t^3]$, since
$$
t^2-x=0,
$$
but $t\notin k[t^2,t^3]$. Thus the cusp is not normal,
and its normalization is $\mathbb{A}^1$.
\end{proof}

\begin{exercise}
\label{exercise-node-xy-normalization}
Let
$$
C=\{xy=0\}\subset \mathbb{A}^2.
$$
Compute $\dim_k\mathfrak m/\mathfrak m^2$ at the
origin and describe the normalization geometrically.
\end{exercise}

\begin{proof}
The local ring at the origin is
$$
\left(k[x,y]/(xy)\right)_{(x,y)}.
$$
The maximal ideal is $\mathfrak m=(x,y)$. Since $xy$
has no linear part, it gives no relation in
$\mathfrak m/\mathfrak m^2$. Hence
$$
\mathfrak m/\mathfrak m^2
=\langle \bar x,\bar y\rangle_k,
\qquad
\dim_k\mathfrak m/\mathfrak m^2=2.
$$
The curve has dimension $1$, so the origin is not
regular.

Geometrically, $C$ is the union of the two coordinate
axes. Its normalization separates the two branches:
$$
\widetilde C=\mathbb{A}^1\sqcup \mathbb{A}^1.
$$
On coordinate rings, the normalization map is
$$
k[x,y]/(xy)\longrightarrow k[x]\times k[y],
$$
where
$$
x\longmapsto (x,0),
\qquad
 y\longmapsto (0,y).
$$
The origin becomes two points, one on each component.
\end{proof}

\begin{exercise}
\label{exercise-projective-jacobian-criterion-plane-curve}
Let
$$
Y=V(F)\subset \mathbb{P}^2
$$
be a plane curve, with $F$ homogeneous of degree $d$.
Show that
$$
Y\text{ is singular at }[p]
\iff
F_X(p)=F_Y(p)=F_Z(p)=0.
$$
\end{exercise}

\begin{proof}
The condition is projectively well-defined because
$F_X,F_Y,F_Z$ are homogeneous of degree $d-1$.
Thus their vanishing does not depend on the choice of
homogeneous representative.

Work in the chart $Z\neq 0$ and write
$$
f(x,y)=F(x,y,1).
$$
Then $Y$ is singular at $[a:b:1]$ if and only if
$$
f_x(a,b)=f_y(a,b)=0.
$$
But
$$
f_x=F_X(x,y,1),
\qquad
f_y=F_Y(x,y,1).
$$
Thus singularity in the affine chart gives
$$
F_X(a,b,1)=F_Y(a,b,1)=0.
$$
Since $F$ is homogeneous, Euler's formula gives
$$
XF_X+YF_Y+ZF_Z=dF.
$$
At a point of $Y$ we have $F=0$. Since also
$F_X=F_Y=0$ and $Z=1$, we get
$$
F_Z=0.
$$
The converse is immediate from the affine Jacobian
criterion in the chart.
\end{proof}

\begin{exercise}
\label{exercise-cone-over-smooth-plane-curve-singularities}
Let
$$
Y=\{F(X,Y,Z)=0\}\subset \mathbb{P}^2
$$
be a smooth plane curve of degree $d>1$. Let
$$
C(Y)=\{F(x,y,z)=0\}\subset \mathbb{A}^3
$$
be the affine cone. Show that
$$
\operatorname{Sing}C(Y)=\{0\}.
$$
\end{exercise}

\begin{proof}
The singular locus of the cone is
$$
\{F=F_x=F_y=F_z=0\}\subset \mathbb{A}^3.
$$
If $q=(a,b,c)\neq 0$ is singular, then $[q]$ is a
point of $Y$. In an affine chart, the vanishing of all
partials of $F$ implies the vanishing of the chart
partials. Thus $[q]$ would be a singular point of
$Y$, contradiction.

At the origin, $F(0)=0$ because $F$ is homogeneous of
positive degree. Since $d>1$, each partial derivative
has degree $d-1\geq 1$, so all partials vanish at the
origin. Thus the origin is singular.
\end{proof}



\section{Blow-up}
\label{section-blow-up}

\noindent
The blow-up of the origin in $\mathbb{A}^2$ is
$$
\operatorname{Bl}_0\mathbb{A}^2
=
\{((x,y),[u:v])\in \mathbb{A}^2\times
\mathbb{P}^1:xv=yu\}.
$$
The projection
$$
\pi:\operatorname{Bl}_0\mathbb{A}^2\to
\mathbb{A}^2
$$
is the first projection. The exceptional divisor is
$$
E=\pi^{-1}(0)\cong \mathbb{P}^1.
$$

\medskip\noindent
In the chart $u\neq 0$, write $[u:v]=[1:v]$.
Then
$$
y=xv.
$$
In the chart $v\neq 0$, write $[u:v]=[u:1]$.
Then
$$
x=uy.
$$
Thus the blow-up replaces the origin by the set of
possible tangent directions.

\medskip\noindent
Let $C=\{f(x,y)=0\}\subset \mathbb{A}^2$ be a
plane curve passing through the origin. The total
transform is
$$
\pi^{-1}(C).
$$
The strict transform is
$$
\widetilde C=
\overline{\pi^{-1}(C\setminus\{0\})}.
$$
In the chart $y=xv$, if
$$
f(x,xv)=x^m g(x,v),
$$
where $x$ does not divide $g$, then
$$
\widetilde C=\{g=0\}.
$$
Indeed, away from the exceptional divisor $x=0$, the
factor $x^m$ is invertible, so the equation is just
$g=0$. Taking the Zariski closure keeps the equation
$g=0$ and allows $x=0$.

\medskip\noindent
For a plane curve the divisor formula is
$$
\pi^*C=\widetilde C+mE,
$$
where $m$ is the multiplicity of $C$ at the origin.
Equivalently, $m$ is the multiplicity with which
$E$ appears in the total transform.

\begin{exercise}
\label{exercise-blow-up-parabola}
Blow up the origin and compute the strict transform
of
$$
C=\{y=x^2\}\subset \mathbb{A}^2.
$$
\end{exercise}

\begin{proof}
Use the chart $y=xv$. Then
$$
y-x^2=xv-x^2=x(v-x).
$$
Thus
$$
\pi^*C=E+\widetilde C,
$$
where
$$
E=\{x=0\},
\qquad
\widetilde C=\{v=x\}.
$$
The strict transform meets $E$ at
$$
x=0,
\qquad
v=0.
$$
So it meets $E$ in one point. The strict transform is
smooth, since it is the line $v=x$ in the chart.
The original curve was already smooth; the blow-up
only records its tangent direction.
\end{proof}

\begin{exercise}
\label{exercise-cuspidal-curve-blow-up}
Let $Y$ be the cuspidal curve $y^2=x^3$ in
$\mathbb{A}^2$.
Blow up the point $O=(0,0)$, let $E$ be the
exceptional curve,
and let $\widetilde Y$ be the strict transform
of $Y$.
Show that $E$ meets $\widetilde Y$ in one point,
and that $\widetilde Y\cong \mathbb{A}^1$.
\end{exercise}

\begin{proof}
Use the chart $y=xv$. Then
$$
y^2-x^3=x^2v^2-x^3=x^2(v^2-x).
$$
Therefore
$$
\pi^*Y=2E+\widetilde Y,
\qquad
\widetilde Y=\{v^2=x\}.
$$
The exceptional divisor is $E=\{x=0\}$, so
$$
E\cap \widetilde Y=\{x=0,v^2=0\}.
$$
This is one point, with multiplicity $2$.
The strict transform is smooth because
$$
F=v^2-x
$$
has
$$
dF=2v\,dv-dx,
$$
which is nowhere zero on $F=0$.
Moreover, the equation $x=v^2$ parametrizes
$\widetilde Y$ by $v$, so
$$
\widetilde Y\cong \mathbb{A}^1.
$$
\end{proof}

\begin{exercise}
\label{exercise-blow-up-tacnode}
Blow up the origin in
$$
C=\{y^2=x^4\}\subset \mathbb{A}^2.
$$
Show that the strict transform has a node.
\end{exercise}

\begin{proof}
Again use $y=xv$. Then
$$
y^2-x^4=x^2(v^2-x^2).
$$
Thus
$$
\pi^*C=2E+\widetilde C,
\qquad
\widetilde C=\{v^2=x^2\}.
$$
The strict transform meets $E=\{x=0\}$ at
$$
v^2=0,
$$
so at one point. But
$$
v^2-x^2=(v-x)(v+x).
$$
Thus the strict transform has two distinct tangent
lines at this point. Hence it has a node.
\end{proof}

\begin{exercise}
\label{exercise-blow-up-node}
Blow up the origin in
$$
C=\{y^2=x^2+x^3\}\subset \mathbb{A}^2.
$$
Show that the strict transform meets the exceptional
curve in two distinct smooth points.
\end{exercise}

\begin{proof}
In the chart $y=xv$,
$$
y^2-x^2-x^3=x^2(v^2-1-x).
$$
Thus
$$
\pi^*C=2E+\widetilde C,
\qquad
\widetilde C=\{v^2-1-x=0\}.
$$
Intersecting with $E=\{x=0\}$ gives
$$
v^2-1=0.
$$
Thus the strict transform meets $E$ at the two points
$$
v=1,
\qquad
v=-1.
$$
The strict transform is smooth because
$$
F=v^2-1-x
$$
satisfies
$$
dF=2v\,dv-dx,
$$
which is nonzero at both points.
\end{proof}

\begin{exercise}
\label{exercise-blow-up-y3-x5}
Blow up the origin in
$$
C=\{y^3=x^5\}\subset \mathbb{A}^2.
$$
Show that the exceptional divisor appears with
multiplicity $3$, and that the strict transform has a
cusp-type double point.
\end{exercise}

\begin{proof}
Use $y=xv$. Then
$$
y^3-x^5=x^3(v^3-x^2).
$$
Thus
$$
\pi^*C=3E+\widetilde C,
\qquad
\widetilde C=\{v^3=x^2\}.
$$
The strict transform meets $E$ at $(x,v)=(0,0)$.
At this point the lowest homogeneous part of
$$
F=v^3-x^2
$$
is
$$
F_2=-x^2.
$$
So the strict transform has multiplicity $2$ and one
repeated tangent line. This is cusp-type, not a node.
\end{proof}

\medskip\noindent
The blow-up of the origin in $\mathbb{A}^3$ is
$$
\operatorname{Bl}_0\mathbb{A}^3
=
\{((x,y,z),[u:v:w])\in \mathbb{A}^3\times
\mathbb{P}^2:\operatorname{rank}
\begin{pmatrix}
x&y&z\\
u&v&w
\end{pmatrix}
\leq 1\}.
$$
Equivalently, the $2\times 2$ minors vanish:
$$
xv-yu=0,
\qquad
xw-zu=0,
\qquad
yw-zv=0.
$$
The exceptional divisor is
$$
E\cong \mathbb{P}^2.
$$

\begin{exercise}
\label{exercise-blow-up-ordinary-double-surface}
Blow up the origin in the surface
$$
X=\{z^2=xy\}\subset \mathbb{A}^3.
$$
Compute the strict transform in the chart
$y=xu$, $z=xv$.
\end{exercise}

\begin{proof}
Substitution gives
$$
z^2-xy=x^2v^2-x^2u=x^2(v^2-u).
$$
Therefore
$$
\pi^*X=2E+\widetilde X,
\qquad
\widetilde X=\{v^2=u\}.
$$
This chart of the strict transform is smooth, since
$$
F=v^2-u
$$
has
$$
dF=2v\,dv-du,
$$
which is nowhere zero.
\end{proof}

\begin{exercise}
\label{exercise-blow-up-quadratic-cone}
Blow up the origin in
$$
X=\{x^2+y^2+z^2=0\}\subset \mathbb{A}^3.
$$
Compute the strict transform in the chart
$y=xv$, $z=xw$.
\end{exercise}

\begin{proof}
We get
$$
x^2+y^2+z^2=x^2(1+v^2+w^2).
$$
Thus
$$
\pi^*X=2E+\widetilde X,
\qquad
\widetilde X=\{1+v^2+w^2=0\}.
$$
The intersection with $E=\{x=0\}$ is
$$
\{x=0,1+v^2+w^2=0\}.
$$
This is part of the projectivized tangent cone inside
$E\cong \mathbb{P}^2$.
The strict transform is smooth in this chart: the only
possible critical point of $1+v^2+w^2$ would be
$v=w=0$, but this point is not on the strict transform.
\end{proof}

\begin{exercise}
\label{exercise-cone-over-smooth-curve-blow-up}
Let
$$
Y=\{F(X,Y,Z)=0\}\subset \mathbb{P}^2
$$
be a smooth plane curve of degree $d>1$.
Let
$$
C(Y)=\{F(x,y,z)=0\}\subset \mathbb{A}^3
$$
be the affine cone over $Y$.
Show that $C(Y)$ is singular only at the vertex, and
that the blow-up of the vertex resolves this
singularity. Show that the exceptional divisor of the
strict transform is naturally isomorphic to $Y$.
\end{exercise}

\begin{proof}
The singular locus of the cone is cut out by
$$
F=F_x=F_y=F_z=0.
$$
If $q=(a,b,c)\neq 0$ is a singular point of
$C(Y)$, then $[q]=[a:b:c]$ is a point of $Y$.
On an affine chart, say $Z\neq 0$, the equation is
$f(x,y)=F(x,y,1)$, and
$$
f_x=F_X(x,y,1),
\qquad
f_y=F_Y(x,y,1).
$$
Thus the vanishing of the cone partial derivatives
would make $[q]$ a singular point of $Y$, impossible.
Therefore the only possible singular point is the
origin.

Since $F$ is homogeneous of degree $d>1$, all partial
 derivatives are homogeneous of degree $d-1\geq 1$.
Hence they vanish at the origin. Thus the vertex is
singular.

Now use the description
$$
\operatorname{Bl}_0\mathbb{A}^3
=
\{(p,[\ell])\in \mathbb{A}^3\times \mathbb{P}^2:
 p\in \ell\}.
$$
If $p=t\ell$, then by homogeneity
$$
F(p)=t^dF(\ell).
$$
The total transform has the exceptional factor $t^d$.
Removing it gives the strict transform
$$
\widetilde{C(Y)}
=
\{(p,[\ell]):p\in \ell, F(\ell)=0\}.
$$
Over the exceptional divisor $p=0$ this becomes
$$
\{0\}\times Y.
$$
So the exceptional divisor of the strict transform is
naturally isomorphic to $Y$.

Moreover, the strict transform is the total space of
$$
\mathcal{O}_{\mathbb{P}^2}(-1)|_Y.
$$
Indeed, over $[\ell]\in Y$, the fiber is the line
$\ell$. Since $Y$ is smooth and a vector bundle over a
smooth variety is smooth, the strict transform is
smooth.
\end{proof}

\medskip\noindent
Now consider the blow-up of a point in the projective
plane. Let
$$
\pi:X=\operatorname{Bl}_p\mathbb{P}^2\to
\mathbb{P}^2.
$$
For example, take $p=[1:0:0]$. Then
$$
X=
\{([X:Y:Z],[U:V])\in \mathbb{P}^2\times
\mathbb{P}^1:YV=ZU\}.
$$
The exceptional divisor is
$$
E=\pi^{-1}(p)\cong \mathbb{P}^1.
$$
It parametrizes tangent directions at $p$.

\medskip\noindent
Let $H$ be the pullback class of a line in
$\mathbb{P}^2$:
$$
H=\pi^*[\text{line}].
$$
If $L$ is a line not passing through $p$, then its
strict transform has class $H$ and is disjoint from
$E$.
Therefore
$$
H^2=1,
\qquad
H\cdot E=0.
$$
If $L$ is a line passing through $p$, then locally it
has multiplicity $1$ at $p$, so
$$
\pi^*L=\widetilde L+E.
$$
Since $\pi^*L$ has class $H$, we get
$$
\widetilde L=H-E.
$$
Take two distinct lines $L,L'$ through $p$. Downstairs
they meet only at $p$. Upstairs their strict transforms
meet the exceptional curve at different points,
because they have different tangent directions.
Thus
$$
\widetilde L\cdot \widetilde L'=0.
$$
Since both have class $H-E$,
$$
0=(H-E)^2=H^2-2H\cdot E+E^2.
$$
Using $H^2=1$ and $H\cdot E=0$, we obtain
$$
E^2=-1.
$$
Thus
$$
H^2=1,
\qquad
H\cdot E=0,
\qquad
E^2=-1.
$$

\medskip\noindent
A final sanity check is the blow-up of a smooth
divisor. Blow up $\mathbb{A}^2$ along
$$
D=\{x=0\}.
$$
The ideal is principal, $I=(x)$. Blowing up a
Cartier divisor is trivial. Indeed, the Rees algebra is
$$
\bigoplus_{n\geq 0}I^nt^n
=
k[x,y]\oplus (x)t\oplus (x^2)t^2\oplus\cdots
=
k[x,y][xt].
$$
There is only one normal direction, so the
projectivized normal direction is $\mathbb{P}^0$ over
each point. Hence no new exceptional geometry appears.

\section{Rational maps}
\label{section-rational-maps}

\noindent
The category of varieties and dominant
rational maps is equivalent to the category
of finitely generated field extensions of $k$
with the arrows reversed.

\section{Normalization}
\label{section-normalization}

\noindent
``All bounded meromorphic functions are
holomorphic''. For exam curves:
find a parametrization of the rational curve.

For the abstract definitions see
Commutative Algebra Section
\ref{commutative-algebra-section-normalization}.
Here we explain how to
parametrize rational curves
and the compute normalizations of some
varieties.

A curve $C$ is rational if its
normalization is $\mathbb P^1$.
A parametrization of $C$ is then the
normalization map
$$
  \nu:\mathbb P^1\to C.
$$

In practice, for plane curves, one often
finds this map using a pencil of lines.

Suppose
$$
  C=V(f)\subset \mathbb A^2
$$
has a singular point at $p=(0,0)$.
Consider the lines through $p$:
$$
  y=tx.
$$
Substitute this into the equation:
$$
  f(x,tx)=0.
$$

Since $p$ is singular, the factor $x=0$
appears with multiplicity at least $2$.
(See Definition
\ref{definition-smoot-ordin-point-curve},
which would match the regular local ring
definition of smoothnes via
the Jacobian criterion,
Proposition
\ref{proposition-jacobian-criterion}.)
Remove this factor. The remaining
equation gives the other intersection
points of the line with $C$.

If this remaining equation can be solved
rationally for $x$ in terms of $t$, then
one gets
$$
  x=x(t),
  \qquad
  y=tx(t).
$$
This gives a parametrization of $C$.

For a singular plane cubic, this method
usually works immediately: a line through
the singular point meets the cubic with
multiplicity $2$ at the singular point,
so there is only one remaining point.

Example: for
$$
  C: y^2=x^2+x^3,
$$
put $y=tx$. Then
$$
  t^2x^2=x^2+x^3.
$$
Thus
$$
  x^2(t^2-1-x)=0.
$$
Ignoring the factor $x^2$, which is the
singular point, gives
$$
  x=t^2-1.
$$
Hence
$$
  y=t(t^2-1).
$$
So a parametrization is
$$
  t\mapsto (t^2-1,t(t^2-1)).
$$

Warning: guessing monomials
$x=t^a$, $y=t^b$ is not a general
method. It works for some cusps, but not
for nodes.

\medskip\noindent
Now I include some
Hartshorne/half-theoretical
exercises I did before.

Recall from commutative algebra
that a (affine or projective) 
variety is called {\it normal at $P$}
is $\mathcal{O}_P$ is integrally closed.
It is called {\it normal} if every point
is normal.

\begin{exercise}
\label{exercise-normal-varieties}
If $Y$ is an affine variety,
$\mathcal{O}_P$ is integrally closed
for all $P \in Y$ if and only if
$\mathcal{O}_Y$ is integrally closed.
\end{exercise}

\begin{proof}
For the converse
let $f/g \in \text{Frac}\mathcal{O}_P$
be integral.
Then there exists $P \in \mathcal{O}_P[t]$
monic killing $f/g$.
We want to show that $f/g \in \mathcal{O}_P$
for which it's enough to put $f/g$ in
$\mathcal{O}_Y$ using the closure-integrality
of $Y$.
Suppose
$$
P=\sum \frac{f_i}{g_i}t^i.
$$
Since this polynomial is not in
$\mathcal{O}_Y[t]$, we cannot use it directly
to prove integrality of $f/g$;
in fact $f/g$ might not be integral at all.
So we get rid of the fractions
considering the least common multiple
of the $g_i$.
But if we just multiply by the lcm
we get a polynomial that is not monic.
We must consider
$$
P\cdot lcm^{\text{deg}P}
\left(\frac{f}{g}\right)
=\sum lcm^{\text{deg}P}\frac{f_i}{g_i}
\left(\frac{f}{g}\right)^i
=\frac{f_i}{g_i}
\left(lcm \frac{f}{g}\right)^{\text{deg}P}
+\frac{f_i}{g_i}lcm
\left(lcm\frac{f}{g}\right)^{\text{deg}P-1}
+…
+\frac{f_0}{g_0}lcm^{\text{deg}P},
$$
which is monic over $\mathcal{O}_Y$
and kills $lcm \frac{f}{g}$,
which implies that $\frac{f}{g}
\in \mathcal{O}_Y$
and thus $g$ is never zero
(or simply a unit of the coordinate algebra),
and thus $f/g$ is in $\mathcal{O}_P$.

The direct implication is actually easier.
Let $f/g \in \text{Frac}\mathcal{O}_Y$
be integral, i.e. there is a polynomial
$P \in \mathcal{O}_Y[t]$ killing $f/g$.
Pick any point $p \in Y$.
Then $P$ can be taken to be a polynomial
in $\mathcal{O}_p[t]$, which in fact
kills the fraction
$\overline{f/g} \in \text{Frac}\mathcal{O}_p$
in the local ring of $p$.
Since $\mathcal{O}_p$ is integrally closed,
$f/g \in \mathcal{O}_P$.
Thus, thinking of $f/g$ as a meromorphic
function, we see that $g(p) \neq  0$ for all $p$,
meaning the fraction $f/g$ is in fact
regular. Alternatively, thinking of $f/g$
as an element of the localization by the
maximal ideal $\mathfrak{m}_P$,
saying that $f/g \in \mathcal{O}_P$
means that $g \not\in \mathfrak{m}_P$
for all $P$, which is to say that $g$
is in no maximal ideal of $k[Y]$,
which means that $g$ is a unit.
\end{proof}

\begin{exercise}
\label{exercise-normal-varieties-examples}
\begin{enumerate}
\item Show that every conic in $\mathbb{P}^2$
is normal.
\item Show that the quadric surfaces
$Q_1,Q_2$ in $\mathbb{P}^3$ given by
$Q_1:xy-zw$ and $Q_2:xy=z^2$ are normal.
\item Show that the cuspidal cubic
$y^2=x^3$ in $\mathbb{A}^2$
is not normal.
\end{enumerate}
\end{exercise}

\begin{proof}
\begin{enumerate}
\item We need to show that every conic
in $\mathbb{P}^2$ is rational.
\item
\item The element $x^2/y$ has a minimal
polynomial in the local ring at the origin,
namely $P(t)=t-x$, but it clearly
does not lie in the local ring since $y$
vanishes at $(0,0)$.
\end{enumerate}
\end{proof}

\begin{exercise}
\label{exercise-universal-prop-normalization}
Let $Y$ be an affine variety.
Show that there is a normal affine
variety $\tilde{Y}$, and a morphism
$\pi:\tilde{Y}\to Y$,
with the property that whenever $Z$
is a (affine) normal variety and 
$\varphi:Z \to Y$ is a dominant morphism,
then there is a unique morphism
$\theta:Z \to \tilde{Y}$ such that
$\varphi=\pi \circ \theta$.
\end{exercise}

\begin{proof}
Let $\tilde{Y}=\Spec \tilde{A}$
where $\tilde{A}$ is the integral closure
of $A=\mathcal{O}_Y$.
Then by Exercise
\ref{exercise-normal-varieties}
it is normal.
Now if the dominant map $\varphi:Z \to Y$
induces an injective algebra morphism
$\varphi^\#:A \to B$
and $B$ is integrally closed,
we can put $\theta^\#:\tilde{A} \to B$
by the inclusion.
Since any element $x \in \tilde{A}$
is the root of a unique monic polynomial
$P \in A[t] \subset B[t]$,
and since $\theta^\#(x)$ must be
the root of a monic polynomial in $B[t]$
because $B$ is integrally closed,
we conclude that $\theta^\#(x)$
is also a root of $P$.
… but what if the roots are different?

A correct solution is as follows.
The inclusion $A \xymatrix{\ar@{^{(}->}[r]&}
B$ induces an inclusion on fraction fields
$\text{Frac}A \xymatrix{\ar@{^{(}->}[r]&}
\text{Frac} B$.
The key observation is that by definition
an element in the integral closure 
$\tilde{A}$ of $A$ is an element in
the fraction field $\text{Frac} A$.
This says that any integral element
can already be seen as an element in
$\text{Frac}B$.
But since $B$ is integrally closed and the
minimal polynomial of any such element
has coefficients in $A \subset B$,
we are done: an integral element of $A$
is mapped to an element of $B$ making
the following diagram commute
$$
\xymatrix{
A\ar[dr]_{\pi^\sharp}
\ar@{^{(}->}[r]^{\varphi^\sharp}
&  B\\
&  \tilde{A}\ar@{-->}[u]^{\theta^\sharp}
}
$$
For uniqueness, suppose that
$\theta':\tilde{A} \to B$ is another
map making the above diagram commute.
Then for any $x=a/b$ integral over $A$ 
we have
$\theta'(b)\theta'(x)=\theta'(a)$.
But this coincides on $A$
with $\theta^\sharp$
we introduced before, giving
$\theta^\sharp(b)(\theta'(x)-\theta^\sharp)
=0$. Since there are no zerodivisors on $B$,
we obtain $\theta\equiv\theta^\sharp$.
\end{proof}

\begin{exercise}
\label{exercise-projectively-normal-vars}
A projective variety $Y \subseteq
\mathbb{P}^n$ is {\it projectively normal} 
(with respect to the given embedding)
if its homogeneous coordinate ring
$S(Y)$ is integrally closed.
\begin{enumerate}
\item If $Y$ is projectively normal
then $Y$ is normal.
\item There are normal varieties 
(i.e. stalks are integrally closed)
in projective space which are not 
projectively normal.
For example, let $Y$ be the twisted
quartic curve in $\mathbb{P}^3$
given parametrically by
$(x,y,z,w)=(t^4,t^3u,tu^3,u^4)$.
\item Show that the twisted quartic curve
$Y$ above is isomorphic to $\mathbb{P}^1$,
which is projectively normal. Thus projective
normality depends on the embedding.
\end{enumerate}
\end{exercise}

\begin{proof}
\begin{enumerate}
\item Take an affine chart $D_+(f)$.
By Proposition
\ref{proposition-standard-open-proj}
the coordinate algebra of $Y$
at this open set is given by
$(S(Y)_f)_0$, the degree-0 elements
in the localization $S(Y)_f$.
By Commutative Algebra Proposition
\ref{prop-localization-integrally-closed},
localization of an integrally closed
ring remains integrally closed,
thus $S(Y)_f$ is integrally closed.
Let $A=(S(Y)_f)_0$.
If $x \in \text{Frac}(A)$ is integral
over $A$, then it is integral over
$S(Y)_f$.
Thus $x \in S(Y)_f$.
But $x$ is a quotient of degree-0
elements, hence has degree 0.
Therefore $x \in (S(Y)_f)_0=A$.

\item By the next item, the twisted quartic
is isomorphic to $\mathbb{P}^1$, hence is
normal.
It remains to see that it is not
projectively normal in this embedding.
Its homogeneous coordinate ring is
\[
R=k[t^4,t^3u,tu^3,u^4]\subset k[t,u].
\]
The element $t^2u^2$ lies in
$\text{Frac}(R)$, since
\[
t^2u^2=\frac{(t^3u)(u^4)}{tu^3}.
\]
Moreover it is integral over $R$, since
\[
(t^2u^2)^2=(t^3u)(tu^3)\in R.
\]
But $t^2u^2\notin R$: the exponent
$(2,2)$ is not in the semigroup generated
by
\[
(4,0),(3,1),(1,3),(0,4).
\]
Thus $R$ is not integrally closed.

\item Let
\[
\varphi:\mathbb{P}^1\to Y
\]
be the map
\[
[t:u]\longmapsto [t^4:t^3u:tu^3:u^4].
\]
We construct its inverse.
On the open subset of $Y$ where
$(x,y)\neq (0,0)$, define
\[
[x:y:z:w]\longmapsto [x:y].
\]
On the open subset where
$(z,w)\neq (0,0)$, define
\[
[x:y:z:w]\longmapsto [z:w].
\]
These opens cover $Y$.
On the image of $\varphi$ we have
\[
[x:y]=[t^4:t^3u]=[t:u]
\]
and
\[
[z:w]=[tu^3:u^4]=[t:u].
\]
Thus the two local formulae agree on
overlaps and define a morphism
$Y\to \mathbb{P}^1$ inverse to
$\varphi$.
Hence $Y\cong \mathbb{P}^1$.
Finally, the standard embedding of
$\mathbb{P}^1$ has homogeneous
coordinate ring $k[t,u]$,
which is integrally closed.
Thus $\mathbb{P}^1$ is projectively
normal.
\end{enumerate}
\end{proof}



\section{Divisors, invertible sheaves}
\label{section-divisors-line-bundles}

\begin{definition}
\label{definition-weil-divisor}
\begin{reference}
\href{https://stacks.math.columbia.edu/tag/%
0BE2}{0BE2},
\cite[p. 130]{hart}
\end{reference}
A {\it Weil divisor} on the variety $X$ is a
formal sum of irreducible
subvarieties over the integers.

More formally, a {\it prime divisor} is an
integral closed
subscheme $Z \subset X$ of codimension 1.
A {\it Weil divisor} is a formal sum $D=\sum
n_Z Z$ where
the sum is over prime divisors of $X$
and the collection
$\{Z|n_Z \neq  0\}$ is locally finite.
If all coefficients are non-negative we say
the divisor is {\it effective}.
\end{definition}

\noindent
A reasonable notion of order of vanishing
should then be introduced. In the case of
Riemann surfaces
this can be done using the order of
vanishings
of the denominator and numerator functions in
$f=g/h$ where $f \in K(X)^*$.
A more general construction is a bit more
involved.

\begin{definition}
\label{definition-principal-weil-divisor}
\begin{reference}
\href{https://stacks.math.columbia.edu/tag/%
0be0}{0be0}
\end{reference}
Let $X$ be a Noetherian integral scheme.
Let $f \in R(X)^*$.
The {\it principal Weil divisor associated to
$f$}
is the Weil divisor
$$
\text{div}(f)=\text{div}_X(f)
=\sum \text{ord}_Z(f)[Z]
$$
where the sum is over prime divisors and
$\text{ord}_Z(f)$ is as in Definition ?.
This makes sense by Lemma ? [a lemma about
finiteness of the sum].
\end{definition}

\begin{definition}
\label{definition-weil-divisor-class-group}
\begin{reference}
\href{https://stacks.math.columbia.edu/tag/%
0BE4}{0BE4}
\end{reference}
Let $X$ be a locally Noetherian scheme.
The {\it Weil divisor class group} of $X$,
denoted
$\text{Cl}(X)$, is the quotient
of the group of Weil divisors by the subgroup
of principal
Weil divisors.
\end{definition}

\noindent
The next section,
\href{https://stacks.math.columbia.edu/tag/%
02SE}{02SE},
shows there is an injective homomorphism
$\text{Pic}(X) \to \text{Cl}(X)$
which is also surjective if $X$ is normal
and the local rings of $X$ are UFDs
(which happens for smooth varieties).

\medskip\noindent
Earlier, in
\href{https://stacks.math.columbia.edu/tag/%
01WQ}{01WQ},
we have

\begin{definition}
\label{definition-cartier-divisor}
Let $S$ be a scheme.
\begin{itemize}
\item A {\it localy principal closed
subscheme} of $S$
is a closed subscheme whose sheaf of ideals
is locally generated by a single element.

\item An {\it effective Cartier divisor} on
$S$
is a closed subscheme $D \subset S$ whose
ideal sheaf
$\mathcal{I}_D \subset \mathcal{O}_S$ is an
invertible
$\mathcal{O}_S$-module.
\end{itemize}
\end{definition}

\noindent
Followed by Lemma
\href{https://stacks.math.columbia.edu/tag/%
01WS}{01WS},
which essentially says that a Cartier
divisor is cut locally by a single function.
After this we would show the well-known
construction of sheaf associated to an
effective Cartier divisor
and eventually construct all group operations
and subtleties
to arrive at Lemma
\href{https://stacks.math.columbia.edu/tag/%
01X0}{01X0},
which says that effective Cartier divisors on
a scheme
are the same as invertible sheaves with fixed
regular global section.

\medskip\noindent
Here's two basic facts I'm not sure how to
prove:
\begin{enumerate}
\item
$$
\mathcal{O}_{\mathbb{P}^n}(1)\cong
\mathcal{O}_{\mathbb{P}^n}(H)
$$
where $\mathcal{O}_{\mathbb{P}^n}(1)$ is
defined
as dual of the tautological bundle,
and $H$ is any hyperplane (zeroes of a
homogeneous degree-1 polynomial
in $n+1$ variables) seen as an effective
Cartier divisor,
so that $\mathcal{O}_{\mathbb{P}^n}(H)$ is
its associated line bundle.

\item
$$
H^0(\mathbb{P}^n,\mathcal{O}_{\mathbb{P}^n}
(1))
\cong
k[x_0,…,x_n]_1,
$$
that is, the space of sections of
$\mathcal{O}_{\mathbb{P}^n}(1)$
is isomorphic to the space of homogeneous
polynomials
of degree 1 in $n+1$ variables
(each of which is a hyperplane in
$\mathbb{P}^n$).
I think this might be
\href{https://www.youtube.com/live/%
bXhgeRZMu8M?si=5BPdnECUREsCUwpV&t=4601}{
here}.
\end{enumerate}

\begin{exercise}
\label{exercise-degree-d-surface}
Show that a hypersurface in $\mathbb{P}^3$
given by the zeroes of a homogeneous
polynomial
of degree $d$ has degree $d$,
i.e., if $S \subset \mathbb{P}^3$, $S=V(f)$
with
$f$ homogeneous of degree $d$,
then
 $$
\mathcal{O}_S(1).\mathcal{O}_S(1)=d.
$$
\end{exercise}

\begin{proof}
The exercise is kind of by definition no?
The intersection form counts the size of the
intersection when
the two classes are just two hypersurfaces
intersecting transversally (as is the case of
two
hyperplane sections in general position,
both of which have $\mathcal{O}_S(1)$ as
associated line bundle). The intersection
has $d$ points because its a line in
$\mathbb{P}^3$
intersecting a degree-$d$ polynomial.
And that's it.
\end{proof}

\begin{exercise}
\label{exercise-intersection-lines}
Let $S \subset \mathbb{P}^3$ be a surface
of degree $d$ and $L \subset S$ a line.
Compute $L^2$ in $S$. Check that when $d=3$,
$L^2=-1$!
\end{exercise}

\begin{proof}
We apply the genus formula:
$$
(L.L+K_S)=2p_a-2
$$
By $L$ being a line we know that $p_a=0$
and that the intersection number of $L$
with any divisor of degree $d'$ will be $d'$.
Indeed, $L.d'H=d'(L.H)=d'$
because $L.H$ is the degree of the line
(by definition of degree, namely
the intersection number of $L$ with as many
hyperplanes as its dimension).
Thus
$$
\begin{aligned}
-2=(L.L+K_S)&=L^2+(L.K_S)=L^2+(L.\mathcal{O}
_S(d-4))
=L^2+d-4\\
\iff L^2=-d+2.
\end{aligned}
$$

\noindent
where the canonical bundle of $S$
is computed using adjunction formula
$K_S=(K_{\mathbb{P}^3}+N)|_S
=(\mathcal{O}_{\mathbb{P}^3}(-4)+\mathcal{O}
_{\mathbb{P}^3}(d))$.
\end{proof}

\begin{exercise}[Bézout]
\label{exercise-bezout}
Show that if $C_d$ and $C_{d'}$ are two plane
curves
(i.e. contained in $\mathbb{P}^2$)
of degrees $d$ and $d'$ then
$C_d.C_{d'}=dd'$.
\end{exercise}

\noindent
One fun application of Bézout's theorem
is the following method for computing
degrees of line bundles in curves:

\begin{lemma}
\label{lemma-line-bundle-degree}
Let $C\subset \mathbb P^n$ be a
projective curve. Then
$$
  \deg \mathcal O_C(m)=m\deg C.
$$
In particular, if $C\subset \mathbb P^2$
is a plane curve of degree $d$, then
$$
  \deg \mathcal O_C(m)=md.
$$
\end{lemma}

\begin{proof}
By definition, $\mathcal O_C(1)$ is the
restriction of the hyperplane bundle
$\mathcal O_{\mathbb P^n}(1)$ to $C$.

A section of $\mathcal O_C(1)$ is obtained
by restricting a hyperplane equation to
$C$. For a general hyperplane $H$, the
zero divisor of this section is
$$
  H\cap C.
$$
By definition, this divisor has degree
$\deg C$. Hence
$$
  \deg \mathcal O_C(1)=\deg C.
$$

Now
$$
  \mathcal O_C(m)
  =
  \mathcal O_C(1)^{\otimes m}.
$$
Since degree is additive under tensor
products of line bundles on curves, we get
$$
  \deg \mathcal O_C(m)
  =
  m\deg \mathcal O_C(1)
  =
  m\deg C.
$$
\end{proof}

\medskip\noindent
The Euler sequence on $\mathbb P^n$ is
$$
  0
  \to
  \mathcal O_{\mathbb P^n}
  \to
  \mathcal O_{\mathbb P^n}(1)^{\oplus n+1}
  \to
  T_{\mathbb P^n}
  \to
  0.
$$

The dual Euler sequence is
$$
  0
  \to
  \Omega^1_{\mathbb P^n}
  \to
  \mathcal O_{\mathbb P^n}(-1)^{\oplus n+1}
  \to
  \mathcal O_{\mathbb P^n}
  \to
  0.
$$

\begin{example}
\label{example-canonical-projective-space}
We compute the canonical bundle of
$\mathbb P^n$.

From the dual Euler sequence
$$
  0
  \to
  \Omega^1_{\mathbb P^n}
  \to
  \mathcal O_{\mathbb P^n}(-1)^{\oplus n+1}
  \to
  \mathcal O_{\mathbb P^n}
  \to
  0,
$$
we take determinants. Since
$$
  \omega_{\mathbb P^n}
  =
  \det \Omega^1_{\mathbb P^n},
$$
we get
$$
  \omega_{\mathbb P^n}
  =
  \det
  \left(
    \mathcal O_{\mathbb P^n}(-1)^{\oplus n+1}
  \right)
  \otimes
  \det(\mathcal O_{\mathbb P^n})^{-1}.
$$
Hence
$$
  \omega_{\mathbb P^n}
  =
  \mathcal O_{\mathbb P^n}(-n-1).
$$

In particular,
$$
  \omega_{\mathbb P^2}
  =
  \mathcal O_{\mathbb P^2}(-3).
$$

Therefore, if
$C\subset \mathbb P^2$ is a plane curve
of degree $d$, adjunction gives
$$
  \omega_C
  =
  \left(
    \omega_{\mathbb P^2}
    \otimes
    \mathcal O_{\mathbb P^2}(d)
  \right)\big|_C
  =
  \mathcal O_C(d-3).
$$
\end{example}


\begin{proof}
Just notice that the line bundles associated
to
$C_d$ and $C_{d'}$ are $dH$ and $d'H$,
then use linearity of the bilinear of form.
To check why the divisor associated to $C_d$
is indeed $dH$ note that the intersection
number
of $C_d$ and a line $H$ is $d$ because $C_d$
is given by the roots of a polynomial of
degree $d$,
thus $C_d.H=d$, and since the Picard group of
$\mathbb{P}^2=\mathbb{Z}H$
we must have $C_d \sim dH$.
\end{proof}
\section{Adjunction formulas}
\label{section-adjunction formulas}

There are several statements called
adjunction formula
in different texts. All of them concern
``subvarieties'',
that is, closed embedded subschemes.

\begin{exercise}[Genus formula for a curve on
a surface]
\label{exercise-genus-formu-curve-surfa}
Let $C \to X$ be a closed embedded subscheme
of dimension $1$ (as a topological space,
i.e. pure dimension)
inside a smooth surface $X$.
Then
$2p_a-2=(\mathcal{O}_X(C),\mathcal{O}_X(C))$.
\end{exercise}

\begin{proof}
Consider the ideal sheaf exact sequence
$$
\xymatrix{
0\ar[r]&\mathcal{O}_X(-C)\ar[r]&\mathcal{O}
_X\ar[r]&\mathcal{O}_C\ar[r]&0
}
$$
This sequence splits since there is an
obvious inverse morphism
to the inclusion $\mathcal{O}_X(-C)\to
\mathcal{O}_X$, namely
mapping a function $f$ to
Then $\mathcal{O}_X\cong \mathcal{O}_X(-C)
\oplus \mathcal{O}_C$.
 $\chi(\mathcal{O}_X)=\chi$
\end{proof}


\section{Linear systems}
\label{section-linear-systems}

\noindent
A basic fact is that
$$
H^0(\mathcal{O}(d)) \cong
\{\text{homogeneous polynomials of degree
$d$}\}.
$$
And more generally,
$$
|D|=\mathbb{P}H^0(\mathcal{O}(D)).
$$

\begin{exercise}
\label{exercise-hodge-numbers-blow-up-9pts}
Let $C_1,C_2$ be two cubic curves in
$\mathbb{P}^2$.
Compute the Hodge numbers of the blowup of
the 9 points
of intersection.
\end{exercise}

\begin{proof}
The solution is:
$$
\begin{aligned}\xymatrix@C=.5em@R=.5em{
&  &   1\\
&  0&  &  0\\
0&&  10&  &  0\\
&  0&  &  0\\
&  &  1
}\end{aligned}
$$
Which follows from:
\begin{itemize}
\item $h^{0,0}=1$ because $X$ is connected.
\item $h^{1,0}=0$ because $X$ is simply
connected.
\item $h^{2,0}=0$ because the canonical
bundle of $X$ has no sections.
\item $h^{1,1}=10$ because in the exponential
exact sequence
we get $\text{Pic} \cong H^2(X,\mathbb{Z})$
(by vanishing irregularity and genus)
and we know
$\text{Pic}=\mathbb{Z}H \oplus
\mathbb{Z}E_1\oplus…\oplus \mathbb{Z}E_9$.
\end{itemize}
\end{proof}

\begin{exercise}
\label{exercise-line-bundle-degree}
Let $C\subset \mathbb P^n$ be a
projective curve. Explain what
$\mathcal O_C(1)$ means and compute its
degree.
\end{exercise}

\begin{proof}
By definition,
$$
  \mathcal O_C(1)
  =
  \mathcal O_{\mathbb P^n}(1)|_C.
$$
Thus $\mathcal O_C(1)$ is the hyperplane
bundle restricted to $C$.

A section of $\mathcal O_C(1)$ is obtained
by restricting a hyperplane equation to
$C$. Its zero divisor is
$$
  C\cap H,
$$
where $H\subset \mathbb P^n$ is a
hyperplane.

Therefore
$$
  \deg \mathcal O_C(1)=\deg C.
$$
More generally,
$$
  \deg \mathcal O_C(m)=m\deg C.
$$
This is Bezout in line-bundle language.
\end{proof}

\begin{exercise}
\label{exercise-plane-section-dimension}
Let $C\subset \mathbb P^2$ be a plane
curve not contained in any line. Compute
$h^0(C,\mathcal O_C(1))$ and describe
$|\mathcal O_C(1)|$.
\end{exercise}

\begin{proof}
The global sections of
$\mathcal O_{\mathbb P^2}(1)$ are the
linear forms
$$
  aX+bY+cZ.
$$
Thus
$$
  H^0(\mathbb P^2,\mathcal O(1))
  =
  \langle X,Y,Z\rangle.
$$

Restricting to $C$, we get sections
$$
  X|_C,\quad Y|_C,\quad Z|_C
$$
of $\mathcal O_C(1)$.

They are linearly independent. Indeed,
if
$$
  aX+bY+cZ
$$
vanished on all of $C$, then $C$ would be
contained in the line
$$
  aX+bY+cZ=0,
$$
contradicting the assumption.

Hence
$$
  h^0(C,\mathcal O_C(1))=3.
$$
Therefore
$$
  |\mathcal O_C(1)|
  =
  \mathbb P(H^0(C,\mathcal O_C(1)))
  \simeq
  \mathbb P^2.
$$
\end{proof}

\begin{exercise}
\label{exercise-sections-map}
Let $L$ be a line bundle on a variety $X$,
and let $s_0,\ldots,s_N$ be global
sections with no common zero. Describe
the morphism induced by these sections.
\end{exercise}

\begin{proof}
The sections define a morphism
$$
  \varphi:X\to \mathbb P^N
$$
by
$$
  p\mapsto [s_0(p):\cdots:s_N(p)].
$$
The condition that the sections have no
common zero guarantees that this point of
projective space is defined for every
$p\in X$.

In particular, for
$C\subset \mathbb P^2$ and
$L=\mathcal O_C(1)$, the sections
$$
  X|_C,\quad Y|_C,\quad Z|_C
$$
give the original embedding
$$
  C\hookrightarrow \mathbb P^2.
$$
\end{proof}

\begin{exercise}
\label{exercise-conic-pullback}
Let
$$
  C=V(Y^2-XZ)\subset \mathbb P^2.
$$
Use the parametrization
$$
  \nu:\mathbb P^1\to C,
  \qquad
  [U:V]\mapsto [U^2:UV:V^2].
$$
Compute $\nu^*\mathcal O_C(1)$ and the
pulled-back coordinate sections.
\end{exercise}

\begin{proof}
The coordinate sections of
$\mathcal O_C(1)$ are
$$
  X|_C,\quad Y|_C,\quad Z|_C.
$$
Pulling back by $\nu$ gives
$$
  \nu^*X=U^2,
  \qquad
  \nu^*Y=UV,
  \qquad
  \nu^*Z=V^2.
$$
These are homogeneous polynomials of
degree $2$ on $\mathbb P^1$. Therefore
$$
  \nu^*\mathcal O_C(1)
  \cong
  \mathcal O_{\mathbb P^1}(2).
$$

Moreover,
$$
  H^0(C,\mathcal O_C(1))
$$
pulls back to the subspace
$$
  \langle U^2,UV,V^2\rangle
  \subset
  H^0(\mathbb P^1,\mathcal O(2)).
$$
This is all of
$H^0(\mathbb P^1,\mathcal O(2))$.
Hence
$$
  h^0(C,\mathcal O_C(1))=3.
$$
\end{proof}

\begin{exercise}
\label{exercise-twisted-cubic}
Let $C\subset \mathbb P^3$ be the
twisted cubic, parametrized by
$$
  \nu:\mathbb P^1\to C,
  \qquad
  [U:V]\mapsto
  [U^3:U^2V:UV^2:V^3].
$$
Compute $\nu^*\mathcal O_C(1)$,
$h^0(C,\mathcal O_C(1))$, and
$|\mathcal O_C(1)|$.
\end{exercise}

\begin{proof}
The coordinate sections pull back to
$$
  U^3,\quad U^2V,\quad UV^2,\quad V^3.
$$
These are homogeneous polynomials of
degree $3$ on $\mathbb P^1$. Hence
$$
  \nu^*\mathcal O_C(1)
  \cong
  \mathcal O_{\mathbb P^1}(3).
$$

The four sections
$$
  U^3,\quad U^2V,\quad UV^2,\quad V^3
$$
form a basis of
$$
  H^0(\mathbb P^1,\mathcal O(3)).
$$
Therefore
$$
  h^0(C,\mathcal O_C(1))=4.
$$
Thus
$$
  |\mathcal O_C(1)|
  \simeq
  \mathbb P^3.
$$
The induced morphism is the original
twisted cubic embedding
$$
  C\hookrightarrow \mathbb P^3.
$$
\end{proof}

\begin{exercise}
\label{exercise-nodal-cubic}
Let $C$ be the nodal cubic
$$
  C: Y^2Z=X^2Z+X^3.
$$
Its normalization is
$$
  \nu:\mathbb P^1\to C,
$$
given by
$$
  [U:V]\mapsto
  [V(U^2-V^2):U(U^2-V^2):V^3].
$$
Compute $\nu^*\mathcal O_C(1)$ and
explain why
$$
  h^0(C,\mathcal O_C(1))=3.
$$
\end{exercise}

\begin{proof}
The coordinate sections pull back to
$$
  V(U^2-V^2),
  \qquad
  U(U^2-V^2),
  \qquad
  V^3.
$$
These are homogeneous polynomials of
degree $3$ on $\mathbb P^1$. Hence
$$
  \nu^*\mathcal O_C(1)
  \cong
  \mathcal O_{\mathbb P^1}(3).
$$

However,
$$
  h^0(C,\mathcal O_C(1))
$$
is not equal to
$$
  h^0(\mathbb P^1,\mathcal O(3))=4.
$$
Only the sections on the normalization
which descend to $C$ are allowed.

The node has two preimages:
$$
  [1:1],
  \qquad
  [-1:1].
$$
A section descending to $C$ must take
compatible values at these two points.

The pulled-back coordinate sections span
the $3$-dimensional subspace
$$
  \langle
  V(U^2-V^2),
  U(U^2-V^2),
  V^3
  \rangle
  \subset
  H^0(\mathbb P^1,\mathcal O(3)).
$$
Therefore
$$
  h^0(C,\mathcal O_C(1))=3.
$$
\end{proof}

\begin{exercise}
\label{exercise-cuspidal-cubic}
Let $C$ be the cuspidal cubic
$$
  C: Y^2Z=X^3.
$$
Its normalization is
$$
  \nu:\mathbb P^1\to C,
  \qquad
  [U:V]\mapsto [U^2V:U^3:V^3].
$$
Compute $\nu^*\mathcal O_C(1)$ and
$h^0(C,\mathcal O_C(1))$.
\end{exercise}

\begin{proof}
The coordinate sections pull back to
$$
  U^2V,
  \qquad
  U^3,
  \qquad
  V^3.
$$
These are homogeneous polynomials of
degree $3$. Hence
$$
  \nu^*\mathcal O_C(1)
  \cong
  \mathcal O_{\mathbb P^1}(3).
$$

The sections descending from $C$ are
spanned by
$$
  U^2V,\quad U^3,\quad V^3.
$$
This is a $3$-dimensional subspace of
$$
  H^0(\mathbb P^1,\mathcal O(3)).
$$
Therefore
$$
  h^0(C,\mathcal O_C(1))=3.
$$
\end{proof}

\begin{exercise}
\label{exercise-plane-quartic-canonical}
Let $C\subset \mathbb P^2$ be a smooth
plane quartic. Compute $K_C$,
$\deg K_C$, $h^0(K_C)$, and describe the
canonical map.
\end{exercise}

\begin{proof}
By adjunction for plane curves,
$$
  K_C=\mathcal O_C(d-3).
$$
Here $d=4$, so
$$
  K_C=\mathcal O_C(1).
$$
Since $C$ has degree $4$,
$$
  \deg K_C
  =
  \deg \mathcal O_C(1)
  =
  4.
$$
Also a smooth plane quartic has genus
$$
  g=\frac{(4-1)(4-2)}2=3.
$$
Thus
$$
  h^0(K_C)=g=3.
$$
The canonical linear system is
$$
  |K_C|=|\mathcal O_C(1)|.
$$
Therefore the canonical map is the given
embedding
$$
  C\hookrightarrow \mathbb P^2.
$$
\end{proof}

\section{Ample line bundles}
\label{section-ample-line-bundles}


\begin{exercise}
\label{exercise-ample-bundle-on-K3}
Let $L$ be an ample bundle on a K3 surface
$M$. Prove that
$\mathcal{L}^{\otimes 2}$ is globally
generated (that is, for each $x\in M$
there exsits a section $h \in
H^{0}(L^{\otimes 2})$ which does not vanish
in
$x$).
\end{exercise}

\begin{proof}
This just asks that the $n$ in Definition
\ref{definition-ample} is $2$ for all
$x\in X$. Because, again, $x\in X_s$ means
that $s(x)\neq 0$ because if it was,
then we could somehow write $s$ as a product
of a vanishing function on
$\mathfrak{m}_x$ and a local frame of
$\Gamma(X,\mathcal{L})$. But I guess for
the exercise do this: a line bundle is {\it
ample} if there is $n$ such that the
canonical embedding (cf Lemma
\ref{lemma-map-into-proj}) is an embedding,
i.e.
that $\mathcal{L}^{\otimes n}$ is {\it very
ample}. (Interestingly, the notion
very ampleness is defined in morphisms.tex.)
\end{proof}


\section{Riemann-Roch theorem for curves}
\label{section-riemann-roch}

\noindent
Let $D=\sum n_p p$ be a divisor on a curve.
How many rational functions $z \in K^*$
satisfy
$\text{div}(z) \geq -D$? That is, rational
functions
which can only have poles of order $\leq n_p$
at the points $p$ with $n_p>0$,
and can only have zeros of order $\geq
|n_p|$
at the points $p$ with $n_p<0$.

For example, fix a point $p$ in the curve.
Then $\mathcal{O}_C(p)$ is the space
of rational functions with a simple pole at
$p$.
$\mathcal{O}_C(2p)$ is the space of rational
functions
with a pole of order at most 2 at $p$, and so
on.
Note that $\mathcal{O}_C(p) \subset
\mathcal{O}_X (2p)
\subset \mathcal{O}(3p)…$.
That is, the more poles we allow, this space
of functions will be larger.

In classical (Riemann's) notation, this space
was denoted
$L(D)$, and its dimension $\ell(D)$.
Modern notation reads $H^0(\mathcal{O}_C(D))$
instead of $L(D)$ and
$h^0(\mathcal{O}_C(D))$
instead of $\ell(D)$.

Riemann-Roch formula for curves is
\begin{equation}
\label{equation-riemann-roch-curves}
h^0(D)-h^0(K-D)=\text{deg}(D)-g+1,
\end{equation}

\noindent
where $K$ is a divisor associated to the
curve canonically.
This is done by choosing a meromorphic
section $s$
of the bundle $\Omega_C^1$
and taking $K:=\text{div}(s)$.
Then one shows that this is independent of
the choice of $s$.

Note that by Serre duality
$h^1(D)=h^0(K-D)$.

\medskip\noindent
Here's the hands-on criterion for
computing the genus of a plane curve:

\begin{proposition}
\label{proposition-plane-genus}
Let $C\subset \mathbb P^2$ be a plane
curve of degree $d$. Then its arithmetic
genus is
$$
  p_a(C)=\frac{(d-1)(d-2)}{2}.
$$
\end{proposition}

\begin{proof}
By definition,
$$
  p_a(C)=1-\chi(\mathcal O_C).
$$
Since $C$ is cut out by one homogeneous
polynomial of degree $d$, we have the
exact sequence
$$
  0
  \to
  \mathcal O_{\mathbb P^2}(-d)
  \to
  \mathcal O_{\mathbb P^2}
  \to
  \mathcal O_C
  \to
  0.
$$
Taking Euler characteristics gives
$$
  \chi(\mathcal O_C)
  =
  \chi(\mathcal O_{\mathbb P^2})
  -
  \chi(\mathcal O_{\mathbb P^2}(-d)).
$$
We use the standard formula
$$
  \chi(\mathcal O_{\mathbb P^2}(m))
  =
  \binom{m+2}{2}.
$$
Thus
$$
  \chi(\mathcal O_{\mathbb P^2})=1
$$
and
$$
  \chi(\mathcal O_{\mathbb P^2}(-d))
  =
  \binom{2-d}{2}
  =
  \frac{(d-1)(d-2)}{2}.
$$
Therefore
$$
  \chi(\mathcal O_C)
  =
  1-\frac{(d-1)(d-2)}{2}.
$$
Hence
$$
  p_a(C)
  =
  1-\chi(\mathcal O_C)
  =
  \frac{(d-1)(d-2)}{2}.
$$
\end{proof}

\begin{exercise}
\label{exercise-riemann-roch-canonical}
Apply Riemann-Roch to $D=K_C$.
\end{exercise}

\begin{proof}
For $D=K_C$, Riemann-Roch gives
$$
  h^0(K_C)-h^0(0)
  =
  \deg K_C-g+1.
$$
We know
$$
  h^0(K_C)=g
  \qquad
  \text{and}
  \qquad
  h^0(0)=1.
$$
Hence
$$
  g-1=\deg K_C-g+1.
$$
Therefore
$$
  \deg K_C=2g-2.
$$
\end{proof}

\begin{exercise}
\label{exercise-high-degree-divisor}
Let $D$ be a divisor on $C$ with
$$
  \deg D>2g-2.
$$
Compute $h^0(D)$.
\end{exercise}

\begin{proof}
Since
$$
  \deg K_C=2g-2,
$$
we have
$$
  \deg(K_C-D)<0.
$$
A line bundle of negative degree has no
nonzero global sections, so
$$
  h^0(K_C-D)=0.
$$
Thus Riemann-Roch gives
$$
  h^0(D)=\deg D-g+1.
$$
\end{proof}

\begin{exercise}
\label{exercise-p1-sections}
Compute $h^0(\mathbb P^1,\mathcal O(d))$
for $d\geq 0$.
\end{exercise}

\begin{proof}
On $\mathbb P^1$, we have $g=0$.
Also
$$
  \deg \mathcal O(d)=d.
$$
Since $d\geq 0>2g-2=-2$, Riemann-Roch
gives
$$
  h^0(\mathcal O(d))
  =
  d-0+1
  =
  d+1.
$$
Therefore
$$
  h^0(\mathbb P^1,\mathcal O(d))=d+1.
$$
\end{proof}

\begin{exercise}
\label{exercise-elliptic-point}
Let $E$ be an elliptic curve, and let
$p\in E$. Compute
$$
  h^0(E,\mathcal O_E(p)).
$$
\end{exercise}

\begin{proof}
An elliptic curve has genus $1$, and
$$
  K_E\sim 0
$$
(because the standard top-degree
holomorphic form $dz^1 \wedge…\wedge dz^n$
descends to $\mathbb{C}^n/\mathbb{Z}^n$).

Apply Riemann-Roch to $D=p$:
$$
  h^0(p)-h^0(K_E-p)
  =
  \deg p-1+1.
$$
Since $K_E\sim 0$, this becomes
$$
  h^0(p)-h^0(-p)=1.
$$
But
$$
  \deg(-p)<0,
$$
so
$$
  h^0(-p)=0.
$$
Therefore
$$
  h^0(E,\mathcal O_E(p))=1.
$$
\end{proof}

\begin{exercise}
\label{exercise-plane-quartic-genus}
Let $C\subset \mathbb P^2$ be a smooth
plane quartic. Compute $g(C)$,
$K_C$, $\deg K_C$, and $h^0(K_C)$.
\end{exercise}

\begin{proof}
For a smooth plane curve of degree $d$,
$$
  g(C)=\frac{(d-1)(d-2)}2.
$$
Here $d=4$, so
$$
  g(C)=\frac{3\cdot 2}{2}=3.
$$

By adjunction,
$$
  K_C=\mathcal O_C(d-3).
$$
Thus
$$
  K_C=\mathcal O_C(1).
$$

Since $C$ has degree $4$,
$$
  \deg \mathcal O_C(1)=4.
$$
Therefore
$$
  \deg K_C=4.
$$

Finally,
$$
  h^0(K_C)=g(C)=3.
$$
\end{proof}

\begin{exercise}
\label{exercise-plane-quintic}
Let $C\subset \mathbb P^2$ be a smooth
plane quintic. Compute:
$$
  g(C),\quad
  K_C,\quad
  \deg K_C,\quad
  h^0(K_C),
  \quad
  h^0(\mathcal O_C(1)),
$$
$$
  \deg \mathcal O_C(2),
  \quad
  h^0(\mathcal O_C(2)).
$$
Also describe the morphism induced by
$|\mathcal O_C(1)|$.
\end{exercise}

\begin{proof}
Since $d=5$,
$$
  g(C)
  =
  \frac{(5-1)(5-2)}2
  =
  6.
$$

By adjunction,
$$
  K_C=\mathcal O_C(5-3)
  =
  \mathcal O_C(2).
$$

Since
$$
  \deg \mathcal O_C(1)=5,
$$
we get
$$
  \deg K_C
  =
  \deg \mathcal O_C(2)
  =
  10.
$$

Also
$$
  h^0(K_C)=g(C)=6.
$$

The sections of $\mathcal O_C(1)$ are
the restrictions of $X,Y,Z$, so
$$
  h^0(\mathcal O_C(1))=3.
$$

Finally,
$$
  h^0(\mathcal O_C(2))
  =
  h^0(K_C)
  =
  6.
$$

The linear system $|\mathcal O_C(1)|$
induces the original embedding
$$
  C\hookrightarrow \mathbb P^2.
$$
\end{proof}

\section{Surfaces}
\label{section-surfaces}

Main example:

$$
  S=\operatorname{Bl}_p\mathbb P^2.
$$

Let

$$
  H=\pi^*(\text{line}),
  \qquad
  E=\text{exceptional divisor}.
$$

Then

$$
  \operatorname{Pic}(S)
  =
  \mathbb ZH\oplus \mathbb ZE.
$$

The intersection table is

$$
  H^2=1,
  \qquad
  H\cdot E=0,
  \qquad
  E^2=-1,
$$
where
$H$ is the total transform of a line,
$E\simeq \mathbb P^1$ is the exceptional
  curve, and
$E^2=-1$. To check that $E^2 = -1$,
let $L$ be a line through $p$, and pull
it back to the blow-up. We obtain:
total transform = strict transform $+ E$,
that is: $H = \tilde{L} + E$,
which gives $1 = H^2 = L^2 + L\cdot E + E^2
=1+1+E^2$, since the strict
transform does intersect $E$.

The canonical divisor is computed by
$$
K_S=\pi^*K_{\mathbb P^2}+E,
$$
which may be seen by pulling back
a top form on $\mathbb{P}^2$,
and when writing its coordinates
at the blow up, like in the routine
exercises, an extra term appears,
i.e. something of the kind $x dx \wedge dy$
(see Exercise
\ref{exercise-canonical-blow-up}).

Since
$$
K_{\mathbb P^2}=-3H,
$$
we get
$$
K_S=-3H+E.
$$
More generally, if we blow up points
$p_1,\ldots,p_r$, then
$$
 K_S=-3H+E_1+\cdots+E_r.
$$
If $C\subset \mathbb P^2$ is a curve of
degree $d$ and multiplicity $m$ at $p$,
then its strict transform has class

$$
\widetilde C=dH-mE.
$$

For several blown-up points,

$$
\widetilde C
=
dH-m_1E_1-\cdots-m_rE_r.
$$
Slogan:

blowing up adds an exceptional divisor
$E$ with $E^2=-1$, and strict transforms
are computed by subtracting the
multiplicity times $E$.

\begin{exercise}
\label{exercise-canonical-blow-up}
Explain locally why, for the blow-up of
a point on a surface,
$$
  K_S=\pi^*K_X+E.
$$
\end{exercise}

\begin{proof}
Use local coordinates
$$
  \pi(x,u)=(x,y)=(x,xu).
$$
A local generator of the canonical bundle
downstairs is
$$
  dx\wedge dy.
$$
Pull it back:
$$
  \pi^*(dx\wedge dy)
  =
  dx\wedge d(xu).
$$
Since
$$
  d(xu)=u\,dx+x\,du,
$$
we get
$$
  \pi^*(dx\wedge dy)
  =
  dx\wedge (u\,dx+x\,du)
  =
  x\,dx\wedge du.
$$
Thus the pulled-back volume form has one
extra zero along
$$
  E=\{x=0\}.
$$
Zeros of top forms define the canonical
divisor. Hence
$$
  K_S=\pi^*K_X+E.
$$
\end{proof}

\begin{exercise}
\label{exercise-blowup-strict-transform}
Let $C\subset \mathbb P^2$ be a curve
of degree $d$ with multiplicity $m$ at
$p$. What is the class of the strict
transform of $C$ in
$\operatorname{Bl}_p\mathbb P^2$?
\end{exercise}

\begin{proof}
The total transform has class $dH$.
Since $C$ has multiplicity $m$ at $p$,
the exceptional divisor appears with
multiplicity $m$ in the total transform.
Thus
$$
  \pi^*C=\widetilde C+mE.
$$
Therefore
$$
  \widetilde C=dH-mE.
$$
\end{proof}

\begin{exercise}
\label{exercise-blowup-line}
On $S=\operatorname{Bl}_p\mathbb P^2$,
compute the self-intersection of the
strict transform of a line through $p$.
\end{exercise}

\begin{proof}
A line through $p$ has degree $1$ and
multiplicity $1$ at $p$. Hence its
strict transform has class
$$
  H-E.
$$
Therefore
$$
  (H-E)^2
  =
  H^2-2H\cdot E+E^2.
$$
Using
$$
  H^2=1,
  \qquad
  H\cdot E=0,
  \qquad
  E^2=-1,
$$
we get
$$
  (H-E)^2=1-1=0.
$$
\end{proof}

\begin{exercise}
\label{exercise-blowup-conic}
On $S=\operatorname{Bl}_p\mathbb P^2$,
compute the self-intersection of the
strict transform of a conic through $p$
with multiplicity $2$.
\end{exercise}

\begin{proof}
The conic has degree $2$ and multiplicity
$2$ at $p$. Thus its strict transform has
class
$$
  2H-2E.
$$
Hence
$$
  (2H-2E)^2
  =
  4H^2-8H\cdot E+4E^2.
$$
Using the intersection table, this gives
$$
  (2H-2E)^2=4-4=0.
$$
\end{proof}

\begin{exercise}
\label{exercise-two-point-blowup}
Let
$$
  S=\operatorname{Bl}_{p,q}\mathbb P^2,
$$
where $p$ and $q$ are distinct points.
Let $H=\pi^*(\text{line})$, and let
$E_p,E_q$ be the exceptional divisors.
Compute the intersection table and
$K_S$.
\end{exercise}

\begin{proof}
The Picard group is generated by
$$
  H,\quad E_p,\quad E_q.
$$
The intersection table is
$$
  H^2=1,
  \qquad
  E_p^2=E_q^2=-1,
$$
and
$$
  H\cdot E_p=H\cdot E_q=0,
  \qquad
  E_p\cdot E_q=0.
$$
The exceptional divisors are disjoint,
so $E_p\cdot E_q=0$.

The canonical divisor is
$$
  K_S=\pi^*K_{\mathbb P^2}+E_p+E_q.
$$
Since $K_{\mathbb P^2}=-3H$, we get
$$
  K_S=-3H+E_p+E_q.
$$
\end{proof}

\begin{exercise}
\label{exercise-two-point-curves}
On
$$
  S=\operatorname{Bl}_{p,q}\mathbb P^2,
$$
compute the classes and self-intersections
of the following curves:
a line through $p$ but not $q$;
the line through both $p$ and $q$;
a conic through $p$ and $q$.
\end{exercise}

\begin{proof}
A line through $p$ but not $q$ has class
$$
  H-E_p.
$$
Its self-intersection is
$$
  (H-E_p)^2=1-1=0.
$$

The line through both $p$ and $q$ has
class
$$
  H-E_p-E_q.
$$
Its self-intersection is
$$
  (H-E_p-E_q)^2=1-1-1=-1.
$$

A conic through $p$ and $q$ has class
$$
  2H-E_p-E_q.
$$
Its self-intersection is
$$
  (2H-E_p-E_q)^2=4-1-1=2.
$$
\end{proof}

\begin{exercise}
\label{exercise-minus-one-curves}
Find the basic $(-1)$-curves on
$$
  \operatorname{Bl}_{p,q}\mathbb P^2.
$$
\end{exercise}

\begin{proof}
The exceptional curves themselves have
self-intersection $-1$:
$$
  E_p^2=E_q^2=-1.
$$
Also, the strict transform of the line
through both points has class
$$
  H-E_p-E_q,
$$
and
$$
  (H-E_p-E_q)^2=-1.
$$
Thus the basic $(-1)$-curves are
$$
  E_p,\qquad E_q,\qquad H-E_p-E_q.
$$
\end{proof}


\begin{exercise}
\label{exercise-quartic-surface}
Let $X\subset \mathbb P^3$ be a smooth
quartic surface. Compute $\dim X$,
$\deg X$, $K_X$, and $H^2$, where
$H=\mathcal O_X(1)$. Classify $X$.
\end{exercise}

\begin{proof}
Since $X$ is a hypersurface in
$\mathbb P^3$, we have
$$
  \dim X=2.
$$
Its degree is
$$
  \deg X=4.
$$
By adjunction,
$$
  K_X
  =
  \mathcal O_X(4-3-1)
  =
  \mathcal O_X.
$$
Also
$$
  H^2=\deg X=4.
$$
Thus $X$ is a smooth surface with trivial
canonical bundle. Moreover, smooth
complete intersections in projective
space satisfy $H^1(X,\mathcal O_X)=0$.
Hence $X$ is a K3 surface.
\end{proof}

\begin{exercise}
\label{exercise-two-three-surface}
Let
$$
  X=(2,3)\subset \mathbb P^4
$$
be a smooth complete intersection.
Compute $\dim X$, $\deg X$, $K_X$,
$H^2$, $K_X^2$, and
$h^0(X,\mathcal O_X(1))$. Classify $X$.
\end{exercise}

\begin{proof}
The codimension is $2$, hence
$$
  \dim X=4-2=2.
$$
The degree of a complete intersection is
the product of the degrees:
$$
  \deg X=2\cdot 3=6.
$$
The normal bundle is
$$
  N_{X/\mathbb P^4}
  =
  \mathcal O_X(2)\oplus \mathcal O_X(3).
$$
Adjunction gives
$$
  K_X
  =
  \mathcal O_X(2+3-4-1)
  =
  \mathcal O_X.
$$
Thus
$$
  K_X^2=0.
$$
Also
$$
  H^2=\deg X=6.
$$

The global sections of $\mathcal O_X(1)$
are the linear forms on $\mathbb P^4$
restricted to $X$. No nonzero linear form
vanishes on all of $X$, because the ideal
of $X$ is generated in degrees $2$ and
$3$. Therefore
$$
  h^0(X,\mathcal O_X(1))
  =
  h^0(\mathbb P^4,\mathcal O(1))
  =
  5.
$$
The linear system $|\mathcal O_X(1)|$
gives the original embedding
$$
  X\hookrightarrow \mathbb P^4.
$$

Since $K_X\simeq \mathcal O_X$ and
$H^1(X,\mathcal O_X)=0$, this is a
complete-intersection K3 surface.
\end{proof}

\begin{exercise}
\label{exercise-two-two-two-surface}
Let
$$
  X=(2,2,2)\subset \mathbb P^5
$$
be a smooth complete intersection.
Compute $\dim X$, $\deg X$, $K_X$,
$H^2$, $K_X^2$, and
$h^0(X,\mathcal O_X(1))$. Classify $X$.
\end{exercise}

\begin{proof}
The codimension is $3$, hence
$$
  \dim X=5-3=2.
$$
The degree is
$$
  \deg X=2\cdot 2\cdot 2=8.
$$
By adjunction,
$$
  K_X
  =
  \mathcal O_X(2+2+2-5-1)
  =
  \mathcal O_X.
$$
Hence
$$
  K_X^2=0.
$$
Also
$$
  H^2=\deg X=8.
$$

Again, no nonzero linear form vanishes
on all of $X$, since the ideal is
generated in degree $2$. Therefore
$$
  h^0(X,\mathcal O_X(1))
  =
  h^0(\mathbb P^5,\mathcal O(1))
  =
  6.
$$
The linear system $|\mathcal O_X(1)|$
gives the original embedding
$$
  X\hookrightarrow \mathbb P^5.
$$

Thus $X$ is a complete-intersection
K3 surface. (For some reason it turns out
that complete intersections have
$h^1(\mathcal{O}_X)=0$).
\end{proof}

\begin{exercise}
\label{exercise-surface-intersection}
Let $X\subset \mathbb P^n$ be a
projective surface, and let
$H=\mathcal O_X(1)$. Show that
$$
  H^2=\deg X.
$$
Deduce that
$$
  \mathcal O_X(a)\cdot\mathcal O_X(b)
  =
  ab\deg X.
$$
\end{exercise}

\begin{proof}
The class $H$ is the hyperplane class
restricted to $X$. Thus $H^2$ is the
intersection of $X$ with two general
hyperplanes. By definition, this number
is the degree of $X$. Hence
$$
  H^2=\deg X.
$$
Since
$$
  \mathcal O_X(a)=aH,
  \qquad
  \mathcal O_X(b)=bH,
$$
we get
$$
  \mathcal O_X(a)\cdot\mathcal O_X(b)
  =
  (aH)\cdot(bH)
  =
  abH^2
  =
  ab\deg X.
$$
\end{proof}

\section{Schemes}
\label{section-shemes}

For a definition of presheaf,
see Categories, Definition
\ref{categories-definition-presheaf}.

\begin{definition}
\label{definition-sheaf}
Let $X$ be a topological space.
\begin{enumerate}
\item A {\it sheaf $\mathcal{F}$ of sets on
$X$} is a presheaf
of sets which satisfies the following
additional property: Given
any open covering $U = \bigcup_{i \in I} U_i$
and any collection
of sections $s_i \in \mathcal{F}(U_i)$, $i
\in I$ such that
$\forall i, j\in I$
$$
s_i|_{U_i \cap U_j} = s_j|_{U_i \cap U_j}
$$
there exists a unique section $s \in
\mathcal{F}(U)$ such that
$s_i = s|_{U_i}$ for all $i \in I$.
\item A {\it morphism of sheaves of sets} is
simply a
morphism of presheaves of sets.
\item The category of sheaves of sets on $X$
is denoted
$\Sh(X)$.
\end{enumerate}
\end{definition}


\medskip\noindent
Let $X$ be a topological space. Let $x \in X$
be a point.
Let $\mathcal{F}$ be a presheaf of sets on
$X$.
The {\it stalk of $\mathcal{F}$ at $x$} is
the set
$$
\mathcal{F}_x
=
\colim_{x\in U} \mathcal{F}(U)
$$
where the colimit is over the set of open
neighbourhoods
$U$ of $x$ in $X$. The set of open
neighbourhoods is
partially ordered by (reverse) inclusion:
We say $U \geq U' \Leftrightarrow U \subset
U'$.
The transition maps in the system are
given by the restriction maps of
$\mathcal{F}$.
See Categories, Section
\ref{categories-section-posets-limits}
for notation and terminology regarding
(co)limits over systems.
Note that the colimit is a directed colimit.
Thus it is easy to describe $\mathcal{F}_x$.
Namely,
$$
\mathcal{F}_x
=
\{
(U, s)
\mid
x\in U, s\in \mathcal{F}(U)
\}/\sim
$$
with equivalence relation given by $(U, s)
\sim (U', s')$ if and only if
there exists an open $U'' \subset U \cap U'$
with $x \in U''$ and
$s|_{U''} = s'|_{U''}$. Given a pair $(U, s)$
we sometimes denote
$s_x$ the element of $\mathcal{F}_x$
corresponding to the equivalence
class of $(U, x)$. We sometimes use the
phrase
``image of $s$ in $\mathcal{F}_x$'' to denote
$s_x$.
For example, given two pairs $(U, s)$ and
$(U', s')$ we sometimes
say ``$s$ is equal to $s'$ in
$\mathcal{F}_x$'' to indicate
that $s_x = s'_x$. Other authors use the
terminology
``germ of $s$ at $x$''.

\medskip\noindent
Now let's discuss closed immersions.

\noindent
Closed immersions of ringes spaces are in
\href{https://stacks.math.columbia.edu/tag/%
01C1}{01C1}.

\begin{definition}
\label{definition-closed-immersion}
A map of ringed spaces
$i:(Z,\mathcal{O}_Z) \to (X,\mathcal{O}_X)$
is a {\it closed immersion}
if $i$ is injective,
its image is closed (I think these two are
``closed immersion of topological spaces'')
and the induced map $\mathcal{O}_X \to
i_*\mathcal{O}_Z$
is surjective; denote its kernel by
$\mathcal{I}$.
To make sure that if $(X,\mathcal{O}_X)$ is a
scheme
then so is $(Z,\mathcal{O}_Z)$ with add the
extra condition
that the $\mathcal{O}_X$-module $\mathcal{I}$
is
locally generated by sections.
\end{definition}

\medskip\noindent
Now let's discuss restricting sheaves to
subschemes.

Upshot about restricting sheaves to
subschemes:
the sheaves on $Z$ correspond essentially
(in the categorical sense)
to sheaves on $X$ with support on $Z$.

More explicitly, from
\href{https://stacks.math.columbia.edu/tag/%
01QY}{01QY}:
the pushforward of the inclusion (a closed
immersion of schemes)
$i:Z \to X$, with kernel on induced map
$\mathcal{O}_Z \to \mathcal{O}_X$
denoted by $I$, is an exact, fully faithful
functor
$$
i_*:\QCoh(\mathcal{O}_Z) \to
\QCoh(\mathcal{O}_X)
$$
with essential image the quasi-coherent
$\mathcal{O}_X$ modules
$\mathcal{G}$ such that
$\mathcal{I}\mathcal{G}=0$.

It makes sense to think of this operation as
$\mathcal{F}|_Z=i^*\mathcal{F}=\mathcal{F}
\otimes\mathcal{O}_Z$
for $\mathcal{F}\in\QCoh(X)$.
Indeed: tensoring by the ideal sheaf would
kill the information on $Z$,
and tensoring with the structure sheaf does
the opposite
--- it's the sheaves on $X$ that are no fun
outside $Z$,
so basically sheaves on $Z$.

\medskip\noindent
Here's a more detailed discussion:

In understanding the cohomology of sheaves
when
we restrict them to submanifolds (see Complex
Geometry
Exercise
\ref{complex-geometry-exercise-%
deformations-curve-in-k3},
I found Stacks Project
\href{https://stacks.math.columbia.edu/tag/%
01AW}{01AW}
section on closed immersions and abelian
sheaves.
First we have a few propositions stating that
sheaves on the subscheme $Z \subset X$ can be
pushed
to sheaves on $X$.

And then there's the converse:
what I would like to call the restriction
functor $\mathcal{H}$
but rather is called the {\it sub-module of
sections with support in $Z$}.
In Remark
\href{https://stacks.math.columbia.edu/tag/%
01AY}{01AY}
(and then in Remark
\href{https://stacks.math.columbia.edu/tag/%
0G6N}{0G6N}
for ringed spaces; these seem to be analogue
constructions…)
we see that for an abelian sheaf
$\mathcal{F}$ on $X$
we can define an abelian sheaf
$\mathcal{H}_Z(\mathcal{F})$
which although is a sheaf on $X$
we can just think it's a sheaf on $Z$ -
that's beacause
$$
\mathcal{H}_Z(\mathcal{F})(U)
=\{s \in \mathcal{F}(U))| \text{the support
of $s$ is contained in }Z\cap U\}.
$$
We can just think that it's the functions
that vanish outside the subvariety.
Which doesn't really make sense:
a function vanishing outside the submanifold
will most likely
vanish also in the submanifold! Just by
continuity,
the submanifold is in the closure of the open
set!
I guess what I mean to say is that I don't
mind if
the functions do not vanish outside $Z$…
I just want to consider them as functions on
$Z$…

\medskip\noindent
Here's another attempt in understanding how
to restrict sheaves to submanifolds.
First and foremost, lacking a reference this,
I understand the symbol $\mathcal{F}|_Z$ for
a sheaf
on $X$ and a subscheme  $Z \subset X$ to mean
the
pullback of $\mathcal{F}$ by the inclusion.
Then look at
\href{https://stacks.math.columbia.edu/tag/%
01QY}{01QY}
(more basically, as above,
\href{https://stacks.math.columbia.edu/tag/%
01AW}{01AW}…
in fact there appear to be several version of
this statement)
to recall that sheaves of the subscheme $Z$
should correspond
to sheaves on $X$ with support on $Z$.
{\bf So essentially $\mathcal{F}|_Z$ should
be a version of $\mathcal{F}$
but with support on $Z$.}
The key point here is that this is non other
than
$\mathcal{F} \otimes \mathcal{O}_Z$.
Why? Because when multiply with the ideal
sheaf of $Z$ we obtain
$(\mathcal{F} \otimes
\mathcal{O}_Z)\mathcal{I}=0$
since sections of $\mathcal{I}$
vanish in
$\mathcal{O}_Z=\mathcal{O}_X/\mathcal{I}$.
That is, the support of $\mathcal{F} \otimes
\mathcal{O}_Z$ indeed
contained in $Z$: if the sections become zero
when
multiplied by sections of $\mathcal{I}$
(sections of $\mathcal{I}$ are only nonzero
outside $Z$)
it means that all the action was inside $Z$.
that outside $Z$

I think the upshot here is that structure
sheaves
act as some sort of ``contour'' for the
geometric object:
they vanish {\it outside} the geometric
object,
they tell us what is the shape of the object
by describing its surroundings, rather than
its interior.
This is exactly how algebraic/analytic sets
are defined,
and this why, when we want to restrict a
sheaf
to a subscheme, we tensor with the structure
sheaf:
tensor by the ideal sheaf would kill the
information on $Z$,
and tensoring with the structure sheaf does
the opposite.














\medskip\noindent
Now let's dicuss Proj of a graded ring.

\noindent
(Stacks Project
\href{https://stacks.math.columbia.edu/tag/%
01M3}{01M3},
Proj of a graded ring.)
This section may be taken as the preamble to
the question: what is projective space
anyway?
Yes, projective space is Proj of a graded
ring $S$.
And yes, projective space is a scheme.
So in particular it's locally affine.

(Stacks Project
\href{https://stacks.math.columbia.edu/tag/%
01MX}{01MX},
Functoriality of Proj.)
This is reminiscent of the olden days.
The point is that Proj is covariant,
i.e. for rings  $S' \subset S$ we have
$\text{Proj}S \subset \text{Proj}S$
(under the right conditions,
which is the point of this section),
meaning that, indeed,
{\bf Projs are projective schemes}.






\medskip\noindent
Noe let's discuss invertible sheaves on Proj
(twists!).

\noindent
As for the construction of the twisted
sheaves:
for a projective space $\text{Proj} S$,
we just compute the twisted $S$ module
$S(m)$ by the formula $S(m)_d=S_{m+d}$
and turn that into a sheaf
$\widetilde{S(m)}$.

Here's the construction of the twisted
sheaves $\mathcal{O}_X(n)$.
The point is that there is a good way
to pass from an $S$-module $M$ to an
$\mathcal{O}_X$-module, called
$\widetilde{M}$.
So basically you just define the twists
by $M(d)_n=M_{n+d}$ and apply that
construction.

Pick an element of degree $d$.
Then think of the module generated by this
element.
The least degree piece of such a module
is the degree $d$ piece of the original
module,
that is we have $M(-d)$.











\medskip\noindent
Lost note on logarithmic derivations:

\noindent
Just to record a notion of logarithmic
derivation
from Vladimiro Benedetti, Danielle Faenzi and
Simone.

Let $F$ be a homogneous polynomial of degree
$d$.
The {\it logarithmic derivations} of $F$ are
$$
\text{Der}_U(-\text{log}F)
=\{\partial \in \text{Der}_U
:\partial F \in \left\langle{}F\right\rangle{}_U\}
$$
where $U=\mathbb{C}[x_1,…,x_e]$.










\medskip\noindent
Lost notes on gonality.

\noindent
From a talk by Juliana Coelho at Bandoleiros
2025:

A curve $C$ is {\it $k$-gonal} if there is a
map $C \to \mathbb{P}^1$ of degree $k$.

A curve is {\it stable} if it is nodal and
$\text{Aut}(C)$ is finite.
Equivalently, the rational components of $C$
meet the
rest of the curve in at least three points.

We may study the moduli of $k$-stable curves,
or the space of $k$-gonal structures.

\section{Reduced schemes}
\label{section-reduced}

\noindent
So far a reduced scheme,
for me,
is just a scheme where there are
no ``repetitions'':
$(x^2)$ is not reduced because
it gives the same information as $(x)$.
Algebraically this means that the
coordinate ring is reduced,
which is defined by asking
that there are no nilpotent elements,
i.e. elements such that some power vanishes.
(There's no definition of reduced ring in
Stacks Project
since it's taken as part of the basic
algebraic knowledge).

The definition is on stalks but the first
lemma
shows it's equivalent to define it on any
open set.
This is in virtue of some ``injectivity''
lemma
from the ring at $U$ to the product of all
stalks parametrized in $U$;
essentially that if a section vanishes when
projected to all
the stalks then it's zero.

\begin{definition}
\label{definition-reduced}
Let $X$ be a scheme. We say $X$ is {\it
reduced} if every local ring
$\mathcal{O}_{X, x}$ is reduced.
\end{definition}


\begin{proof}
Assume that $X$ is reduced.
Let $f \in \mathcal{O}_X(U)$ be a section
such that $f^n = 0$.
Then the image of $f$ in $\mathcal{O}_{U, u}$
is zero for
all $u \in U$. Hence $f$ is zero, see
Sheaves, Lemma
\ref{sheaves-lemma-sheaf-subset-stalks}.
Conversely, assume that $\mathcal{O}_X(U)$ is
reduced
for all opens $U$. Pick any nonzero element
$f \in \mathcal{O}_{X, x}$.
Any representative $(U, f \in
\mathcal{O}(U))$  of $f$ is nonzero and
hence not nilpotent. Hence $f$ is not
nilpotent in $\mathcal{O}_{X, x}$.
\end{proof}

\section{Flat morphisms}
\label{section-flat-morphisms}

\noindent
The essential technical property
for for defining flatness
is the preservation of exact sequences.
Right-exactness is true in general;
it follows from currying in category
of commutative rings.
But the functor $-\otimes_R N$ does not
preserve injectivity of maps---
that's the point of flatness.

\begin{lemma}[Internal Hom for R-modules]
\label{lemma-hom-from-tensor-product}
For any three $R$-modules $M, N, P$,
$$
\Hom_R(M \otimes_R N, P) \cong \Hom_R(M,
\Hom_R(N, P))
$$
\end{lemma}

\begin{proof}
An $R$-linear map $\hat{f}\in \Hom_R(M
\otimes_R N, P)$ corresponds to an
$R$-bilinear map $f : M \times N \to P$. For
each $x\in M$ the mapping $y\mapsto f(x, y)$
is $R$-linear by the universal
property. Thus $f$ corresponds to a
map $\phi_f : M \to \Hom_R(N, P)$. This map
is $R$-linear since
$$
\phi_f(ax + y)(z) =
f(ax + y, z) = af(x, z)+f(y, z) =
(a\phi_f(x)+\phi_f(y))(z),
$$
for all $a \in R$, $x \in M$, $y \in M$ and
$z \in N$. Conversely, any
$f \in \Hom_R(M, \Hom_R(N, P))$ defines an
$R$-bilinear
map $M \times N \to P$, namely $(x, y)\mapsto
f(x)(y)$.
So this is a natural one-to-one
correspondence between the
two modules
$\Hom_R(M \otimes_R N, P)$ and $\Hom_R(M,
\Hom_R(N, P))$.
\end{proof}

\begin{lemma}[Tensor product is right exact]
\label{lemma-tensor-product-exact}
Let
$$
\begin{aligned}
M_1\xrightarrow{f} M_2\xrightarrow{g} M_3 \to
0
\end{aligned}
$$
be an exact sequence of $R$-modules and
homomorphisms, and let $N$ be any
$R$-module. Then the sequence
\begin{equation}
\label{equation-2ndex}
M_1\otimes N\xrightarrow{f \otimes 1}
M_2\otimes N \xrightarrow{g \otimes 1}
M_3\otimes N \to 0
\end{equation}
is exact. In other words, the functor $-
\otimes_R N$ is
{\it right exact}, in the sense that
tensoring
each term in the original right exact
sequence preserves the exactness.
\end{lemma}

\begin{proof}
For every $R$-module $P$
we apply the functor $\Hom(-, \Hom(N, P))$ to
the first exact
sequence. We obtain
$$
0 \to
\Hom(M_3, \Hom(N, P)) \to
\Hom(M_2, \Hom(N, P)) \to
\Hom(M_1, \Hom(N, P))
$$
which is exact by Lemma \ref{lemma-hom-exact}
(1).
By Lemma \ref{lemma-hom-from-tensor-product}
this becomes the sequence
$$
0 \to \Hom(M_3 \otimes N, P) \to
\Hom(M_2 \otimes N, P) \to \Hom(M_1 \otimes
N, P)
$$
which is therefore also exact. Then using
Lemma \ref{lemma-hom-exact} (1) again, we
arrive at the desired exact sequence.
\end{proof}

\begin{remark}
\label{remark-tensor-product-not-exact}
However, tensor product does NOT preserve
exact sequences in general.
In other words, if $M_1 \to M_2 \to M_3$ is
exact, then it is not necessarily true that
$M_1 \otimes N \to M_2 \otimes N \to M_3
\otimes N$
is exact for arbitrary $R$-module $N$.
\end{remark}

\begin{example}
\label{example-tensor-product-not-exact}
Consider the injective map $2 : \mathbf{Z}\to
\mathbf{Z}$
viewed as a map of $\mathbf{Z}$-modules.
Let $N = \mathbf{Z}/2$. Then the induced map
$\mathbf{Z} \otimes \mathbf{Z}/2 \to
\mathbf{Z} \otimes \mathbf{Z}/2$
is NOT injective. This is because for
$x \otimes y\in \mathbf{Z} \otimes
\mathbf{Z}/2$,
$$
(2 \otimes 1)(x \otimes y) = 2x \otimes y = x
\otimes 2y = x \otimes 0 = 0
$$
Therefore the induced map is the zero map
while $\mathbf{Z} \otimes N\neq 0$.
\end{example}

\begin{definition}
\label{definition-flat-module}
For $R$-modules $N$, if the
functor $-\otimes_R N$ is exact, i.e.
tensoring
with $N$ preserves all exact
sequences, then $N$ is said to be {\it flat}
$R$-module.
We will discuss this later in Section
\ref{section-flat}.
\end{definition}

\medskip\noindent
Epiphany:
in the category of commutative rings
pushout is tensor product.
So think of $\Spec $ as a functor
from $\mathit{CRing}^{\text{op}}$ to $\Sch$,
then pushout goes to pullback,
and what's an example of a pullback?
Fibre!
So, the coordinate ring of a fiber
is essentially given by the residue field
at the point that parametrizes it!
(tensored with the coordinate ring
of the deformation space).

There is a lot of information on Stacks
Project about flatness.
It looks like the heart of the concept is
captured in the commutative-algebraic notion
of preserving
exact sequences:

\begin{definition}
\label{definition-flat}
Let $R$ be a ring.
\begin{enumerate}
\item An $R$-module $M$ is called {\it flat}
if whenever
$N_1 \to N_2 \to N_3$ is an exact sequence of
$R$-modules
the sequence $M \otimes_R N_1 \to M \otimes_R
N_2 \to M \otimes_R N_3$
is exact as well.
\item An $R$-module $M$ is called {\it
faithfully flat} if the
complex of $R$-modules
$N_1 \to N_2 \to N_3$ is exact if and only if
the sequence $M \otimes_R N_1 \to M \otimes_R
N_2 \to M \otimes_R N_3$
is exact.
\item A ring map $R \to S$ is called {\it
flat} if
$S$ is flat as an $R$-module.
\item A ring map $R \to S$ is called {\it
faithfully flat} if
$S$ is faithfully flat as an $R$-module.
\end{enumerate}
\end{definition}

\medskip\noindent
Recall that a module $M$ over a ring $R$ is
{\it flat} if the functor
$-\otimes_R M : \text{Mod}_R \to
\text{Mod}_R$ is exact. A ring map
$R \to A$ is said to be {\it flat} if $A$ is
flat as an $R$-module.

\section{Hilbert polynomial}
\label{section-Hilbert-polynomial}

\noindent
The following lemma will be made obsolete by
the more general
Lemma
\ref{lemma-numerical-polynomial-from-euler}.

\begin{lemma}
\label{lemma-hilbert-polynomial}
Let $k$ be a field. Let $n \geq 0$. Let
$\mathcal{F}$ be a coherent sheaf
on $\mathbf{P}^n_k$. The function
$$
d \longmapsto \chi(\mathbf{P}^n_k,
\mathcal{F}(d))
$$
is a polynomial.
\end{lemma}

\begin{proof}
We prove this by induction on $n$. If $n =
0$, then
$\mathbf{P}^n_k = \Spec(k)$ and
$\mathcal{F}(d) = \mathcal{F}$.
Hence in this case the function is constant,
i.e., a polynomial
of degree $0$. Assume $n > 0$. By
Lemma
\ref{lemma-euler-chara-exten-base-field}
we may assume $k$ is infinite. Apply
Lemma \ref{lemma-exact-sequence-induction}.
Applying Lemma
\ref{lemma-euler-characteristic-additive}
to the twisted sequences
$0 \to \mathcal{F}(d - 1) \to \mathcal{F}(d)
\to i_*\mathcal{G}(d) \to 0$
we obtain
$$
\chi(\mathbf{P}^n_k, \mathcal{F}(d)) -
\chi(\mathbf{P}^n_k, \mathcal{F}(d - 1)) =
\chi(H, \mathcal{G}(d))
$$
See Remark
\ref{remark-exact-seque-induc-cohom}.
Since $H \cong \mathbf{P}^{n - 1}_k$
by induction the right hand side is a
polynomial.
The lemma is finished by noting that any
function
$f : \mathbf{Z} \to \mathbf{Z}$ with the
property that the map
$d \mapsto f(d) - f(d - 1)$ is a polynomial,
is itself a polynomial.
We omit the proof of this fact (hint: compare
with
Algebra, Lemma
\ref{algebra-lemma-numerical-polynomial}).
\end{proof}

\begin{definition}
\label{definition-hilbert-polynomial}
Let $k$ be a field. Let $n \geq 0$. Let
$\mathcal{F}$ be a coherent sheaf
on $\mathbf{P}^n_k$. The function
$d \mapsto \chi(\mathbf{P}^n_k,
\mathcal{F}(d))$ is called the
{\it Hilbert polynomial} of $\mathcal{F}$.
\end{definition}

\noindent
The Hilbert polynomial has coefficients in
$\mathbf{Q}$ and not
in general in $\mathbf{Z}$. For example the
Hilbert polynomial
of $\mathcal{O}_{\mathbf{P}^n_k}$ is
$$
d \longmapsto {d + n \choose n} =
\frac{d^n}{n!} + \ldots
$$
This follows from the following lemma and the
fact that
$$
H^0(\mathbf{P}^n_k,
\mathcal{O}_{\mathbf{P}^n_k}(d)) = k[T_0,
\ldots, T_n]_d
$$
(degree $d$ part) whose dimension over $k$ is
${d + n \choose n}$.

\begin{lemma}
\label{lemma-hilbert-polynomial-H0}
Let $k$ be a field. Let $n \geq 0$. Let
$\mathcal{F}$ be a coherent sheaf
on $\mathbf{P}^n_k$ with Hilbert polynomial
$P \in \mathbf{Q}[t]$.
Then
$$
P(d) = \dim_k H^0(\mathbf{P}^n_k,
\mathcal{F}(d))
$$
for all $d \gg 0$.
\end{lemma}

\begin{proof}
This follows from the vanishing of cohomology
of high enough twists
of $\mathcal{F}$. See
Cohomology of Schemes,
Lemma
\ref{coherent-lemma-coherent-projective}.
\end{proof}


\medskip\noindent
For completeness I include earlier notes
from \cite{har} on the matter.

The fact that $M$ is finitely generated is
what makes the following two
definitions make sense.

\begin{definition}
\label{definition-Hilbert-function}
The {\it Hilbert function} of a finitely
generated graded $S=k[x_0,\ldots,x_r]$
-module $M$ is
$$
H_M(d)=\dim_kM_d
$$
\end{definition}

\begin{definition}
\label{definition-sysygy}
Define $F_0$ to be the free $S$-module on the
generators of $M$. Elements in the
 kernel $M_1$ of the inclusion are called
 {\it sysygies}. By Hilbert's basis
theorem, $M_1$ is also finitely generated, so
we may choose a set of generators
and repeat this process.
\end{definition}

\begin{theorem}[Hilbert Syzygy Theorem]
\label{theorem-Hilbert-syzygy}
\begin{reference}
\cite[Theorem 1.1]{sys}
\end{reference}
Any finitely generated $S$-module $M$ has a
finite graded free resolution
$$
\xymatrix{
0\ar[r]&F_m\ar[r]^{\varphi_m}&\ar[r]&F_{m-q}
\ar[r]&\cdots\ar[r]&
F_1\ar[r]^{\varphi_1}&F_0
}
$$
Moreover, we may take $m\leq r+1$, the number
of variables in $S$.
\end{theorem}

\begin{lemma}
\label{lemma-Hilbert-function}
Suppose that $S=k[x_0,\ldots,x_r]$ is a
polynomial ring. If the graded
$S$-module $M$ has finite free resolution
$$
\xymatrix{
0\ar[r]&F_m\ar[r]^{\varphi_m}&F_{m-1}
\ar\cdots[r]&F_1\ar[r]^{\varphi_1}&
\ar[r]&F_0
}
$$
with each $F_i$ a finitely generated free
module,
$F_i=\bigoplus_{j}S(-a_{i,j})$, then
\begin{equation}
\label{equation-Hilbert-function}
H_M(d)=\sum_{i=0}(-1)^i\sum_{j}\binom{r+d-
a_{i,j}}{r}
\end{equation}
\end{lemma}

\begin{lemma}
\label{lemma-hilbe-funct-becom-polyn}
There is a polynomial $P_M(d)$ called the
{\it Hilbert polynomial} such that, if
$M$ has free resolution as above, then
$P_M(d)=H_M(d)$ for
$d\geq\text{max}_{i,j}\{a_{i,j}-r\}$.
\end{lemma}

\begin{proof}
When $d$ satisfies this bound then the
binomial coefficients in Eq.
\ref{equation-Hilbert-function} are
polynomials of degree $r$ in $d$.
\end{proof}

\begin{theorem}[Hilbert-Serre]
\label{theorem-Hilbert-Serre}
\begin{reference}
\cite[I, Theorem 7.5]{hart}
\end{reference}
Let $M$ be a finitely generated graded
$S=k[x_0,\ldots,x_n]$. Then there exists
a unique polynomial $p_M$ such that
$p_M(\ell)=\dim S_\ell$ for large enough
$\ell$.
\end{theorem}

\begin{definition}
\label{definition-Hilbert-polynomial}
\begin{reference}
\cite[I, p. 52]{har}
\end{reference}
The polynomial $P_M$ of Hilbert-Serre Theorem
\cite{Hilbert-Serre} is the {\it
Hilbert polynomial} of the finitely generated
$k[x_0,\ldots,x_n]$-module $M$.
\end{definition}

\begin{definition}
\label{definition-degre-proje-varie}
\begin{reference}
\cite[p. 52]{hart}
\end{reference}
If $Y\subset \mathbb{P}^n$ is an algebraic
set of dimension $r$, we define the
{\it degree of $Y$} to be $r!$ times the
leading coefficient of the Hilbert
polynomial of the homogeneous coordinate ring
$S(Y)$.
\end{definition}

\begin{exercise}
\label{exercise-very-ample-bundl-self-inte}
\begin{reference}
\cite[V, Exercise 1.2]{har}
\end{reference}
Let $H$ be a very ample divisor on the
surface $X$, corresponding to a
projective embedding
$X\subseteq\mathbb{P}^N$. If we write the
Hilbert
polynomial of $X$ as
$P(z)=\frac{1}{2}az^2+bz+c$, show that
$a=H^2$,
$b=\frac{1}{2}H^2+1-\pi$, where $\pi$ is the
genus of a nonsingular curve
representing $H$, and $c=1+p_a$.
\end{exercise}

\section{Nakai-Moishezon Criterion}
\label{section-Nakai-Moishezon-criterion}

\begin{theorem}[Nakai-Moishezon Criterion]
\label{theorem-Nakai-Moishezon-criterion}
\begin{reference}
\cite[V, Theorem 1.10]{hart}
\end{reference}
A divisor $D$ on the surface $X$ is ample if
and only if $D^2>0$ and $D.C>0$ for
all irreducible curves $C$ in $X$.
\end{theorem}

\begin{proof}
The direct implication is easy: since $D$ is
ample,  $mD$ is very ample for some
$m$, so that $m^2D^2$ is the
self-intersection number of $mD$. By exercise
\ref{exercise-very-ample-bundl-self-inte},
 $D^2$ is the leading coefficient of the
 Hilbert polynomial of $X$ as a
subscheme of $\mathbb{P}^n$. This means that
$D^2$ is twice the leading
coefficient of the Hilbert polynomial of a
projective variety for large enough
$m$, so that it must be a positive number
(it's the dimension of one of the
graded components of the coordinate ring of
the surface).
\end{proof}

\section{Hilbert scheme}
\label{section-Hilbert-scheme}

{\bf Upshot \cite[p. 6]{HarrMorr}.}
We wish to parametrize subschemes of a
projective space (or
perhaps a more general scheme?). Since there
are too many such subschemes we
restrict ourselves to schemes with a given
Hilbert polynomial, since the latter
``encodes the most important numerical
invariants of schemes''. The Hilbert
scheme is introduced via a theorem by
Grothendieck
as the object that represents the functor
$\mathbf{Hilb}_{P,r}$
that maps a reduced scheme $B$ to the
set of proper flat families
$$
\xymatrix{
\mathcal{X}\ar@{^{(}->}[r]^i\ar[rd]&
\mathbb{P}^r
\times B \ar[d]^{\pi_B}\\
& B
}
$$
with $\mathcal{X}$ having Hilbert polynomial
$P$.

\begin{theorem}[Grothendieck, '66]
\label{theorem-Grothendieck}
The functor $\mathbf{Hilb}_{P,r}$ is
representable by a projective scheme
$\mathcal{H}_{P,r}$.
\end{theorem}

SEE Hilbert schemes of subschemes.

\section{Deformation theory}
\label{section-deformation-theory}

\noindent
\cite{Sernesi-Overview} has a great way of
motivating deformation
theory via the moduli problem.
Here let's just put the following important
meaning for
``deformation theory'' (as opposed to
studying flat families!):

\begin{quote}
``{\it Deformation theory} is the study of
infinitesimal deformaitons
as a tool to understand the local structure
of the moduli space.
The goal is to be able to describe the
restriction
of the universal family to a small
neighborhood of $m \in \mathcal{M}$,
or, more precisely, its restriction to the
germ of $M$ at $m$.

What is interesting here is that we can study
order and infinitesimal
deformations even though the functor $F$ is
not representable or simply we don’t
yet know it is. This is the most frequent
case. Such an investigation will
reveal the infinitesimal properties at $[X]$
of a yet unknown global structure
on $M$ which will be hopefully understood at
a subsequent stage of the
investigation.  In order words it turns out
to be possible and convenient to
separate the {\it global moduli problem} from
the {\it local moduli problem},
and deformation theory studies the latter,
with the purpose of constructing a
family of deformations of a given object
parametrized by the spectrum of a local
ring, and having properties as close as
possible to a universal property.''
\end{quote}




\noindent
I start by reading Stacks Project.

\medskip\noindent
The first notion is {\it thickening of ringed
spaces},
which I ultimately think of as a closed
subscheme
$(X,\mathcal{O}_X) \to (X',\mathcal{O}_{X'})$
with a nilpotent ideal sheaf,
which in an imprecise way means
that $\mathcal{I}_X^n=0$ for some $n$,
and in a precise way its given on sections.

\medskip\noindent
The following definition is from
[lucas-defos], which in turn comes from
\cite{Sernesi-deformations}

\begin{definition}
\label{definition-deformation}
Let $X$ be an algebraic $\mathbb{C}$-scheme.
\begin{enumerate}
\item A {\it deformation} of $X$ is a
Cartesian diagram $\xi$
$$
\xymatrix{
X\ar[r]\ar[d]&\mathcal{X}\ar[d]^{\pi}\\
\Spec \mathbb{C}\ar[r]^s&S
}
$$
where $\pi$ is a flat surjective morphism of
algebraic $\mathbb{C}$-schemes and
$S$ is connected. (Recall that flatness
accounts for ``continuity''.)
\item A {\it local deformation} of $X$ is a
deformation $\xi$ where
$S=\Spec A$ for $A$ a noetherian local
$\mathbb{C}$ algebra with residue
field $\mathbb{C}$.
\item An {\it infinitesimal deformation of
$X$} is a local deformation with $A$
an artinian local $\mathbb{C}$-algebra with
residue field $\mathbb{C}$. $X$ is
called {\it rigid} if all infinitesimal
deformations are trivial.
\item An {\it inifinitesimal deformation of
order $n$} is an infinitesimal
deformation when $S=\Spec
(\mathbb{C}[t]/(t^{n+1})$.
\end{enumerate}
\end{definition}

{\bf Upshot.} An interpretation of the
so-called {\it dual numbers}
$k[t]/(t^2)$ (see
\cite{Hartshorne-deformation}) as the tangent
space of
something is thinking of Taylor polynomials:
after quotienting by $t^2$ we loose
the tails of the polynomials and are left
with the first derivative information
only - this justifies the use of the ``first
order deformations'' term.

There is the following interpretation in
[continued-fractions] p. 39: the space
of first order deformation classes of $X$ is
$D(\mathbb{C}[t]/(t^2)$. This is
said to ````represent'' the tangent space
$\mathbb{T}^1_X$ of the hypothetical
deformation space of $X$''. (I put double
quotations because the word
``represent'' is quoted in the original
text.) Further, if $X$ is nonsingular
and compact, then
$\mathbb{T}^1_X=H^{1}(X,T_X)$.

Which basically I interpret as: the dimension
of the deformation space of a
smooth compact variety is $H^{1}(X,T_X)$.

\medskip\noindent
After a few months of the previous text in
this section
I think I have understood the construction of
Kodaira-Spencer map for embedded
deformations,
which is \cite[Proposition
3.2.1(ii)]{Sernesi-deformations}.

But I will actually briefly outline a
particular case
of this, which is easier: \cite[Examples
3.2.4(ii)]{Sernesi-deformations}.
This the case when $X \subset Y$ is a Cartier
divisor.
Then $X$ is cut out locally by a set of
functions $f_i$
defined on an affine cover $U_i$ of $Y$.
Then we have a line bundle defined by the
transition functions $f_{ij}=f_i/f_j$.
This is the line bundle $\mathcal{O}_Y(X)$,
which
by the adjunction formula is also the normal
bundle $N_{X/Y}$
of $X$ in $Y$.

Now consider a first order deformation
$$
\xymatrix{
X\ar[r]\ar[d]
&\mathcal{X}\subset Y\times
B_\varepsilon\ar[d]^\pi\\
0\ar[r]
&B_\varepsilon
}
$$

The great thing about first order
deformations
is that the structure sheaf of $\mathcal{X}$
on an affine open $U_i$ of $Y$
looks like the tensor product
$\mathcal{O}_{U_i}\otimes
\mathbb{C}[\varepsilon]/(\varepsilon^2)$
(\href{https://stacks.math.columbia.edu/tag/%
01I4}{01I4}).
By definition of tensor product this just
says
that sections look like $f+\varepsilon g$
for $f,g \in \Gamma(\mathcal{O}_{U_i})$.

If $\mathcal{X}$ is a Cartier divisor on $Y
\times B_\varepsilon$
(it is; but I'm missing a formal proof;
the intuition is that an ``infinitesimal
thickening'' $\mathcal{X}$
of $X$ has the same underlying topological
space
but different scheme structure;
\href{https://stacks.math.columbia.edu/tag/%
01JL}{01JL}
\href{https://stacks.math.columbia.edu/tag/%
01I4}{01I4}),
it is given as the vanishing locus
of some local sections $f_i+\varepsilon g_i$.
(Perhaps the $f_i$ coincide with the
distinguished
section of $X$.)
The line bundle associated to $\mathcal{X}$
is given by the cocycle
$$
F_{ij}=f_{ij}+\varepsilon g_{ij}
$$
for $g_{ij}=\frac{g_i-f_{ij}g_j}{f_j}$.
Indeed, note that
$\frac{1}{f_j+\varepsilon
g_j}=\frac{1}{f_j}-\varepsilon
\frac{g_j}{f^2_j}$
and compute $F_{ij}=(f_i+\varepsilon
g_j)/(f_j+\varepsilon g_j)$.

Then
$$
f_i+\varepsilon g_i=(f_{ij}+\varepsilon
g_{ij})(f_j+\varepsilon g_j)
$$
gives
$$
g_i=f_{ij}g_j+g_{ij}f_j
$$
and since $f_j$ vanishes along $X$,
we are left with $g_i=f_{ij}g_j$
which is how a section of the sheaf $N_{X/Y}$
is constructed by definition of sheaf.


\begin{example}
\label{example-modul-space-nonsi-riema}
\begin{reference}
\cite[Example 3.1]{continued-fractions}
\end{reference}
Fix $g \geq 2$. The dimension of the
deformation space of a nonsingular
projective curve $X$ is $3g-3$. This is ``the
dimension of the moduli space of
curves of genus $g$''.

We can compute this number by Riemann-Roch
formula on the bundle $-K_X$.
Indeed, since $X$ is a curve,
$\Omega_X^1=K_X$ and thus $-K_X=T_X$. We get
$$
\begin{aligned}
h^0(-K)-h^0(K-(-K))&=\text{deg}(-K)-g+1\\
&=-2g+2-g+1\qquad \substack{\text{degree
additive and}\\ \text{deg}K=2g-2}\\
&=-3g+3
\end{aligned}
$$
Now by Serre duality
$h^0(K-(-K))=h^1(-K)=h^1(T_X)$, and it turns
out that a
Riemann surface of genus $g \geq 2$ has no
holomorphic vector fields, so that
$h^0(-K)=h^0(T_X)=0$.
\end{example}

\begin{exercise}
\label{exercise-deformations-of-curves-in-K3}
Let $C$ be a smooth genus $g$ curve which can
be embedded in a K3 surface  $M$,
and $X$ the family of all deformations of $C$
in $M$.
\begin{enumerate}
\item Prove that $\dim X \leq  g$.
\item Let $\mathcal{X}_g$ be the space of all
curves of genus $g$ (smooth?)
which can be possibly embedded to a K3
surface. Prove that each irreducible
component $Z$ of $\mathcal{X}_g$ satisfies
$\dim_\mathbb{C} \leq  g+19$. Deduce
that there exists a compact complex curve
which cannot be embedded in a K3
surface.
\end{enumerate}
\end{exercise}

\begin{exercise}
\label{exercise-}
Let $C$ be a smooth curve embedded in a K3
surface $X$. Show that the dimension
of the deformation space of $C$ is $\leq g$.
\end{exercise}

\begin{proof}
% \begin{enumerate}
% \item
The deformation space of a variety is the
space of isomorphism classes of
deformations as explained above. It turns out
that there is a way to associate
1-cocyles of the tangent sheaf to
deformations, so that in fact the deformation
space $\text{Def}_1$ is isomorphic to
$H^{1}(X,\mathcal{T}_X)$ for any variety
$X$.

For our curve $C$ we thus know that the
dimension of the space of deformations
(deformations not necessarily contained in
$X$) is $h^1(\mathcal{T}_C)$.  The
family of deformations of $C$ that are
contained in $X$ is the Hilbert scheme of
curves with fixed Hilbert polynomial $P(t)$
after quotienting by $\mathbb{C}s$,
where $s$ is the section whose vanishing
locus is  $C$.
This says that the number we are looking for
is $h^{0}(X,\mathcal{O}(C))-1$.

Before showing that in fact
$h^0(X,\mathcal{O}(C))=1+g$,
let's explain why this Hilbert scheme
construction actually works.
Intuitively, any curve lying in the same
component
as $C$ would be deformation-equivalent to $C$
--- at least we know that it's Hilbert
polynomial is the same.
But it's not clear what's the deformation
involving these two
(this question should be addressed!).
Further, this doesn't explain Vladimiro's
suggesion of ``taking quotient by
$\mathbb{C}s$''.
To me it looks like I would just compute the
dimension
of the tangent space of $\text{Hilb}^P(M)$ at
$C$.

Now I will show that
$h^0(X,\mathcal{O}(C))=1+g$ (see
\cite[Lemma 1.2.1, Remark 1.2.2]{huk}).

Consider the ideal sheaf exact sequence
twisted by
$\mathcal{O}(C)$:
$$
\xymatrix{
0\ar[r]&\mathcal{O}_X\ar[r]&\mathcal{O}(C)
\ar[r]
&\mathcal{O}_X(C) \otimes
\mathcal{O}_C=\mathcal{O}(C)|_{C}\ar[r]&0
}
$$
By $X$ being a K3 we know that
$H^{1}(\mathcal{O}_X)=0$, so that we have the
short exact sequence in cohomology
$$
\xymatrix{
0\ar[r]&H^{0}(\mathcal{O}_X)\ar[r]&H^{0}
(\mathcal{O}(C))\ar[r]
&H^{0}(\mathcal{O}(C)|_{C})\ar[r]&0
}
$$
so that
$h^0(\mathcal{O}(C))=1+h^1(\mathcal{O}(C)
|_{C})$.
A version of adjunction formula says
$\omega_C \cong (K_X
\otimes\mathcal{O}(C)|_{C}$, and using that
$K_X=\mathcal{O}_X$ we obtain
$h^0(\mathcal{O}(C)|_{C})=h^0(\omega_C)=g$.

To

\medskip\noindent
For the record I put other thoughts I went
through in solving this exercise.

Recall from \cite[p. 146]{gri} that the
normal bundle
$\mathcal{N}$ of a hypersurface of a smooth
variety $X$ satisfies
$\mathcal{N}^\vee\cong\mathcal{O}_X(-C)|_{C}
$.
Taking duals we get that
$\mathcal{N}\cong\mathcal{O}_X(C)|_{C}$.

By adjunction formula
$2g-2=(\mathcal{O}(C),\mathcal{O}(C))$.
Applying Riemann-Roch
to $\mathcal{O}(C)$ (which is by definition
the dual of the ideal sheaf of $C$,
which is a line bundle on $X$), we obtain
that
$\chi(\mathcal{O}(C))=2+\frac{1}{2}
(\mathcal{O}(C),\mathcal{O}(C))$.
Then
$\chi(\mathcal{O}(C))=g+1$.

Now recall that
$\chi(\mathcal{O}(C))=h^0(\mathcal{O}(C)
-h^1(\mathcal{O}(C))+h^2(\mathcal{O}(C))$. By
Serre duality and $X$ being a K3
surface we see that $h^2(\mathcal{O}(C))\cong
h^{0}(\mathcal{O}(-C))$, which is
the ideal sheaf of $C$. Any section of such a
sheaf would vanish along $C$, and
since $X$ is compact we conclude there cannot
be any such section.

Now we show that also $h^1(\mathcal{O}(C))=0$
to conclude that
$h^0(\mathcal{O}(C))=g+1$.

We thus conclude that $h^0(\mathcal{O}(C))$




Since $C$ is smooth we can use the normal
exact sequence
$$
\xymatrix{
0\ar[r]&\mathcal{T}_C\ar[r]&\mathcal{T}
_X|_{C}\ar[r]&\mathcal{N}\ar[r]&0
}
$$
Taking Euler characteristic we see that
$\chi(\mathcal{T}_C)+\chi(\mathcal{N})=
\chi(\mathcal{T}_X|_{C})$.


This means that we should be done once we
compute $\chi(\mathcal{T}_X|_{C})$.
For this we can use Riemman-Roch formula for
coherent sheaves on a curve, which
tells us that
$$
\text{deg}(\mathcal{T}_X|_{C})=
\chi(\mathcal{T}_X|_{C})
-\text{rk}(\mathcal{T}_X|_{C})
\cdot\chi(\mathcal{O}_C)
$$
But of course we know that since $C$ is a
curve, by Serre duality we get that
$\chi(\mathcal{O}_C)=1-g$. (Indeed:
$h^1(\mathcal{O}_C)=h^0(\Omega^1_C):=p_a(C)
$.)
And now the question is what is
the degree. Apparently this is just the first
Chern class $c_1$ of the
restricted tangent bundle. So what is it? And
what is $h^0(\mathcal{T}_C)$ if
the genus is 0 or 1, and that's it.

I
know that $h^0(\mathcal{T}_X)=0$ and
$h^1(\mathcal{T}_X)=20$ by $X$ being a K3
surface, but I'm not sure what happens when
we restrict to $C$.

If $g \geq 2$ we know that
$h^0(\mathcal{T}_C)=0$, so that
$\boxed{\chi(\mathcal{T}_C)=h^1(\mathcal{T}
_C)}$.
\bigskip
To compute the
latter Euler characteristic, which is given
by definition by
$\chi(\mathcal{T}_X|_{C})=h^0(\mathcal{T}
_X|_{C})-h^1(\mathcal{T}_X|_{C})$,
we
first note that $h^0(\mathcal{T}_X|_{C})=0$,
because this is the dimension of
the global holomorphic vector fields on $X$
restricted to $C$, which is constant along
$C$ since $X$ is
smooth. And then this number is actually zero
by Hodge numbers of a K3 surface.

and taking cohomology long exact sequence we
obtain
$$
\xymatrix{
\cdots\ar[r]&H^{0}(T_X)\ar[r]&H^{0}
(\mathcal{N})\ar[r]&H^{1}(T_C)\ar[r]&\\
H^{1}(T_X)\ar[r]&H^{1}(\mathcal{N})\ar[r]&
\cdots
}
$$
\medskip\noindent
Or is it $h^1(\mathcal{N})$? See
\cite{HarrMorr}. If this was the case, then
we
can use adjunction formula as above to get
that
$2g-2=(\mathcal{N},\mathcal{N})=\text{deg}_C
\mathcal{N}$. Then we may find
$h^0(\mathcal{N})$ via Riemann-Roch:
$$
h^0(\mathcal{N})-h^0(K_C-\mathcal{N})=
\text{deg}\mathcal{N}-g+1
$$
Note that
$h^0(N_C-\mathcal{N})=h^1(-K_C+N+K_C)=
h^1(\mathcal{N})$
via Serre
duality, and by Riemann-Roch on a surface as
above we see that
$$
g+1=\chi(\mathcal{N})=h^0(\mathcal{N})-
h^1(\mathcal{N})
\implies
h^1(\mathcal{N})=-g-1+h^0(\mathcal{N})
$$
so that
$$
\begin{aligned}
h^0(\mathcal{N})-(-g-1+h^0(\mathcal{N}))&=
\text{deg}\mathcal{N}-g+1\\
\implies \text{oops! I lost
$h^0(\mathcal{N})$ in this operation…}
\end{aligned}
$$
Maybe if I just use normal exact sequence and
realise that
$$
h^0(\mathcal{N})=h^0(\mathcal{T}_X|_{C})-
h^0(\mathcal{T}_C)
$$
I know that $h^0(\mathcal{T}_C)=0$ for $g>1$,
so the question is how to compute
the restricted holomorphic vector fields.
% \end{enumerate}
\end{proof}

\medskip\noindent
Look for the ``Perfect obstruction theory''.
This is used by Bruzzo for deformations of
(super)stacks,
work with E. Paiva and ---.

\medskip\noindent
Until I have a better candidate of where to
put this…

\begin{definition}
\label{definition-extension}
\begin{reference}
\cite[p.7]{Sernesi-deformations}
\end{reference}
Let $A \to R$ be a ring homomorphism
(i.e. $R$ is an $A$-algebra).
An {\it $A$-extension of $R$ by $I$} is a
short exact sequence
$$
\xymatrix{
0\ar[r]
&I\ar[r]
&E\ar[r]^\varphi
&R\ar[r]
&0.
}
$$
where $E$ is an $R'$-algebra
and $\varphi$ is a morphism of $A$-algebras
whose kernel is $I$.
We also require that $I^2=(0)$;
this condition implies that $I$ has the
sturcture of $R$-module.
\end{definition}

\begin{example}
\label{example-dual-numbers-extension}
$$
\xymatrix{
0\ar[r]
&(\varepsilon)\ar[r]
&A[\varepsilon]\ar[r]
&A\ar[r]
&0
}
$$
where $A[\varepsilon]$ is the algebra of dual
numbers over $A$.
\end{example}

\section{Cotangent modules}
\label{section-cotangent-modules}

\noindent
My aim in this section is to cover the
minimal
theory needed to apply results in \cite{jan2}
to my deformation problem. The main theorem
I expect to use is the following:

\begin{theorem}
\label{theorem-deformations-threefolds}
\begin{reference}
\cite[Theorem 5.7]{jan2}
\end{reference}
If $\mathcal{K}$ is a 3-dimensional manifold
then
$$
\dim T^1_{\mathbb{P}(K)}
=11d_3+5e_3+3e_4+e_{\geq 5}+c_{\geq
6}+5f_1^{(3)}+2f_1^{(4)}+h^2(\mathcal{K}).
$$
\end{theorem}

\noindent
I need to do some interpretation first.
From \cite[p. 1]{jan3}:

\begin{quotation}
  ``For each $\mathbb{C}$-algebra $A$,
there is a cohomology theory providing
modules $T^i_A$.
However, only $T^1_A$ and $T^2_A$ are
relevant for the deformation
theory of $\Spec A$. The module $T^1_A$
collects the infinitesimal
deformations of $A$ and $T^2_A$ contains the
obstructions
for lifting these deformations to decent
parameter spaces.
Eventually, $\text{Der}_{\mathbb{C}}(A,A)$
will be called $T_A^0$."
\end{quotation}

\cite[p. 3]{jan3}:
\begin{quotation}
``In general, for a finitely generated
$\mathbb{C}$-algebra $A$,
the modules $T_A^i$ allow the following ad
hoc definitions:
Let $P=\mathbb{C}[x_0,…,x_n]$ mapping onto
$A$ so that $A \simeq P/I$
for an ideal $I$. The $T_A^1$ is the cokernel
of the natural map
$\text{Der}_{\mathbb{C}}(P,P) \to
\Hom_P(I,A)$.
Moreovr, […] and we obtain $T^2_A$ as the
cokernel
of the induced map
$\Hom_P(P^m,A)\to\Hom_A(R/R_0,A)$.''
\end{quotation}

\begin{remark}[dani]
\label{remark-dani1}
Let's think more about that definition of
$T^1_A$
as the cokernel of the ``natural map''
$\text{Der}_\mathbb{C}(P,P) \to \Hom_P(I,A)$.

Pick a derivation $\chi \in
\text{Der}_{\mathbb{C}}(P,P)$.
We want a map $I \to A=P/I$.
Pick an element $a \in I$ and map it to
$\chi(a)$.
And that's it. It might be zero or not.
If the $\chi$ we picked maps any $a \in I$ to
some element in $I$,
then this $\chi$ is in the kernel of the
natural map.
So $T^1_A$ is the image of the natural map
modulo the kernel.
\end{remark}

\noindent
Unfortunately this doesn't say much about the
following crucial fact:
the $T^i_A$ modules are
$\mathbb{Z}^{n+1}$-graded.
But here goes an explanation.
We use the notation that $e_p$ is a ``face
not in the complex'',
which of course corresponds to a generating
monomial $x^p$
of the SR ideal $I$.
It has a degree, of course (given by the
``support'' of $p$).
Now here's a rather sloppy step but:
since the cotangent modules are built upon
$I$,
who is given by $e_p$,
we just notice that the degree of $e_p$
along with a (very natural, too) notion of
degree among relations,
induce a grading in all the cotangent modules
$T_A^i$.
Thus, the grading is morally just the degree
of the monomials,
which is in fact the support of $p$, a vector
$\mathbf{c} \in \mathbb{Z}^{n+1}$.

Next in line is the decomposition
$\mathbf{c}=\mathbf{a}-\mathbf{b}$
into positive minus negative part. But how
exactly…?








\section{Fano varieties}
\label{section-Fano-varieties}

\begin{definition}
\label{definition-Fano-variety}
A {\it Fano variety} is a projective variety
with $-K_X$ ample.
\end{definition}

\noindent
The idea of the Fano index is that the
larger the index, the simpler is the variety.
In a talk by Jõao, it is defined as the
larger integer
dividing $-K_X$ in $\text{Pic}(X)$.
\begin{definition}
\label{definition-Fano-index}
$$
r(X):=\text{min}\{r:\frac{c_1(X)}{r}\in
H^{2}(X,\mathbb{Z})\}
$$
\end{definition}

\begin{exercise}
\label{exercise-fano-vanis-highe-cohom}
By Kodaira vanishing theorem
\ref{theorem-Kodaira-vanishing},
you can show that the cohomology $H^{i}(X,L)$
for
a Fano variety $X$ vanishes. You just have to
put $L=\mathcal{O}(k)$ with $k\geq
-r$, where $r$ is the Fano index.
\end{exercise}

\begin{exercise}
\label{exercise-Pic-H2-Fano}
Show that  $\text{Pic}(X)\cong
H^{2}(X,\mathbb{Z})$ holds for Fano
varieties.
\end{exercise}

\begin{remark}
\label{remark-deriv-categ-fano-3-folds}
\begin{reference}
Marcos Jardim, CIMPA 2025 Florianópolis,
Lecture 2.
\end{reference}
If $H^3(X,\mathbb{Z})=0$ of a Fano 3-fold,
then its derived category is
generated by 4 elements.
\end{remark}

\medskip\noindent
Araujo-Druel em 2017 classify
Fano foliations relating the Fano index (of
the foliation)
with its rank.
João found conditions
(integrability of the leaves)
to describe del Pezzo foliations
(a kind of Fano foliation)
with Picard rank 1.
Araujo (2019) classifies the leaves of a
algebraically integrable
del Pezzo foliation (with no restriction on
the (log
canonica) singularities!).


\section{Quivers}
\label{section-quivers}

\begin{definition}
\label{definition-quiver}
A {\it quiver} is a set of vertices $Q_0$, a
set of arrows $Q_1$ equipped with
the maps of source $s$ and target $t$ that to
each arrow they assign the point
that is source or target of the arrow.
\end{definition}

\begin{definition}
\label{definition-representation-of-quiver}
A {\it representation} of a quiver is a set
of finite dimensional vector spaces
equipped with maps between them realising a
given quiver (incomplete…).
\end{definition}

There is a notion of projective
representation, which I missed to write. But
it
is analogous to the injective representation:

\begin{definition}
\label{definition-injec-repre-quive}
Given a quiver $Q$, the {\it injective
representation} of $Q_0$ is given by, for
$i \in Q_0$,
$$
I(i)_j=\begin{cases}
k\qquad &i=j \\
k^{d'}\qquad &j \neq i
\end{cases}
$$
where $d'$ is the number of paths from $j$ to
$i$.
\end{definition}

\section{Stacks}
\label{section-stacks}

My first definition of stack can be extracted
from

\begin{definition}
\label{definition-superstack}
A {\it superstack} is a stack over
the étale site $\text{SSch}$ of superschemes,
i.e. it is a category fibered in groupoids
over the category of superschemes,
the latter equipped with the
étale topology,
satisfying the descent condition.
\end{definition}

Here are some other definitions:

\begin{definition}
\label{definition-algebraic-stacks}
Let  $\mathfrak{X}$ be a stack over
$\text{Sch}_{\text{ét}}$.
An {\it algebraic space} is
such that there exists morphism
$\mathcal{U} \to \mathfrak{X}$
where $\mathcal{U}$ is a scheme, that is
schematic, étale and injective (check this
one).

$\mathfrak{X} \to y$ is {\it representable}
if
there exists a scheme $\mathcal{U}$ and a map
$\mathcal{U} \to y$ such that the
fibered product
$$
\xymatrix{
\mathcal{U} \times_y
\mathfrak{X}\ar[r]\ar[d]\ar@{}[dr]|-
{\lrcorner}&\mathfrak{X}\ar[d]\\
\mathcal{U}\ar[r]&y
}
$$
is an algebraic space.

Finally, a  stack is {\it algebraic} (resp.
{\it Deligne-Mumford})
is there exists a
representable surjective morphism
$\mathcal{U} \to \mathfrak{X}$
that is smooth (resp. étale).

A {\it stable map} over a projective
variety $X$ is an element of the first
Chow group $\beta \in A_1$, where
 $(C,g)$ is an algebraic curve and
$f:C \to X$ with $[f(C)]=\beta$.
\end{definition}

\noindent
The curves that are points under this map
(contractible) are {\bf stable}.

\section{Moduli}
\label{section-moduli}

\noindent
Pseudo-admissible covers may be the way
Mumford orifinally developed stability
notions.

\noindent
This is a two-session minicourse by Ugo
Bruzzo
in LEGALzinho 2025.

Let $X$ be a smooth complex projective
variety,
$X \overset{j}{\hookrightarrow}\mathbb{P}^n$.
$\mathcal{O}_{\mathbb{P}^n}(1)$ is the class
of hyperplanes of $\mathbb{P}^n$.
$\mathcal{L}=j^*\mathcal{O}_{\mathbb{P}^n}(1)
$.
$c_1(\mathcal{L})\in H^2(X,\mathbb{Z})$ is
a {\it polarization}.

Let $\mathcal{E}$ be a coherent
$\mathcal{O}_X$-module.
For example, $\mathcal{E}$ may be a locally
free $\mathcal{O}_X$-module.
Recall that $H^k(X,\mathcal{E})$ are finite
dimensional vector spaces
and they vanish for $k>\dim X=n$.
Recall that the {\it Euler characterstic} of
$\mathcal{E}$
is given by
$\chi(\mathcal{E})=\sum_{k=0}^n(-1)
^k\dim_{\mathbb{C}}H^k(X,\mathcal{E})$.

The {\it reduced Hilbert polynomial} of
$\mathcal{E}$ is
$$
p_{\mathcal{E}}(k)=\frac{1}{\text{rk}
\mathcal{E}}
\chi(\mathcal{E}\otimes_{\mathcal{O}_X}
\mathcal{L}^k).
$$

\begin{definition}
\label{definition-stable-module}
Let $\mathcal{E}$ be a coherent
$\mathcal{O}_X$-module
with no torsion and positive rank.
$\mathcal{E}$ is {\it stable} if for all $f
\subset \mathcal{E}$
($0 \text{rk} f<\text{rk}\mathcal{E}$),
$$
p_f(k)<p_{\mathcal{E}}(k)\qquad \text{for }
k\gg 0
$$
and {\it semi-stable} if we put $\leq$
instead of $<$.

\begin{lemma}
\label{lemma-stable}
If $\mathcal{E}$ is stable then $\mathcal{E}$
is irreducible
($\mathcal{E} \neq  \mathcal{E}_1 \oplus
\mathcal{E}_2$)
and $\text{Aut}(\mathcal{E})\simeq
\mathbb{C}^*$.
\end{lemma}

\end{definition}

The {\it functor of moduli} is
$$
\begin{aligned}
\mathcal{M}: \text{Sch}_\mathbb{C}
&\longrightarrow \mathsf{Set} \\
S
&\longmapsto
\left\{
\substack{\text{stable family}  \\ \text{of
sheaves with} \\
\text{Hilbert polynomial} f\\
\text{parametrized by $S$}}
\right\}\Big/\sim.
\end{aligned}
$$

\medskip\noindent
Riemann-Roch.
Let $C$ be a smooth, projective, connected
curve
of genus $g$.
Let $\mathcal{E}$ be a coherent sheaf over
$H^2(C,\mathbb{Z})$
with generator of $H^2(X,\mathbb{Z})$
$w=[pt] \in H^2(C,\mathbb{Z})$.
$c_1(\mathcal{E})=\text{deg}\mathcal{E}\cdot
w$.

\begin{theorem}
\label{theorem-euler-characteristic-moduli}
$$
\begin{aligned}
\chi(\mathcal{E})
&=\int_C(1+(1-g)w)(\text{rk}\mathcal{E}+
\text{deg}\mathcal{E}w)\\
&=\int_C
\text{deg}\mathcal{E}+w+(1-g)\text{rk}
\mathcal{E}\cdot
w\\
&=\text{deg}(\mathcal{E})+1-g
\text{rk}\mathcal{E}.
\end{aligned}
$$
\end{theorem}

\begin{example}
\label{example-hilbert-polynomial-sheaf}
\end{example}

\begin{theorem}
\label{theorem-moduli-suave}
Se $(r,d)$, $r>0$ are coprime,
the functor of moduli of stable bundles
over $C$ of rank $r$ and degree $d$
is representable, and the moduli space
$\mathcal{M}_C(\mathcal{E},d)$ is smooth.
\end{theorem}

\noindent
In the non coprime case this does not
necessarily hold.

\begin{definition}
\label{definition-}
$\mathcal{M}$ is a {\it grosseiro} moduli
space
if the funtor is not representable.
\end{definition}

\noindent
This means that although there is not
a universal family,
there is still a correspondence
of the closed points of $\mathcal{M}$ and the
equivalence classes of semistable sheaves
of rank $r$ and degree $d$.

\medskip\noindent
Now let's fix $\dim X=2$.
Let $\mathcal{E}$ be a stable
$\mathcal{O}_X$-module.
Let $r=\text{rk}\mathcal{E}$,
and the {\it discriminant}
$$
\Delta(\mathcal{E})=c_2(\mathcal{E})-
\frac{\mathcal{E}-1}{2r}c_1(\mathcal{E})^2
\in H^4(X,\mathbb{Q})\cong\mathbb{Q}.
$$
Here enter $c_1(\mathcal{E}) \in
H^2(X,\mathbb{Z})$
and $c_2(\mathcal{E}) \in H^4(X,\mathbb{Z})$.

\begin{theorem}[Bogomolov]
\label{theorem-bogomolov}
$\Delta(\mathcal{E}) \geq 0$.
\end{theorem}

\noindent
It could be interesting to study the proof of
this theorem.

\medskip\noindent
Miyaoka (1987). Consider a curve $C$ and a
bundle $E$.
We can for any fibre, which is a vector space
$V$,
consider its projectivization
$V\setminus\{0\}/\mathbb{C}^*=\mathbb{P}(V)$.
Then we can bundlize this construction ad
obtain
$\mathbb{P}E \to C$.
Consider $\mathcal{O}_{\mathbb{P}E}(1)$.
We define
$$
\lambda=rc_1(\mathcal{O}_{\mathbb{P}E}(1))
-\pi^*c_1(E) \in H^2(\mathbb{P}E,\mathbb{Z}).
$$

\begin{theorem}
\label{theorem-semistable-nef}
$E$ is semistable if $\lambda$ is
semicorrente efectiva.
\end{theorem}

\noindent
\begin{definition}
\label{definition-nef}
Consider again a projective smooth connected
curve $\Gamma$ and let $f:\Gamma \to
\mathbb{P}E$.
We say $\lambda$ is {\it nef} if
$$
\int_{\Gamma} f^*\lambda\geq 0.
$$
\end{definition}


\noindent
The following theorem was proved by Nakayama
(who is a different author than that of
Nakayama's lemma)
in 1999 but was unkown to Bruzo and Daniel
HR.
\begin{theorem}[Nakayama '99, Bruzzo-HR '04]
\label{theorem-bruzzo-daniel-nakayama}
The following are equivalent:
\begin{enumerate}
\item $\lambda$ is nef.
\item For all $f:\Gamma \to X$,
$f^*E$ is semi-stable.
\item $E$ is semistable with respect to
the polarization, and  $\Delta(E)=0$.
\end{enumerate}
\end{theorem}

The polarization says how the surface is
embedded in projective space,
it is $H=c_1(\mathcal{L})$. The stability
notion is
$$
\mu(E)=\frac{c_1(E)\cdot
H^{n-1}}{\text{rk}E}.
$$
\medskip\noindent
They have proved that for Higgs bundles 1
$\iff$ 2, $3 \implies 1,2$.
And also for $1,2 \implies 3$ but only for
$r=2$ using the Higgs Grassmannian,
but the problem is very complicated for
general rank.

\medskip\noindent
Second lecture. I think this might be a more
general
definition of the Hilbert polynomial:
$$
\sum_{k=0}^n(-1)^k \dim
H^k(X,\mathcal{E})=\chi(\mathcal{E})
=\int_X\text{ch}(\mathcal{E})\text{td}X
$$
And we also have:
$$
\chi(\mathcal{E}\otimes \mathcal{L}^*)
=\int_X (1+(1-g)w)(r+kr
\underbrace{c_1(\mathcal{L})}_{\text{deg}
\mathcal{L}\cdot
w}
+\text{deg}\mathcal{E}\cdot w.
$$
And then
$$
p_{\mathcal{E}}=k \text{deg} \mathcal{L}
+1-g+\mu(\mathcal{E}).
$$



\medskip\noindent
Now we discuss Higgs bundles.
Let $X$ be ac onnected complex projective
variety.

\begin{definition}
\label{definition-higgs-bundle}
A {\it Higgs bundle} on $X$
is a pair $\mathfrak{C}=(E,\phi)$ with
\begin{itemize}
\item $E$ a vector bundle on $X$
\item $\phi:$ a morphism (connection-like…)
\end{itemize}
\end{definition}

For the following see Nitsuze, FGA explained.
There following scheme is also representable
(I think this is the ``universal quotient'')

We have also consider the Grassmanization of
a vector bundle.
We obtained a short exact sequence
and tried to compute the tensor product with
the holomorphic differentials.
The idea was to pullback the grassman bundle
given a map $f:Y \to X$.
The result is not satisfactory, the pulled
back
thing may be strange. This is called the
{\it Higgs Grassmannian}.


Here's some notion also studied by Campana,
Demailly:
numerically flat bnef bundle $E$.
Then we have that
\begin{lemma}
\label{lemma-proof-hugo-bruzzo}
If $c_1=0$ then numerically flat bundle iff
semistable.
\end{lemma}

\medskip\noindent
This is the beginning of a talk by Alan Muniz
at Bandoleros 2025 titled
``Rank 2-bundles on $\mathbb{P}^3$ with odd
determinant''.
(Work in progress with A. Fontes.)

Problem: Classify rank-2 bundles on
$\mathbb{P}^3$.

We consider bundles $E$ with chern classes
$c_1=-1$, $c_2=2n$ and $c_3=0$.
Let $\mathcal{B}(-1,2n)$ be the moduli space
of $\mu$-stable rank 2-bundles.

\begin{itemize}
\item $\mathcal{B}(-1,2)$ is irreducible of
dimension 11
[Hartshorne-Sols '82, Manolache '81].

\item $\mathcal{B}(-1,4)$ has two components
[Barnica-Manolache '85].

\item $c_2\geq 6$, Tikhomirov (2014)
Fontes-Jardim (2023, 2024, 2025).
$\mathcal{B}(-1,6)$ has $\geq 4$ components.
\end{itemize}

\medskip\noindent
Now let's discuss Serre correspondence.

How can we describe these spaces?
Let $E$ be a vector bundle of rank 2
on $\mathbb{P}^3$ and $\sigma \in H^0(E(r))$.
Then $C=(\sigma_0)\subset \mathbb{P}^3$ is
a curve * equipped
with an ``extension class''
$\zeta \in H^0(\omega_C(4-2r-c_1))$.
In fact, this correspondence
is 1-1, and we call it Serre correspondence.
We have that $\text{deg}C=c_2(E(r))$ and
$p_a(C)=1=\frac{c_2(E(r))(c_1(E(r))-4)}{2}$
To pass from the curve to the bundle
and viceversa we use
$$
\xymatrix{
0\ar[r]
&\mathcal{O}(-r)\ar[r]
&E\ar[r]
&\mathcal{I}_C(c_1+r)\ar[r]
&0
}
$$


\medskip\noindent
Upshot about Deligne-Mumford stacks:
the sit below a stack via a map that is
étale.


















\medskip\noindent
Now let's discuss framed moduli.

\noindent
Framed instantons on the blow-up of
$\mathbb{P}^3$ at a point.

Let $X$ be a projective variety and fix a
polarization $H$.
Let $\mathcal{E}$ be a sheaf on $X$.
A {\it framed pair} is a
$(\mathcal{F},\alpha)$ where
$\mathcal{F}$ is a coherent sheaf on $X$
and $\alpha$ is a {\it framing datum}, i.e.
a map $\alpha:\mathcal{F} \to \mathcal{E}$.

Set
$$
\mathcal{E}(\alpha)
=\begin{cases}
	1\qquad &\alpha\neq 0 \\
	0\qquad &\text{otherwise} .
\end{cases}
$$
Define the Hilbert polynomial by
$$
P(k)=\chi(\mathcal{F} \otimes
\mathcal{O}(kH)).
$$
For fied $\delta \in \mathbb{Q}[t]$ rational
positive,
define
$$
P_{(\mathcal{F},\alpha)}=P_{\mathcal{F}}(k)-
\varepsilon(\alpha)\delta.
$$

We can also define induced framings from
subsheaves $\mathcal{F}' \subset
\mathcal{F}$, by asking that the inclusion
commutes
with the framing datum $\alpha$.
Then we put
$$
\varepsilon(\alpha')
=\begin{cases}
	1\qquad &\mathcal{F}' \subseteq Ker \alpha
	\\
	0\qquad &\text{otherwise} .
\end{cases}
$$
We say that $\mathcal{F}$ is {\it
$\delta$-(semi)stable}
if for any subpair $(\mathcal{F}',\alpha')$
with
$\text{rk}\mathcal{F}'=r'<r=\text{rk}
\mathcal{F}$
one has
$$
\frac{1}{r'}P_{(\mathcal{F}',\alpha')}(k)
\leq  \frac{1}{r}P_{(\mathcal{F},\alpha)}(k).
$$
If $\text{deg}(\delta)\leq \dim X$,
there are moduli functors
$\mathcal{M}^{st}(X,\mathcal{E},\delta)
\subseteq
\mathcal{M}^{\text{ss}}(X,\mathcal{E},\delta)
$.

\begin{remark}
\label{remark-degree-greater-than-dimension}
If $\text{deg} \delta>\dim X$
then $\alpha$ is injective or zero, and
$\mathcal{F}$ is semistable.
This gives quot schemes.
\end{remark}

\noindent
Huybrechts-L. In $\dim X\leq 2$
there exists a quasi-projective fine moduli
space
for $\delta \in \mathbb{Q}[t]$ and $p \in
\mathbb{Q}[k]$.

\medskip\noindent
Then there is a notion of good framing by
Bruzz-M,
depending on $D$, an effective divisor on
$X$.
They proved that having good framing is
enough
to obtain a fine quasi-projective moduli
space.

\begin{definition}
\label{definition-instanton}
A coherent sheaf $\mathcal{E}$ on
$X=\tilde{\mathbb{P}^3}$
is an {\it instanton} if it is
$\mu_{\mathcal{L}}$-semistable,
$c_1(\mathcal{E})=0$ and
$H^0(\mathcal{E})=0=H^0(Ec(-3,2))$,
$H^1(\mathcal{E}(-2,1))=0=H^2(\mathcal{E}(-2,
1))$
and
$H^3(\mathcal{E}(0,-1))=H^2(\mathcal{E}(-1,1)
)$.
\end{definition}

\noindent
Some examples are locally free sheaves of
rank 2 (via Harrshorne-Serre
correspondence),
torsion-free rank 2 constructed via
elementary
transformations.

\medskip\noindent
General result: for rank 2 instantons on Fano
threefolds
that have a plane $E$,
$\mathcal{E}|_{E}$ is $\mu$-semistable
if and only if $h^1(\mathcal{E}(-2,1) \otimes
\mathcal{O}(-E))=0$.




\section{Stability}
\label{section-stability}

\noindent
Stability helps us select the
building blocks of an abelian
category, every object in the
category can be reconstructed
from those.
Wall and chamber structure: inside a
compact set $K \subset
\text{Stab}(\mathcal{T})$,
only finitely many walls appear.
If $C$ is a connected component of
the complement of the walls in $K$,
then the stable objects are constant
on $C$.
Think of an ant walking in
$\text{Stab}(\mathcal{T})$ with a
backpack of objects and a detector:
when it crosses a wall and some object
stops being stable, the red light fires.
Stable objects change only when you
cross a wall.

\begin{center}
\setlength{\unitlength}{1mm}
\begin{picture}(98,78)(0,0)
\scriptsize

% axes
\linethickness{0.35pt}
\put(6,7){\vector(1,0){88}}
\put(88,7){\vector(0,1){68}}
\put(92,10){$\beta$}
\put(91,68){$\alpha$}

% Gieseker/PT chamber boundary
\linethickness{0.45pt}
\qbezier(18,7)(18,13)(24,15)
\qbezier(24,15)(31,16)(38,19)
\qbezier(38,19)(49,23)(60,22)

\qbezier(60,22)(57,27)(57,35)
\qbezier(57,35)(50,39)(47,45)
\qbezier(47,45)(41,53)(52,56)
\qbezier(52,56)(58,57)(68,58)
\qbezier(68,58)(75,62)(77,74)

% diagonal wall, blue in the photo
\linethickness{0.9pt}
\qbezier(22,73)(46,40)(70,7)

% semicircular wall, red in the photo
\linethickness{0.75pt}
\qbezier(45,7)(45,23)(60,22)
\qbezier(60,22)(75,23)(75,7)

% marked point and labels
\linethickness{0.4pt}
\put(59.35,21.35){\circle*{1.3}}
\put(61.5,23.5){$(\beta^\nu,\alpha^\nu)$}

\put(20,36){Gieseker chamber}
\put(47,59){PT chamber}
\put(39,42){$\theta^-_\nu$}

\end{picture}
\end{center}

Let $\mathcal{A}$ be an abelian
category with Grothendieck group
$K_0(\mathcal{A})$.

\begin{definition}
\label{definition-stability-function}
An additive homomorphism
$Z: K_0(\mathcal{A}) \to \mathbb{C}$
is a {\it stability function} if for all
$0 \neq E \in \mathcal{A}$,
$$
\text{Im}Z(E) \geq 0\qquad
\text{and}\qquad
\text{Re}Z(E) <0 \text{ if }
\text{Im}Z(E)=0.
$$
The {\it slope} of $E \in \mathcal{A}$
with respect to $Z$ is
$$
M(E)=
\begin{cases}
-\frac{\text{Re}Z(E)}{
\text{Im}Z(E)}\qquad &
\text{if }\text{Im}Z(E) \neq 0 \\
\infty\qquad &
\text{if }\text{Im}Z(E)=0.
\end{cases}
$$
A nonzero object $E \in \mathcal{A}$
is called {\it stable} if for all
proper non trivial subobjects
$F \subset E$, we have
$$
M(F) < M(E)
$$
and {\it semistable} if we put
$\leq $ instead.
\end{definition}

\begin{definition}
\label{definition-stability-condition}
A pair $(\mathcal{A},Z)$ is called
a {\it stability condition} if
\begin{enumerate}
\item $Z$ is a stability function, and
\item every $0 \neq E \in \mathcal{A}$
has a Harder-Narasimhan filtration.
\end{enumerate}
\end{definition}

\begin{lemma}
\label{lemma-harder-narasimhan-existence}
Let $Z:K_0(\mathcal{A})\to \mathbb{C}$
be a stability function.
Assume that
\begin{enumerate}
\item $\mathcal{A}$ is Noetherian, and
\item the image of $\text{Im}Z$
is discrete in $\mathbb{R}$.
\end{enumerate}
Then every $0 \neq E \in \mathcal{A}$
has a Harder-Narasimhan filtration
with respect to $Z$.
\end{lemma}

\noindent
Let $\mathcal{D}$ be a triangulated
category. Fix a finite rank lattice
$\Lambda$, a surjective homomorphism
$v:K_0(\mathcal{D}) \to \Lambda$
and a norm $\|\cdot\|$ in $\Lambda$.

\begin{definition}
\label{definition-bridgeland-stability}
A {\it Bridgeland stability condition}
$\mathcal{D}$ is a pair $(\mathcal{A},Z)$
where $\mathcal{A}$ is the heart of a
bounded $t$-structure on $\mathcal{D}$,
such that
\begin{enumerate}
\item $(\mathcal{A},Z)$ is a stability
condition,
\item the following support property
holds:
$$
\text{inf}\left\{
\frac{|Z(E)|}{
\|v(E)\|}
:0 \neq  E \in \mathcal{A}
\text{semistable}\right\}
>0.
$$
\end{enumerate}
\end{definition}

\begin{example}
\label{example-stability-curve}
Let $X$ be a smooth projective
curve and $\mathcal{D}=D^b(X)$.
Take $\mathcal{A}=\Cohstack(X)$ as the
heart of the standard
$t$-structure.
Let
$$
Z(E)=-\deg(E)+i\,\mathrm{rk}(E).
$$
Then $\sigma$-stability coincides
with slope stability.
\end{example}

\noindent
Let $Stab^{Br}_\Lambda(\mathcal{D})$
be the set of Bridgeland stability
conditions and let
$Stab_\Lambda(\mathcal{D})$ be the set
of stability conditions on $\mathcal{D}$
with respect to $(\Lambda,v)$.

\begin{theorem}[Bridgeland]
\label{theorem-bridgeland}
The natural map
$$
Stab_\Lambda(\mathcal{D})
\to \Hom(\Lambda,\mathbb{C})
$$
is a local homeomorphism.
$Stab_\Lambda(\mathcal{D})$ is a complex
manifold and
its dimension is the rank of $\Lambda$.
In particular,
$$
Stab_\Lambda^{Br}(\mathcal{D})
\subset
Stab(\mathcal{D}).
$$
\end{theorem}

\noindent
Let $(\mathcal{A},Z)$ be a stability
condition and $M$ be the slope with
respect to $Z$.
For each $\beta \in \mathbb{R}$
we define the following pair of
subcategories on $\mathcal{A}$:
$$
\begin{aligned}
\mathcal{T}^\beta(A,Z)&=\mathcal{T}^\beta
=\{\mathcal{E} \in \mathcal{A}
: M(\mathcal{G})>\beta
\text{ whenever }
\mathcal{E} \to\!\!\!\!\!\to
\mathcal{G}\}\\
\mathcal{F}^\beta(\mathcal{A},Z)
&=\mathcal{F}^\beta
=\{\mathcal{E} \in \mathcal{A}
:M(\mathcal{F})\leq \beta
\text{ for every }
0 \neq  \mathcal{F} \subset \mathcal{E}\}.
\end{aligned}
$$

\begin{definition}
\label{definition-tilt}
The {\it tilt}
$\mathcal{A}^*
=\left\langle{}\mathcal{F}[1],\mathcal{T}\right\rangle{}
\subset\mathcal{D}$
of $\mathcal{A}$ with respect to
$(\mathcal{T}^\beta,\mathcal{F}^\beta)$
is the smallest extension closed full
subcategory of $\mathcal{D}$ containing
$\mathcal{T}^\beta$ and
$\mathcal{F}^\beta[1]$.
\end{definition}

\noindent
$\mathcal{A}^*$ is the heart of a bounded
$t$-structure on $\mathcal{D}$.

\begin{definition}
\label{definition-gieseker-stability}
Let $\mathcal{E}$ be a torsion free
sheaf.
We say that $\mathcal{E}$ is
{\it Gieseker semistable} if
for every subsheaf $F \subset \mathcal{E}$
with $0 < \mathrm{rk}(F) <
\mathrm{rk}(\mathcal{E})$,
$$
\frac{P_F(m)}{\mathrm{rk}(F)}
\leq
\frac{P_{\mathcal{E}}(m)}{
\mathrm{rk}(\mathcal{E})}
\qquad \text{for } m \gg 0.
$$
It is {\it Gieseker stable}
if we replace $\leq$ by $<$.
Here $P_{\mathcal{E}}(m)$ is the
Hilbert polynomial of $\mathcal{E}$.
\end{definition}

\begin{definition}
\label{definition-2-gieseker-stability}
Let $\mathcal{E}$ be a torsion free
sheaf.
We say that $\mathcal{E}$ is
{\it 2-Gieseker semistable} if
for every subsheaf $F \subset \mathcal{E}$
with $0 < \mathrm{rk}(F) <
\mathrm{rk}(\mathcal{E})$,
$$
\frac{P_F(m)-\chi(F)}{\mathrm{rk}(F)}
\leq
\frac{P_{\mathcal{E}}(m)-\chi(\mathcal{E})}
{\mathrm{rk}(\mathcal{E})}
\qquad \text{for } m \gg 0.
$$
It is {\it 2-Gieseker stable}
if we replace $\leq$ by $<$.
\end{definition}

\noindent
The Gieseker chamber is the region
where this is the relevant stability.
The PT chamber is the neighboring region
where the PT condition is the relevant
one.

\begin{definition}[PT stability]
\label{definition-pt-stability}
\end{definition}

\begin{theorem}
\label{theorem-pt-stable-triple}
Assume $\gcd(\mathrm{ch}_0,
\mathrm{ch}_1,\mathrm{ch}_2)=1$.
Then every PT stable object comes from
a stable triple $(B,P,\eta)$:
\begin{enumerate}
\item $B$ is a 2-Gieseker
      semistable sheaf with
      $\mathrm{hd}(B)\geq 1$,
\item $P$ is a $0$-dimensional sheaf,
\item $\eta \in
      \Hom_{D^b(X)}(P,B[2])=
      \Ext^2(P,B)$ is nontrivial.
\end{enumerate}
Moreover,
$$
A_\eta := \operatorname{cone}(\eta)[-1].
$$
\end{theorem}

\begin{remark}
\label{remark-wall-crossing-triples}
Stable triples generalize rank $1$
PT stable points.
\end{remark}

\begin{theorem}[Jardim et al.]
\label{theorem-single-wall-crossing}
There is a single wall crossing
that takes the DT moduli space to
the PT moduli space.
Equivalently, the Gieseker and PT
chambers are separated by one wall.
\end{theorem}

\begin{proof}
\begin{enumerate}
\item Understand the stable
objects at the boundary,
labelled $\Theta_v^-$
in the picture.
\item Study the asymptotic stable
objects, that is,
what you find if you go to
infinity.
\item Identify what the PT
chamber is.
This may also be an asymptotic
analysis.
\item Show that there are no other
walls between the $\Theta$ wall
and infinity on the Gieseker side,
and no other walls in the
direction of the PT side.
\end{enumerate}
\end{proof}

\medskip\noindent
History.
\begin{itemize}
\item Pandharipande and Thomas
(2009) suggested that DT/PT
correspondence could be a wall
crossing phenomenon in
Bridgeland stability space.
\item Bayer (2009) established
that the rank $1$ case is a
wall crossing for polynomial
stability conditions.
\item Toda.
\item \dots
\end{itemize}

\medskip\noindent
Case studies.

\noindent
Let $X$ be a smooth threefold with
Picard rank $1$.
\begin{itemize}
\item Collapsing wall.
For
$v=(r,0,0,-n)$,
$$
\mathcal{G}(v)=
\mathrm{Quot}^n(
\mathcal{O}_X^{\oplus r})/\!/
\mathrm{GL}(r),
$$
while $\mathcal{T}(v)=\emptyset$.
\item Fake wall.
For $v=(v_0,v_1,v_2,v_{3,\max})$,
$\mathcal{G}(v)=\mathcal{T}(v)$.
This is a wall that changes nothing.
\end{itemize}

\noindent
Now set $X=\mathbb{P}^3$.
\begin{itemize}
\item For
$v=(2,-1,-\tfrac{1}{2},-\tfrac{1}{6})$,
$\mathcal{G}(v)$ is birational to a
$\mathbb{P}^1$-fibration over
$\mathbb{P}^3\times\mathbb{P}^3$,
while $\mathcal{T}(v)\cong\mathbb{P}^3$.
The wall crossing induces a
retraction from the $6$-dimensional
space $\mathcal{G}(v)$ to
$\mathbb{P}^3\times\mathbb{P}^3$,
and
$\mathcal{T}(v)\simeq \mathbb{P}^3$.
\end{itemize}

\medskip\noindent
There is a map
$$
\xymatrix{
\mathcal{G}(-v)\ar[dr]^\gamma
&&
\mathcal{T}(v)\ar[dl]^\tau
\\
&
\mathcal{M}_{\bar\beta,\bar\alpha,s}(v)
&
}
$$
Pavel and Tajakka show that
the PT moduli stack and the
Bridgeland moduli stack at the wall
are projective schemes.
So the map above is a map of
schemes.
\noindent
Formula of the map missing.

\medskip\noindent
Some questions:
\begin{itemize}
\item (Existence.) Given a projective
variety $X$, are there stability
conditions on $D^\text{b}(X)$?
Positive answer for surfaces,
Fano 3folds with Picard rank 1,
abelian 3folds, quintic 3folds,
products of curves, $\mathbb{P}^n$,
\dots
\item (Moduli spaces.) Is
$M_\sigma(v)$ a projective scheme?
Cannot use usual git techniques to
study. Yes for K3, $\mathbb{P}^2$,
and $\mathbb{P}^1\times \mathbb{P}^1$.
The parameter space is a stack
(ref. missing, involves Toda).
\item What do stable objects
look like?
\end{itemize}

\medskip\noindent
\medskip\noindent
Following are very sketchy notes picked
up along the way; don't hessitate to skip

\begin{enumerate}
\item Stable objects in an abelian category
are the ``building blocks":
we can reconstruct the whole category from
them.
\item  An abelian subcategory (hart)
$\subset$ a
triangulated
\item Stability defined via stability f
unction on $\mathcal{A}$.
\item (Example.) $\mathcal{A}=\text{Coh}X$ is
heart
of $D^b(X)$ with stability function.
\item Bridgeland Stabl:= the stability
conditions are a complex manifold
of complex codimension $\text{rk}\Lambda$:
$$
\mathcal{Z}:\text{Stab}(\mathcal{T})
\longrightarrow
\text{Hom}(\Lambda,\mathbb{C})
$$
\item There's a chamber structure;
moduli space changes across chambers.
\item Picture: blue + black are walls.
Q. What are $\beta$ and $\alpha$?
\item thm: bridgeland stable
= gieseker stable ?
\item Q. slope stability
= bridgeland stability?
A. Not always.
\item DT/PT correspondance:
only one wall between PT and G chambers
\medskip
\item Polynomial stability function.
This is an asymptotic version of BS.
\item There are some $\rho$'s.
Arrangements of $\rho_i$ are polynomial
stability conditions on a threefold.
\item Pata-Thomas introduced stability
for rank 1 objects.
Bayer compares the---wall. Q.
Same for Bridge S---only one wall?
\medskip
\item Def. A {\it stable triple}: when
$\text{gcd}(\text{ch}_0,\text{ch}_2,\text{ch}
_3)=1$,
every PT stable object
comes from three conditions (missing).
\item What happens when
you cross the blue wall?
Both $\mathcal{G}$ and $\mathcal{T}$ are
projective.
\medskip
\item For $X$ smooth threefold with
$\text{rk}\text{Pic}=1$, $\mathcal{G}(v)$,
$v=(r,0,0,-n)$, $\mathcal{G}(v)$ is a known
sheaf object and
$\mathcal{T}=\emptyset$.
\item $X$ sm 3 rkpic1, Fake wall;
$\mathcal{G}=\mathcal{T}$.
\item (Extra.) red circle is a wall for a
weaker form of stability.
\end{enumerate}

\begin{definition}[Talk at impa]
\label{definition-slope-stability}
A rank-2 sheaf $\mathcal{F}$ is
{\it semistable (stable)} if
$H^{0}(\mathcal{F}(t))=0$ for $-t \geq(>)
\frac{c_1F}{2}$
\end{definition}

\noindent
Compare with

\begin{definition}[moduli-curves.tex]
\label{definition-semistable}
Let $f : X \to S$ be a family of curves.
We say $f$ is a {\it semistable
family of curves} if
\begin{enumerate}
\item $X \to S$ is a prestable family of
curves,
and
\item $X_s$ has genus $\geq 1$ and
does not have a rational tail for all $s \in
S$.
\end{enumerate}
\end{definition}

\section{Instantons}
\label{section-instantons}


\begin{itemize}
\item The twistor diagram.
$$
\xymatrix{
&  \ell_\infty \ar@{.>}[dr]\ar[dl]\\
\mathbb{C}P^{2} \ar@{^{(}->}[rr]& &
\mathbb{C}P^{3}\ar[dd]
_{\substack{\text{twistor}
\\ \text{map}}}\\ \\
\mathbb{C}^2\ar@{^{(}->}[u]
\mathbb{R}^4 \ar[u]\ar[rr]& &S^4
}
$$

\item Instantons are solutions to some
Yang-Mills solution.
\item Expository paper
of Donaldson arXiv:2205.08639 \item ADHM
construction, 1978. The first appearance of
Algebraic Geometry in Mathematical Physics.
See Hitchin-Kobayashi correspondence.
\item Donaldson: ``unashamedly
computational''.
\item Expository paper by Simon.
\item Take bundle $(E,\nabla)$ with
an anti-self dual
(ADS) connection on $S^4$ and pullback to
$\mathbb{C}P^{3}$ via the twistor map

$$
\tau[x:y:z:w]=[x+jy:z+jw] \qquad \text{note:
$\tau^{-1}(p)=\mathbb{C}P^{1}$}
$$
Then:
\begin{itemize}
\item Restriction to fibres are trivial.
\item Invariant under anti holomorphic
involution (check this formula!)
$[x:y:z:w] \mapsto [-y:x:-w:z]$.
\item $\mathcal{E}$ is also an instanton
bundle.
\item Penrose Transform:
 $H^1(\mathcal{E}(-2)) \cong \Ker \Delta=0$
 where
$\Delta$ is a Laplacian.
\item Definition of instanton sheaf
on $\mathbb{P}^3$ via $c_1(E)=0$ and some
vanishing of cohomologies.
\item Passage from
[differential equations? algebraic geometry?]
to linear algebra: via {\it monads}.
The point is that instanton sheaf is
equivalent to  ``$E$ being the cohomology
of a linear monad''; theorem by Horps
in the 60's, and is the main tool used
by ADHM. Indeed, ADHM equations come from
the cohomology sequences of the so-called
monads.
\item Mathamatical inst bund:=
locally free instanton sheaves.
\end{itemize}
\end{itemize}

\noindent
{\bf Properties.}
\begin{itemize}
\item The only instanton of rank 1 on
$\mathbb{P}^3$ is
the structure sheaf.
\item non-trivial rank 2 locally free
instanton sheaf is  $\mu_0$ stable
\item double dual is locally free and
also instanton
\item non-trivial rank 2 instanton sheaf
is Fieseker stable
therefore it makes sense to define
moduli space of instanton sheaves as
an open subset of
$\mathcal{M}(c)=\mathcal{G}(k,0,2,0)$.
\end{itemize}

Then studied the irreducibility (Tikhomirov)
and smoothness (Jardim-Verbitsky, 2014. Uses
``3rd hyperkähler quotient'') of
$\mathcal{I}(c)$, the moduli space of rank 2
locally free instanton sheaves of charge $c.$
But nobody likes this results; want new
proofs.

In contrast, $\mathcal{M}(c)$ of rank
2-instanton sheaves of charge $c$ is not
irreducible in general! $\mathcal{M}(1)$ and
$\mathcal{M}(2)$ are irreducible,
$\mathcal{M}(3)$ has exactly 2 irreducible
components of dimension 21; $\mathcal{M}(4)$
has 4 irredicuble components: the locally
free is irreducible, and the other 3 that
intersect the closure of the locally free,
$\overline{\mathcal{I}(c)}$. 3 components of
dimension 29 and one of dimension 32

\begin{remark}
\label{remark-gieseker}
In general the $\mu$ moduli space is
not projective, but the Gieseker is.
\end{remark}

\noindent
{\bf Is $\mathcal{M}(c)$ connected?} Use
$\mathbb{C}^*$ action. True for $c \leq 4$;
every component intersects
$\overline{\mathcal{I}(c)}$ in this range!

\begin{definition}[Elementary transformation]
\label{definition-elementary-transformation}
$F$ of rank 2 locally free instanton,
$Q$ of rank 0 instanton with 1 dimensional
sheaf $h^p Q(-2)$ for $p = 0,1$.
So in this conditions if we have an
epimorphism
 $$
F \overset{\varphi,\text{ surj}}{\to}\} Q
$$
we get that $\Ker \varphi$ is an instanton.
\end{definition}

So  we might be interested in
$$
E \hookrightarrow E^{*
*}\xrightarrow{\text{surj.}} Q_E
$$
Let's have a look at
$\mathcal{M}(3)=\overline{\mathcal{I}(3)}\cup
\overline{C(0,1,3}$. Consider a generic line
bundle of degree 0 on a planar cubic (which
is encoded in that we have intersection of
something of cimension 1 and something of
simension 3), $L inn \text{Pic}^0(C)$ so we
have an epimorphism
$$
\xymatrix{
0\ar[r]&E\ar[r]&\mathcal{O}_{\mathbb{P}^3}
\oplus
\mathcal{O}_{\mathbb{P}^3}\ar[r]&(i_*
L)\ar[r]&0
}
$$
yielding a bundle $E$ as previously outlined.

So we are studying the components via
pushforwards of sheaves on complete
intersection curves inside $\mathbb{P}^3$

``when we do alementary transformation of
rational (not rationl?) we get
something on the boundary of locally free.''

\medskip\noindent
{\bf Perverse instanton sheaves.} Like an
instnaton sheaf in monad description but with
some restrictions on the cohomologies. This
leads to definition of {\it $0$-rimensional
instanton}, a perverse instanton such that
$\mathcal{H}^0=0$.

\medskip\noindent
{\bf Instantons and quivers.}

\medskip\noindent
{\bf Framed instantons.} Fix a line $j:L\to
\mathbb{P}^3$ and $E$ a perverse
sheaf; a {\it framing} at $L$ is an
isomorphis [missing].

Apply GIT to the ADHM data to construct a
moduli space
$$
\mathcal{P}(r,c)=\mathcal{V}(r,c)^{\text{st}}
/\!/\text{GL}(V_c)
$$
and it will follow that $\mathbb{P}(?,?)$ is
connected.

\medskip\noindent

Considerations on quaternionic spaces lead to
generalization of what has been discussed so
far to higher dimensions. I.e. an {\it
instanton sheaf} on $\mathbb{P}^n$ is… [other
cohomological conditions] sheaf

Why are instantons interesting? They are the
simplest; may provide examples for Bridgeland
stability.

In Kuznetsov (2012) and Faenzi (2014)
introduced {\it rank 2 instanton bundles on
Fano 3-folds}. An {\it instanton bundle} on
$X$ is a $\mu$-stable … and some Chern class
is called the {\it charge}. An {\it instanton
sheaf} (introduced by Marcos-Gaia) is ….

 ``Since we are imposing $\mu$-stability on
the defintion we can consider the moduli
$\mathcal{I}_X(c) \subset
\mathcal{G}_X(2,-r_X,c,0)$''.

There's also monad representations as an
ingredient.

\medskip\noindent
Here are two questions that invite us to join
the instanton fever:

\medskip\noindent
{\bf Task 1.} Construct rank 2-instanton
sheaves that do not deform into locally free
ones, and obtain the new irreducible
components of $\mathcal{G}(2,-r_X,c,0)$.

\medskip\noindent
{\bf Task 2.} Nonlocally instanton sheaves
that can be deformed into non-locally free
ones: the {\it instanton boundary}
$\overline{\mathcal{I}(c)}/\mathcal{I}(X)
$[formula
right?]


\medskip\noindent
{\bf Recipe to construct your own instanton.}
\begin{enumerate}
\item (Make a bunch of instantons.) Find an
appropriate curve to
use Serre correspondence to find some
rank 2 instanton sheaf:
$$
\xymatrix{
0\ar[r]&\mathcal{O}_{\mathbb{P}^3}(-1)\ar[r]&
\underbrace{E}_{\exists
}\ar[r]&\mathcal{I}_{\sqcup
\text{lines}}\ar[r]&0
}
$$
where $E$ is a locally free instanton
of charge = the number of lines $-1$.

\item (Is your family of instantons generic?)
You look at the $\text{Ext}$s.
``Therefore, the family of instantons only
defines a locally closed subset
within a generically smooth irreducible
component of $\mathcal{G}$''.

\item ``Find suitable rank 0 instanton
sheaves
 to perform an elementary
transformations on the examples obtained in
Step 1.'' But they are non-locally
free.

\item Now I have my instantons, I know they
are locally free: but how to prove
that the elementary transformations deform
to locally free ones? Looks like
the challenge is to prove that the
deformation is locally free.
\end{enumerate}

In the papers by the group there are several
particular cases when the
deformations are locally free. But they don't
have a general result that would
work for 3-folds.

\medskip\noindent
{\bf What you need to call a thing an
instanton.}
\begin{itemize}
\item Minimal cohomology possible;
try to kill as much cohomology as you can.
\item Fixing $c_1$ (which may determine other
Chern classes).
\item Some stability condition
like $\mu$-stability or quiver stability.
Here is an example that is not
$\mu$-semistable:
$T\mathbb{P}^3(-1)\oplus\Omega_{\mathbb{P}^3}
(1)$
\item Whenever possible,
look for a monadic representation.
(The monadic
representation comes from ADHM ---
the beginnings of this theory.
And it's still here!)
\end{itemize}

\section{Jardim-Muniz}
\label{section-jardim-muniz}

\noindent
Reflexive sheaves
are vector bundles except for a small
locus.

\begin{definition}
\label{definition-reflexive}
Let $X$ be an integral locally Noetherian
scheme. Let $\mathcal{F}$
be a coherent $\mathcal{O}_X$-module.
The {\it reflexive hull}
of $\mathcal{F}$ is the
$\mathcal{O}_X$-module
$$
\mathcal{F}^{**} = \SheafHom_{\mathcal{O}_X}(
\SheafHom_{\mathcal{O}_X}(\mathcal{F},
\mathcal{O}_X), \mathcal{O}_X)
$$
We say $\mathcal{F}$ is {\it reflexive}
if the natural map
$j : \mathcal{F} \longrightarrow
\mathcal{F}^{**}$
is an isomorphism.
\end{definition}

\begin{lemma}
\label{lemma-reflexive-torsion-free}
Let $X$ be an integral locally Noetherian
scheme. Let $\mathcal{F}$
be a coherent $\mathcal{O}_X$-module.
\begin{enumerate}
\item If $\mathcal{F}$ is reflexive, then
$\mathcal{F}$
is torsion free.
\item The map $j : \mathcal{F}
\longrightarrow \mathcal{F}^{**}$
is injective if and only if $\mathcal{F}$
is torsion free.
\end{enumerate}
\end{lemma}

\begin{remark}[Talk at IMPA, 11 June]
\label{remark-reflexive-talk}
Torsion could also be defined so that the
sheaf can inject onto its dual. In this talk
we discussed the moduli space of
reflexive/torsion-free sheaves, which turned
out to be parametrized by $c_1, c_2$ and
$c_3$. This was denoted by $R(c_1,c_2,c_3)$.
Actually I think it may have been Manolache
that proved the existence of this moduli
space.

Alan Muniz

Nesta palestra discutiremos a classificação
de feixes reflexivos de posto dois e seus
espaços de módulos. Apresentaremos algumas
ferramentas básicas usadas na construção e
determinação de tais feixes. Aplicaremos
estas técnicas para o caso de feixes com
segunda classe de Chern igual a quatro,
obtido recentemente em colaboração com Marcos
Jardim.
\end{remark}




\medskip\noindent
Now let's discuss
distributions on manifolds.

Here's the abstract from a talk by Marcos
Jardim at Geometric Structures:

``I will revise the work done over the past
10 years with various collaborators on
distributions and foliations on 3-folds,
especially on the projective space, with a
focus on properties of the tangent sheaf and
singular scheme."

Here are two key ideas: if the distribution
is codimension 1 we can write:
$$
\xymatrix{
0\ar[r]&F\ar[r]&TX\ar[r]^{\omega}&I_Z \otimes
L\ar[r]&0
}
$$
where $L$ is a line bundle and
$\omega \in H^{0}(\Omega_X \otimes L)$, and
$Z=\{p:\omega(p)=0\}$.

When codimension is 2 then $\mathcal{D}$ is
given by
a holomorphic vector field
$\nu$: $T_p=\left\langle{}\nu(p)\right\rangle{}$.

It can be encoded as an exact sequence
$$
\xymatrix{
	0\ar[r]&L\ar[r]^{\nu}&TX\ar[r]&N\ar[r]&0
}
$$
where $L$ is a line bundle
and $\nu \in H^{0}(TX \otimes L^\vee)$;
  $Z=\{p | \nu(p) = 0\}$.

\begin{remark}
\label{remark-stauration}
Saturation means that $Z \subset X$ is a
union
of curves and points.
\end{remark}

And again, distributions are parametrized by
Chern classes.

Two interesting open questions:
\begin{enumerate}
\item {\bf Conjeture.} if $\mathcal{D}$ is a
codimension 1 foliation of degree
$d$ on $\mathbb{P}^3$, then $c_2(F)\leq
d^2-d+1$ and
bound is attained
by rational foliations of type $(1,d+1)$.
(True for $d \leq 2$.)
\item {\bf Conjecture (with Pepe Seade).}
$\mathcal{D}$ is a codimension 1
foliation on a smooth projective 3-fold,
then $\text{Sing}\mathcal{D}$ is connected.
\end{enumerate}

\begin{theorem}[Jardim-Muniz]
\label{theorem-jardim-muniz}
Conditions on Chern classes used to
understand moduli space $R(c_1,c_2,c_3)$.
$c_2=4$ gives (?). For $c_3\leq 6$, possible
``spectrum" exists…
\end{theorem}









\section{Campana}
\label{section-campana}

A \textit{special} variety should be thought
of
as the farthest notion possible from being of
general type.

\begin{definition}[Campana]
\label{definition-special-variety}
$X$ complex projective.
$X$ is \textit{special} if $\forall i \leq  p
\leq \dim X$,
for all line bundle $L \subset \Omega^P_X$,
$\kappa(L)<P$ (Kodaira dimension).
\end{definition}

\noindent
A conjecture by Campana says that
$X$ is special if and only if there exists
$f:\mathbb{C}\to X$
Zariski dense.
Compare this with Kobayashi hyperbolicity!
There are no lines on …?

The main tool for proving this kind of
``Bloch-Ochai" statements (i.e.
something of the kind: if $h^0(\Omega_X)>
\dim X$
then there are no nonconstant
$f:\mathbb{C}\to X$
Zariski dense)
is by passing through the Albanese variety.

In a talk [fill details]
we study a theorem of this kind but with the
bound
given by $q^+(X)$ defined as the supremum of
the $h^0(\hat{X},\pi^*\Omega_{X}$
for $\hat{X}\to X$ some (étale) cover.
Can we prove a BO stement with the condition
$q^+(X)>\dim X$?

Notion of special comes in: indeed
 $$
q^+(X)>\dim X \implies X \text{ not special
}\implies
\not\exists f:\mathbb{C}\to X \text{Zariski
dense}.
$$

The following theorem is not at all easy to
prove:
\begin{theorem}[Campana]
\label{theorem-campana}
$\hat{X}\xrightarrow{\text{étale} }X$, then
$\hat{X}$ is special iff $X$ is special.
\end{theorem}

\noindent
My question: why does the abelianess of
the Albanese variety is important in this
construction?

\bibliography{my}
\bibliographystyle{amsalpha}




\end{document}
